\section{対象と、今回確定する主張}

以下、$X=V(F)\subset\PP^4_{\CC}$ は滑らかなcubic threefoldとする。
「非可換cubic threefold」には複数の定式化があり、変形対象を明示する必要がある。

\begin{center}
\small
\begin{tabular}{p{.18\textwidth}p{.36\textwidth}p{.36\textwidth}}
\toprule
水準 & 変形する対象 & 今回の結論\\
\midrule
構造層 &
$\OO_X$ のmarkedな平坦結合代数層 &
古典的形式基底を5本の二次式で記述\\
次数付き代数 &
$\CC[x_0,\ldots,x_4]/(F)$ の次数を保つ変形 &
一般の形式構成と、Fermatの代数的1パラメータ族\\
残余圏 &
$\Ku(X)=\langle\OO_X,\OO_X(1)\rangle^\perp$ のdg変形 &
古典的Morita変形の10次元基底、固定ファイバー、Fermat族の二次類\\
\bottomrule
\end{tabular}
\end{center}

基底環は可換な局所Artin $\CC$-代数とし、特殊ファイバーとの同一視を固定する。
従って、ここでの「非可換」は変形後の構造代数についての語である。
非可換Artin環をパラメータとするLaudal型変形理論を扱うものではない。

構造層の変形はYekutieliの定義\cite[Definitions 2.7--2.8]{Yek}に従う。
$A\in\Art$ に対し、$\NC_X(A)$ を、位相空間 $X$ 上の平坦な結合 $A$-代数層
$\mathscr B$ と同一視 $\mathscr B/\frakm_A\mathscr B\simeq\OO_X$ の組の、
同一視を保つ代数層同型による同値類とする。
$\Pois_X(A)$ は、可換積そのものも変形でき、特殊ファイバーで括弧が零となる
Poisson代数層の対応する関手である。

\medskip
\noindent\textbf{結論の位置づけ。}
第2節の形式環と第5節の代数族は第1版の内容である。
第7節では、第1版で条件付きだった圏論的比較を、古典的Artin基底上の
compactly generated Morita変形の意味で完成する。
第8節以下のHesse比較では、単なる点曲線の一致で済ませず、
次数シフトと三次周期性を含むdg模型を比較する。
第\ref{sec:gluing-objects}--\ref{sec:points}節は第3版の追加であり、
貼り合わせの局所変形問題と、明示的な点加群スキームを扱う。
第\ref{sec:derived-gluing}--\ref{sec:reconstruction}節が第4版の追加であり、
導来構造、代数的モジュライでの実現、楕円曲線の復元を扱う。
第\ref{sec:general-tangent-comparison}--\ref{sec:five-generator}節が第5版の追加である。
一般の接錐の比較とBrill--Noether境界の六次接錐、
明示的族の評価核による記述、五生成元代数の内在的分類を扱う。
第\ref{sec:picard-compatibility}--\ref{sec:boundary-torelli}節が第6版の追加である。
相対Picardによる適合部分空間の全次数の記述、一般評価核点の境界の形式局所構造、
および六次境界によるTorelliを扱う。
第\ref{sec:global-bn-singular}--\ref{sec:boundary-higher-jets}節が第7版の追加である。
境界の全特異部分、埋め込み解消と対数標準性、普遍同変方程式の9次までの斉次性、
および10次の候補を扱う。
第\ref{sec:tate-boundary-tenth}節が第8版の追加である。
Tate行列による実際の対象・incidence・境界の構成、係数 $1/96$ の厳密計算、
全ての滑らかな三次式でのJacobian障害、10次jetでの非同型、および
incidenceの形式法束模型の非存在を証明する。
第\ref{sec:tate-algebraization-twelfth}節が第9版の追加である。
Tate再帰の有限代数化、収束評価、実際のétale変形族、および12次項と座標変更の作用を扱う。
第\ref{sec:contact-twelve-thirteen}節が第10版の追加である。
12次商の次元・生成元・係数汎関数を決定し、実際の12次類の消滅を証明する。
また13次の接触商の消滅により、実際の境界を13次まで正規化する。
第\ref{sec:contact-fourteen}節が第11版の追加である。
14次の接触階数・商・係数汎関数を決定し、低次安定化群の作用を含めて、
実際の14次項の消去を証明する。対角評価の限界と非対角微分の役割も明示する。
第\ref{sec:all-order-parity}--\ref{sec:contact-fifteen}節が第12版の追加である。
Clifford反転により全奇数次を同時に消去し、12次・14次の正規化との両立を証明する。
またPfaffian微分の全次数の全射性、15次の接触商の消滅、
実際の15次係数の厳密計算、およびKoszulホモロジーの非消滅を扱う。

第\ref{sec:algebraic-parity}--\ref{sec:contact-sixteen}節が第13版の追加である。
全次数偶化の代数性と収束性、16次接触商と低次安定化群の作用、
一般テンソルによる厳密な階数証明、および実際の16次係数の判定を扱う。




\section{構造層の形式変形空間：17次元のPfaffian特異点}

\subsection{主張}

$0\leq i<j\leq4$ に対して10個の変数 $a_{ij}$ をとり、
\[
 q_{ijkl}=a_{ij}a_{kl}-a_{ik}a_{jl}+a_{il}a_{jk}
 \qquad(0\leq i<j<k<l\leq4)
\]
とおく。これは5個のPlücker二次式である。

\begin{theorem}[markedな古典的形式変形基底]\label{thm:base}
非標準的な関手の同型
\[
 \NC_X(A)\simeq\operatorname{Hom}_{\mathrm{cont,loc},\CC}(R_{\NC},A)
\]
が存在し、
\begin{equation}\label{eq:base}
 R_{\NC}=
 \frac{\CC[[t_1,\ldots,t_{10},a_{01},a_{02},a_{03},a_{04},
                         a_{12},a_{13},a_{14},a_{23},a_{24},a_{34}]]}
 {(q_{0123},q_{0124},q_{0134},q_{0234},q_{1234})}.
\end{equation}
従って古典的形式変形空間は
\[
 \widehat{\Aff}^{10}\times
 \widehat{C(\Gr(2,5))}_{\,0}
\]
である。
\end{theorem}

これはYekutieliの量子化定理\cite[Corollary 11.2]{Yek}と、
MayanskiyのSchouten括弧計算\cite[Section 2.2]{May}の相対版から導く帰結である。
以下では非被約な基底上の議論を含めて証明する。

\subsection{証明：Poisson変形への移行}

三次超曲面の標準的なコホモロジー計算から
\[
 H^1(X,\OO_X)=H^2(X,\OO_X)=0.
\]
YekutieliのCorollary 11.2は
\[
 \Pois_X(A)\xrightarrow{\;\sim\;}\NC_X(A)
\]
を与え、パラメータ環の準同型との可換性を明記している。
従ってこれはArtin変形関手の同型である。
この定理は、固定した $\OO_X[[h]]$ 上の大域star積だけに限定されず、
可換な構造層の変形も含むことが重要である。

Poisson変形の可換代数層は通常のスキーム変形を与える。
実際、滑らかなアフィン開集合上では、同論文のGeometrization Theorem 3.6により、
基礎の可換代数層は $A\otimes\OO_U$ に局所的に同型である。
これらは可換代数層の同型で貼り合わされるので、通常の $A$-スキーム $X_A$ を与える。
従って、$X$ の通常の変形と、その上の相対Poisson括弧を分類すればよい。

\subsection{証明：通常の変形の普遍形式族}

Euler列と法束列から
\[
 H^0(T_X)=H^2(T_X)=0,\qquad h^1(T_X)=10.
\]
特に通常のmarked変形には無限小自己同型がなく、障害もない。
$W_3=H^0(\PP^4,\OO(3))$ は35次元であり、
\[
 \mathfrak{gl}_5\longrightarrow W_3,\qquad M\longmapsto M(F)
\]
は単射である。実際、核の元は $X$ のベクトル場を与え、その消滅からスカラー行列となるが、
スカラー行列 $\lambda I$ は $F$ を $3\lambda F$ に写す。
従って核は零である。

$\mathfrak{gl}_5\cdot F$ の10次元の補空間をとり、その基底を $G_1,\ldots,G_{10}$ とする。
\[
 S=\CC[[t_1,\ldots,t_{10}]],\qquad
 F_t=F+\sum_{\alpha=1}^{10}t_\alpha G_\alpha
\]
は形式的なsliceを与える。これは $GL_5$ の作用の微分と補空間の直和が同型であること、
および形式的逆関数定理による。

$\OO_X(1)$ の延長の障害・曖昧さはそれぞれ $H^2(\OO_X)$、$H^1(\OO_X)$ に入り、
ともに零である。その切断も延長し、変形は再び $\PP^4$ 内の三次超曲面となる。
従ってこのsliceは抽象的な通常のmarked変形をも表す。
特殊ファイバーで恒等となる自己同型は、Artin環の長さに関する帰納法と
$H^0(T_X)=0$ から恒等である。

\subsection{証明：相対括弧と冪零基底}

$p:X_S\to\Spf S$ を上の普遍族とする。
線形形式の自由加群 $V_S=p_*\OO_{X_S}(1)$ と随伴の同一視を選ぶ。
基底からの可逆層の因子が現れても、局所形式基底では自明化できる。
相対Euler列・余法束列から、任意のArtin基底変換と可換な同型
\[
 p_*\Lambda^2T_{X_S/S}\simeq\Lambda^2 V_S,\qquad
 p_*\Lambda^3T_{X_S/S}\simeq\Sym^2V_S
\]
を得る。階数はそれぞれ10、15である。
例えば $\Lambda^2T\simeq\Omega^1(2)$ と
\[
 0\longrightarrow\OO_{X_S}(-1)
 \longrightarrow\Omega^1_{\PP^4_S/S}(2)|_{X_S}
 \longrightarrow\Omega^1_{X_S/S}(2)\longrightarrow0
\]
を用いる。左端は特殊ファイバーで全コホモロジーが零であり、Artin基底上でも同じである。

$\epsilon_{ij}=x_j\,dx_i-x_i\,dx_j$ をこの同型を通じて双ベクトルとみなし、
$\pi=\sum_{i<j}a_{ij}\epsilon_{ij}$ と書く。
Mayanskiyの計算は、規約をこのように固定すると
\begin{align}
 \frac14[\pi,\pi]
 ={}&q_{1234}(F_t)_0-q_{0234}(F_t)_1+q_{0134}(F_t)_2 \notag\\
    &-q_{0124}(F_t)_3+q_{0123}(F_t)_4, \label{eq:bracket}
\end{align}
を与える。ここで $(F_t)_i=\partial F_t/\partial x_i$ である。
この式は微分と斉次Euler公式から得る多項式恒等式なので、相対的にも成立する。

滑らかさから $F_0,\ldots,F_4$ は線形独立である。
定数係数の非自明な線形関係があれば、座標変換後に $F$ が一変数を含まず、
$X$ が特異な射影錐になるからである。
15個の二次単項式で勾配を表した係数行列の、ある $5\times5$ 小行列式は特殊ファイバーで非零である。
従って形式基底上で可逆であり、
\[
 S^5\longrightarrow\Sym^2V_S,\qquad (b_i)\longmapsto\sum_i b_i(F_t)_i
\]
は分裂単射となる。この単射性は任意のArtin基底変換の後も保たれる。

よって \eqref{eq:bracket} が零であることは、非被約基底上でも
\[
 q_{0123}=q_{0124}=q_{0134}=q_{0234}=q_{1234}=0
\]
と同値である。量子化による関手同型と合わせて命題\ref{thm:base}を得る。
\hfill$\square$

\subsection{次元、特異点、および主張の限界}

$\Gr(2,5)$ の次元は6なので、そのアフィン錐は7次元である。
従って
\[
 \dim R_{\NC}=17,\qquad
 \dim_\CC\frakm_{R_{\NC}}/\frakm_{R_{\NC}}^2=20.
\]
5本の二次式は最小生成系だが、余次元は3である。
従って「20変数から5を引いて15次元」と数えてはならず、この環は完全交叉ではない。
例えば $a_{01}\ne0$ のチャートでは、3本の式から $a_{23},a_{24},a_{34}$ を解け、
残りの2本は従属する。

錐は頂点以外で滑らかであり、被約な特異点集合は $a=0$、すなわち10次元の可換変形の部分である。
Jacobianイデアルが定める特異部分スキームの厚みまでは、この言い方では主張していない。

HKR分解と上のコホモロジー計算は
\[
 \HH^2(X)\simeq H^1(T_X)\oplus H^0(\Lambda^2T_X)\simeq\CC^{10}\oplus\CC^{10},
 \qquad \HH^3(X)\simeq\CC^{15}
\]
を与える\cite{May}。$\HH^3$ は障害の入り得る空間であり、最小方程式の本数ではない。
また、\eqref{eq:base} は\textbf{古典的集合値関手}の環であり、
derived formal moduliの全ての高次構造をこの通常の環に同定する主張ではない。
markingを忘れれば $\operatorname{Aut}(X)$ の有限群作用も考慮する必要がある。
積分解はsliceと自明化の選択に依存する。

\section{同時に動かせる非可換方向は最大4次元}

17次元の形式空間の特異性には、具体的な意味がある。
固定したcubic上のPoissonテンソルは、$\Lambda^2V$ の分解可能な元、
すなわちGrassmannian錐の点に対応する\cite[Theorem 1]{May}。
任意の二つのPoisson方向を独立なパラメータで同時に量子化できるとは限らない。

\begin{lemma}\label{lem:linear}
$V$ を5次元とし、線形部分空間 $L\subset\Lambda^2V$ の全ての元が分解可能とする。
このとき、ある直線 $\ell\subset V$ または3次元部分空間 $W\subset V$ に対して
\[
 L\subset\ell\wedge V
 \quad\text{または}\quad
 L\subset\Lambda^2W.
\]
特に $\dim L\le4$ である。
\end{lemma}

\begin{proof}
非比例な $a,b\in L$ に対し $(a+b)^2=0$ なので $a\wedge b=0$。
対応する二つの2平面は直線で交わり、基底を選んで
$\langle e_0,e_1\rangle$、$\langle e_0,e_2\rangle$ と書ける。
双方と交わる別の2平面は、$\langle e_0\rangle$ を含むか、
$W=\langle e_0,e_1,e_2\rangle$ に含まれる。
後者で $e_0$ を含まないものが一つでもあれば、他の全ての2平面も $W$ に含まれる。
そうでなければ全てが $\langle e_0\rangle$ を含む。
\end{proof}

従って、可換点を通る滑らかな形式基底
$\Spf\CC[[s_1,\ldots,s_r]]$ 上の族の非可換接方向の像は最大4次元である。
実際、$a_{ij}=a_{ij}(s)$ をPlücker式に代入した次数2の項は、
一次項の像が補題の $L$ になることを意味する。
幾何学的方向も含めた接写像の階数は最大14となる。
これは\textbf{基底の次元や一般点での階数の上限ではなく、可換点での接写像の上限}である。
$\ell$ を固定して $\ell\wedge V$ を動かすと、4方向の上限は達成される。

もう一つの帰結として、非零Poissonテンソル $\pi$ における
$[\pi,-]:H^0(\Lambda^2T_X)\to H^0(\Lambda^3T_X)$ の階数は3である。
これはPlücker式の微分と分裂単射な勾配写像の合成だからである。
MayanskiyのPoissonコホモロジー公式\cite[Theorem 2]{May}を代入すると
\[
 \dim\HP^0=1,\quad \dim\HP^1=0,\quad
 \dim\HP^2=17,\quad \dim\HP^3=12.
\]
従って障害群が非零でも、Poisson対の通常のモジュライはこの点で滑らかであり得る。

\section{一般のcubic：次数を保つ形式量子化}

\subsection{線形束に対応するJacobian括弧}

座標を $(u,v,x,y,z)$ とする。三次式 $F=F(u,v,x,y,z)$ に対して
\begin{equation}\label{eq:jac}
 \{x,y\}=F_z,\qquad \{y,z\}=F_x,\qquad \{z,x\}=F_y,
 \qquad \{u,-\}=\{v,-\}=0
\end{equation}
とおく。これは $\CC[u,v]$ 上の三変数Jacobian Poisson括弧であり、
$F$ はCasimirである。括弧の係数は二次式なので
$\deg\{f,g\}=\deg f+\deg g$ となる。

この括弧は $F=0$ に降り、次数0の局所化を通して $X$ 上のPoisson構造を与える。
対応する分解可能な2形式は $u\,dv-v\,du$ に比例する。
任意の非零分解可能な方向は線形座標変換でこの形にできる。
ただしArtin基底上の普遍Plücker解が常に二つの線形形式の外積として選べる、
という主張ではない。

\subsection{中央の三次式を同時に量子化する}

\begin{theorem}[各Poisson方向の形式的な次数付きモデル]\label{thm:formalgraded}
\eqref{eq:jac} に対して、次数を保つ結合形式変形 $B_h$ と次数3の中心元
$\widehat F=F+O(h)\in B_h$ が存在する。
$u,v$ は厳密に中心的であり、
\[
 A_h^{\mathrm{form}}=B_h/(\widehat F)
\]
は各次数で有限自由な $\CC[[h]]$-加群となる。そのHilbert級数は
\[
 H_{A_h^{\mathrm{form}}}(t)=\frac{1-t^3}{(1-t)^5}
\]
であり、一次交換子は \eqref{eq:jac} を与える。
\end{theorem}

\begin{proof}
アフィン5空間上のKontsevichの多項式formality写像 $\mathcal U$ を使う
\cite[Sections 4.6, 6, 8]{Kon}。
対応する正規化star積
\[
 f*g=fg+\frac h2\{f,g\}+O(h^2)
\]
の各グラフでは、$n$ 個の双ベクトル頂点が係数の次数を $2n$ 増し、
$2n$ 本の辺が微分で同じ次数を減らす。
従って積は次数を保つ。双ベクトルには $u,v$ 方向の微分がないため、
各双微分作用素は $x,y,z$ のみを微分する。
単位の正規化から $u*f=uf=f*u$、$v*f=vf=f*v$ が成り立つ。

twistした $L_\infty$ 写像の接写像は、Poisson複体からHochschild複体への鎖写像である。
$[\pi,F]=0$ なので
\[
 \widehat F=\mathcal U^{h\pi}_1(F)
 =\sum_{n\ge0}\frac{h^n}{n!}\mathcal U_{n+1}(\pi,\ldots,\pi,F)
\]
はstar積について中心的である。係数は三次斉次式である。
実際、$F$ 頂点を加えたグラフにも辺は $2n$ 本あり、
入力の多項式次数は $2n+3$ だから、出力は次数3となる。

最後に、各次数 $d$ の乗法写像
\[
 B_{h,d-3}\xrightarrow{\ *\widehat F\ }B_{h,d}
\]
は有限自由 $\CC[[h]]$-加群の間の写像であり、$h=0$ で非零多項式 $F$ による単射となる。
従って最大階数のある小行列式が単元であり、分裂単射で余核は自由である。
余核の階数は通常の三次超曲面と一致する。
\end{proof}

ここで完備化は各次数で行う。命題は形式級数としての構成であり、
$h\in\CC^\times$ への収束・代数化を与えてはいない。
また、商で使うのは一般に補正された $\widehat F$ である。
可換な方程式 $F$ をそのまま量子化代数で割ればよい、とは限らない。
次節は、この補正を有限な式で全て書ける例である。

\section{Fermat cubicの明示的な代数的非可換族}

\subsection{定義と中心元}

$h$ を次数0の中心パラメータとし、
\begin{equation}\label{eq:skly}
 \mathscr S=
 \frac{\CC[h]\langle x,y,z\rangle}
 {(xy-yx-hz^2,\ yz-zy-hx^2,\ zx-xz-hy^2)}
\end{equation}
とおく。$x,y,z$ の次数は1とする。
\[
 D=x^3+y^3+z^3,\qquad
 T=xyz+xzy+yxz+yzx+zxy+zyx
\]
と書く。

\begin{theorem}[次数3の中心元の明示式]\label{thm:center}
$\mathscr S$ において
\begin{equation}\label{eq:center}
 G_h=D-\frac{h^2}{3}T
\end{equation}
は中心的である。等価な表示は
\begin{equation}\label{eq:centerordered}
 G_h=(1-h^3)(x^3+z^3)+(1+h^3)y^3-2h^2zyx
\end{equation}
である。この恒等式は全ての $h$ に対して成り立つ。
\end{theorem}

Sklyanin型代数の中心と三次ポテンシャルについては
Etingof--Ginzburg\cite{EG}の体系的な研究がある。
ここでは規約の符号に依存しないよう、上の表示から直接証明する。

\begin{proof}
$C=y^3-z^3$ とおく。\eqref{eq:skly} から
\[
 [yz,x]=[zy,x]=hC
\]
であり、さらに
\begin{align*}
 [D,x]
 &=h(zy^2z-yz^2y)\\
 &=h[zy,yz]
 =h^2[yz,x^2]
 =h^3(Cx+xC).
\end{align*}
6個の単語を $zyx$ の順に整理すると
\[
 T=6zyx+3h(x^3-C).
\]
従って
\[
 [T,x]=6hCx-3h[C,x]=3h(Cx+xC).
\]
ゆえに $[D-h^2T/3,x]=0$。
$D,T$ と関係式は巡回置換で不変なので、$y,z$ との交換子も零である。
上の $T$ の表示を代入すれば \eqref{eq:centerordered} を得る。
$h$ で割る操作は用いていない。
\end{proof}

この計算は、元の三次式 $D$ をそのまま使う問題点も示している。
$[D,x]=h^3(Cx+xC)$ であり、一般に $D$ 自身は中心元ではない。
中心性を保証する $-h^2T/3$ の補正が本質的である。

\subsection{非可換cubic代数族と平坦性}

次数1の中心変数 $u,v$ を加え、
\begin{equation}\label{eq:family}
 \boxed{\quad
 \mathscr A=\mathscr S[u,v]/(G_h+u^3+v^3).
 \quad}
\end{equation}
$S_\lambda=\mathscr S/(h-\lambda)$、
$A_\lambda=\mathscr A/(h-\lambda)$ と書く。

\begin{theorem}[全パラメータで平坦な次数付き族]\label{thm:family}
$\mathscr A$ は各次数で有限自由な $\CC[h]$-加群であり、
\[
 A_0\simeq
 \CC[x,y,z,u,v]/(x^3+y^3+z^3+u^3+v^3).
\]
全ての $\lambda\in\CC$ に対して
\[
 H_{A_\lambda}(t)=\frac{1-t^3}{(1-t)^5}.
\]
また $\lambda\ne0$ なら $A_\lambda$ は非可換である。
\end{theorem}

\begin{proof}
$S_\lambda$ のSklyaninパラメータは
\[
 (p,q,r)=(1,-1,-\lambda)
\]
である。少なくとも二つのパラメータは非零で、
$p^3=1$ と $q^3=-1$ は異なる。
従ってArtin--Tate--Van den Berghの定理、または
Iyudu--Shkarinの代数的証明\cite{IS}により、
全ての $\lambda$ で $S_\lambda$ は3次元AS正則・Koszulで、
$H_{S_\lambda}(t)=(1-t)^{-3}$ である。

$\mathscr S_n$ は有限生成 $\CC[h]$-加群で、その各閉ファイバーの次元は一定である。
PID上の有限生成加群の構造定理から、ねじれ部分は零となり、$\mathscr S_n$ は自由である。
次に $u^3+G_h+v^3$ は中心変数 $u$ についてmonicであるため、割り算により
\[
 \mathscr A\simeq
 \mathscr S[v]\oplus\mathscr S[v]u\oplus\mathscr S[v]u^2
\]
という加群の直和分解を得る。従って次数ごとの自由性と
\[
 H_{A_\lambda}(t)
 =\frac{1+t+t^2}{(1-t)^4}
 =\frac{1-t^3}{(1-t)^5}
\]
が従う。

$\lambda\ne0$ では $[x,y]=\lambda z^2$ である。
$S_\lambda$ の次数2で $z^2\ne0$ は関係式から直接分かり、
次数3の斉次元で割っても次数2の成分は変わらない。
従って $A_\lambda$ は非可換である。
\end{proof}

$B_\lambda=S_\lambda[u,v]$ はAS正則次元5・Gorensteinパラメータ5である。
$G_\lambda+u^3+v^3$ は中心的な非零因子なので、change-of-ringsの議論
\cite[Lemma 1.2]{HeYe}から
$A_\lambda$ はAS Gorenstein次元4・パラメータ2となる：
\[
 \Ext^i_{A_\lambda}(\CC,A_\lambda)=
 \begin{cases}
 \CC(2)&i=4,\\
 0&i\ne4,
 \end{cases}
\]
右加群についても同様である。
この性質とHilbert級数は、次数付き超曲面モデルとしてcubic threefoldに適合する。
一方、$\qgr A_\lambda$ の有限大域次元やdg滑らかさには別の議論が必要である。
前者は次節で開集合 $\lambda^6\ne-27$ 上について証明する。

\subsection{半古典極限と点加群の三次曲線}

一次交換子は
\[
 \{x,y\}=z^2,\qquad\{y,z\}=x^2,\qquad\{z,x\}=y^2,
 \qquad \{u,-\}=\{v,-\}=0
\]
を与える。これはFermat式の、$u,v$ に相対的なJacobian括弧の $1/3$ 倍である。
射影チャート $u\ne0$ では $v/u$ がCasimirとなり、
線形束 $[u:v]$ がPoisson葉の幾何を支配する。

$S_\lambda$ の関係式を多重線形化すると、
点 $(x:y:z)$ と次の点 $(x':y':z')$ の間に
\[
 M_\lambda
 \begin{pmatrix}x'\\y'\\z'\end{pmatrix}=0,\qquad
 M_\lambda=
 \begin{pmatrix}
 -y&x&-\lambda z\\
 -\lambda x&-z&y\\
 z&-\lambda y&-x
 \end{pmatrix}
\]
を得る。直接計算により
\[
 \det M_\lambda=-\lambda\bigl(x^3+y^3+z^3+\lambda^2xyz\bigr).
\]
$\lambda\ne0$ で現れるHesse曲線は
\begin{equation}\label{eq:elliptic}
 E_\lambda:\quad x^3+y^3+z^3+\lambda^2xyz=0.
\end{equation}
その三つの偏微分に共通射影零点があるための条件は
$\lambda^6=-27$ である。従って
$\lambda\ne0,\ \lambda^6\ne-27$ でこの曲線は滑らかである。
$\lambda=0$ では $S_0$ 自体の点スキームは $\PP^2$ なので、
極限の三次曲線と点スキームを同一視してはいけない。

曲線の式に最初の変化が $\lambda^2$ で現れることは、
一次情報だけでは捉えられない幾何の一例である。
ただし、これだけで $\Ku(A_\lambda)$ の変形を計算したことにはならない。

\subsection{標準中心元との一致}

Walton--Wang--Yakimov\cite[Lemma 2.19]{WWY}に記載された標準的なSklyanin中心元に
$(a,b,c)=(1,-1,-\lambda)$ を代入すると、
\[
 g=(\lambda^4-\lambda)y^3+(1+\lambda^3)yxz
       +(\lambda^3-1)xyz-(\lambda+\lambda^4)x^3.
\]
関係式を用いて整理すると厳密に $g=-\lambda G_\lambda$ となる。
従って $\lambda\ne0,\ \lambda^6\ne-27$ で
\[
 S_\lambda/(G_\lambda)\simeq B(E_\lambda,\mathcal L,\sigma),
 \qquad
 \qgr(S_\lambda/(G_\lambda))\simeq\operatorname{coh}E_\lambda.
\]
$B$ はtwisted homogeneous coordinate ringである。
$\lambda=0$ では通常のFermat平面三次曲線の斉次座標環となる。
従って $\lambda^6\ne-27$ の全域で、中心商のtailsは滑らかな楕円曲線である。

\section{Fermat族の射影的正則性：大域次元3}

\begin{theorem}\label{thm:smooth}
$\lambda^6\ne-27$ とする。このとき
\[
 \operatorname{gldim}(\qgr A_\lambda)=3.
\]
ここで大域次元はcoherent対象間のExtの消滅で測る。
従って $A_\lambda$ は、tailsの有限大域次元を定義とする意味で
非可換graded isolated singularityである。
\end{theorem}

以下の証明は、matrix factorizationの定理と非可換Serre双対を組み合わせるものである。
この結論を、そのままperfect diagonalの意味でのdg滑らかさと言い換えない。

\begin{proof}
$S=S_\lambda,\ g=G_\lambda,\ A=A_\lambda,\ C=S/(g)$ と略記する。
$S$ は左右noetherianな3次元AS正則代数である。
滑らかな楕円曲線の場合のこの性質と中心商の記述はSklyaninの標準的結果
\cite[Section 2.3]{WWY}であり、$\lambda=0$ は可換多項式環である。

\textbf{双対性の準備。}
三次tensor
\[
 w=xyz+yzx+zxy-(xzy+zyx+yxz)-\lambda(x^3+y^3+z^3)
\]
は巡回不変で、その左偏微分は $S$ の三つの関係式である。
AS正則性とsuperpotentialの判定\cite[Lemma 2.15]{IM}から $S$ は3-Calabi--Yauである。
中心多項式拡大 $B=S[u,v]$ と中心正則元 $W=g+u^3+v^3$ による商を考えると、
$A$ のbalanced dualizing complexは
\[
 A(-2)[4]
\]
となる。従って $\chi$、有限局所コホモロジー次元、および
perfect対象上のSerre双対が使える\cite[Theorem 4.2, Corollary 4.3]{YZ}。
ここでは有限大域次元を先に仮定していない。

\textbf{特異圏の消滅。}
$\pi_A:\operatorname{gr}A\to\qgr A$ を商関手とし、
\[
 \mathcal P=\operatorname{thick}\{\pi_AA(n)\mid n\in\mathbb Z\},
 \qquad \mathcal Q=D^b(\qgr A)/\mathcal P
\]
とおく。$\mathcal Q=0$ を示す。

任意のcoherent対象 $M=\pi_A N$ の有限生成graded lift $N$ を選ぶ。
$B$ はnoetherian AS正則であり、$W$ は中心正則元なので、
Cassidy--Conner--Kirkman--Moore\cite[Theorems 4.7, 5.5]{CCKM}により
$\Dsg^{\mathrm{gr}}(A)$ の対象を $W$ のmatrix factorizationで表せる。
その奇微分を $d$ と書くと $d^2=W$ であり、
中心変数による形式微分から
\[
 d(\partial_ud)+(\partial_ud)d=3u^2,\qquad
 d(\partial_vd)+(\partial_vd)d=3v^2.
\]
従って特異圏では $u^2,v^2$ の作用が零となる。
$D^b(\operatorname{gr}A)\to\mathcal Q$ はgraded perfect複体を零にするので、
この零化は $\mathcal Q$ にも降りる。
導来商関手の本質的全射性を仮定する必要はなく、選んだ実際のlift $N$ に適用すればよい。

$N$ に対する二重Koszul複体 $K(u^2,v^2;N)$ のコホモロジーは
$(u^2,v^2)$ に殺される。
$I=(u,v)$ とすれば $I^3\subset(u^2,v^2)$ なので、
各コホモロジーは長さ高々3の $I$-filtrationを持ち、その商は $C$-加群となる。

一方、$D^b(\qgr C)=D^b(\operatorname{coh}E_\lambda)$ は $\pi_CC(n)$ で厚く生成される。
これはample sequenceの生成性からも分かるが、
大域次元1と二段の有限表示を使えば直接に示せる。
二項複体の負次数コホモロジーから零次コホモロジーへの連結写像は
$\Ext^2$ に入り零なので、零次コホモロジーがその二項複体の直和因子となるからである。

$u,v$ は $A$ 上の中心正則列である。
$u$ の非零因子性はmonic表示から、$A/(u)$ 上の $v$ の非零因子性は $g$ の正則性から従う。
従ってKoszul分解により $\pi_A C(n)\in\mathcal P$。
前段のfiltrationとexactな制限関手から
\[
 \pi_AK(u^2,v^2;N)\in\mathcal P
\]
を得る。

しかし $\mathcal Q$ では $u^2,v^2$ の作用は零なので、この二重coneは
$M$ の四つのshift・twistの直和になる。
その直和が零である以上、$M=0$ である。
全てのcoherent対象について成立するので $D^b(\qgr A)=\mathcal P$。

\textbf{一様な大域次元の上界。}
双対性を全てのperfect対象へ拡張する正確な結果として、
Chen\cite[Theorem 3.10]{Chen}のgeneralized Serre functor
$-\otimes_A^{\mathbf L}R_A[-1]$ を用いる。
$\pi_A A(n)$ はコホモロジー次元3から有限射影次元を持つので適用でき、
有限coneと直和因子を通して $\mathcal P$ 全体に拡張する。
いま全てのcoherent対象が $\mathcal P$ に入るので
\[
 \Ext^i(M,N)^*\simeq\Ext^{3-i}(N,M(-2)).
\]
$i>3$ なら右辺は零である。
さらに $\Ext^3(\OO,\OO(-2))\simeq\CC$ なので、大域次元はちょうど3となる。
\end{proof}

$\lambda^6=-27$ の6点については、この証明は適用しない。
第\ref{sec:exceptional}節では非perfectな行列因子化を構成して、
これらのファイバーの大域次元が無限であることを示す。

\section{Kuznetsov成分の形式Morita変形と固定ファイバー}

\subsection{Hochschild変形のhomotopy abelian性}

滑らかなcubic threefoldでは
\[
 D^b(X)=\langle\Ku(X),\OO_X,\OO_X(1)\rangle.
\]
Kuznetsovの計算\cite[Theorem 8.9]{Kuz}は
\[
 \HH^i(\Ku(X))=
 \begin{cases}
 \CC&i=0,\\
 \CC^{10}&i=2,\\
 0&i\ne0,2
 \end{cases}
\]
を与える。また自然な幾何学的写像
\[
 H^1(T_X)\longrightarrow\HH^2(\Ku(X))
\]
は同型である\cite[Theorem 4.4(2)]{JLLZ}。

このコホモロジーの次数だけで、Hochschild dg Lie代数
$L=CC^\bullet(\Ku(X))[1]$ のhomotopy abelian性が分かる。
その最小 $L_\infty$ モデルの台は次数 $-1$ と $1$ のみである。
$n$ 項演算 $\ell_n$ の次数は $2-n$ なので、
次数 $-1$ の入力が $m$ 個あると、出力次数は
\[
 (n-2m)+(2-n)=2-2m
\]
となる。これは偶数であり、台には存在しない。
従って全ての高次括弧が零である。

よってHochschild/Maurer--Cartan形式変形問題の古典的基底は
$\Spf\CC[[s_1,\ldots,s_{10}]]$ である。
$\HH^0$ に由来する高次自己同型は、集合値の基底には表れていない。
これだけでuncurved Morita変形との比較が済むわけではない。
Keller--Lowen\cite{KL}の曲率に関する注意を踏まえ、以下では実在する
幾何学的変形を使って比較を証明する。

\subsection{実在するMorita変形の形式基底}

$\mathcal C=\Ind(\Ku(X))$ とし、$\CatDef^c_{\mathcal C}$ を
compactly generatedな線形圏の標識付き変形空間とする。
以後は可換局所Artin $\CC$-代数に制限し、$\pi_0$ で同値類を取る。

\begin{theorem}[古典的形式categorical Torelli]\label{thm:formal-torelli}
自然な写像
\[
 \Def_X\longrightarrow\pi_0\CatDef^c_{\mathcal C}
\]
は同型である。従って右辺は滑らかな10次元の形式空間で表される。
\end{theorem}

\begin{proof}
$\mathcal C$ はproperなcompact generatorを持つのでtamely compactly generatedである。
Blanc--Katzarkov--Pandit\cite[Section 3.3]{BKP}が記述するLurieの比較定理により
\[
 \theta:\CatDef^c_{\mathcal C}\longrightarrow\CatDef^\wedge_{\mathcal C}
\]
は $(-1)$-truncatedである。特に各Artin基底で $\pi_0\theta$ は単射であり、
右辺の形式完備化はHochschild複体で制御される。
前小節から $\pi_0\CatDef^\wedge_{\mathcal C}$ は滑らかな10次元形式空間である。

通常の三次三様体の変形に相対例外対の直交補を対応させる写像を考える。
その合成
\[
 \Def_X\longrightarrow\pi_0\CatDef^c_{\mathcal C}
 \longrightarrow\pi_0\CatDef^\wedge_{\mathcal C}
\]
の微分は $H^1(T_X)\xrightarrow{\sim}\HH^2(\Ku(X))$ である。
両端は滑らかな10次元の表現可能な形式変形関手なので、形式逆関数定理から合成は同型。
従って $\pi_0\theta$ は全射でもあり、結論を得る。
\end{proof}

ここで使う文献の入力は、全Artin基底に対する $\theta$ の単射性である。
\cite[Theorem 4.27]{BKP}の一変数形式族に対する全射性だけから、
上の多変数の主張を結論しているわけではない。
また、圏の変形の高次空間まで離散的な $\Def_X$ と同一視する主張ではない。

\subsection{非可換層から残余圏への自然な写像}

構造層の結合代数変形から、Lowen\cite[Theorem 3.26]{LowenPre}により
$\operatorname{Qch}(X)$ の平坦なGrothendieck圏の変形 $\mathcal D_A$ を得る。
\[
 \Ext^1(\OO_X(i),\OO_X(i))=
 \Ext^2(\OO_X(i),\OO_X(i))=0
\]
なので各 $\OO_X(i)$ は一意な同型類として持ち上がる。持ち上げを $E_i$ とする。
$G=\bigoplus_{i=0}^4\OO_X(i)$ は $\PP^4$ のBeilinson生成子の制限であり、
随伴と閉埋め込みの押し出しの保守性によりcompact generatorである。
De Deken--Lowen\cite[Propositions 5.8--5.10]{DDL}から
$G_A=\bigoplus_{i=0}^4E_i$ もcompact generatorとなる。

Derived base changeとNakayamaから $(E_0,E_1)$ は相対例外対である。
その右直交成分は $P_A(G_A)$ で生成される。ただし
$P_A=L_{E_0}L_{E_1}$ である。
\cite[Lemma 5.14]{DDL}からmutationは基底変換と可換なので、
\[
 \RHom_{\mathcal D_A}(P_A(G_A),P_A(G_A))
\]
は残余圏の実在するMorita変形を与える。生成子の選択を変えてもMorita同値類は変わらない。
従ってArtin基底上の自然な写像
\[
 \Phi:\NC_X\longrightarrow\pi_0\CatDef^c_{\mathcal C}
\]
が構成される。この議論で $\OO_X(i)$ がfull helixであるとは仮定しない。

\begin{lemma}\label{lem:bivector-zero}
接写像 $\HH^2(X)\to\HH^2(\Ku(X))$ は、通常のHKR分解の
$H^0(\Lambda^2T_X)$ を零に写す。
\end{lemma}
\begin{proof}
JLLZ\cite[Theorems 3.10, 4.4(2)]{JLLZ}のTodd補正付きHKRの加群構造と
$\gamma:\HH^2(\Ku(X))\to\operatorname{Hom}(H^{2,1},H^{1,2})$ の単射性を使う。
通常のHKRとの差も確認する。局所体積形式 $\mu_i$ を選び、
$\iota_{D_i}\mu_i=d(\iota_\pi\mu_i)$ によりベクトル場 $D_i$ を定める。
$\mu_j=f_{ij}\mu_i$ なら $D_j-D_i=\pm\pi^\sharp d\log f_{ij}$ となる。
従って $\pi\lrcorner c_1(K_X)$ は $H^1(T_X)$ で零である。
残りのTodd補正は $H^2(\OO_X)=0$ に入るので、双ベクトルに対して両HKRは一致する。
その $H^{2,1}=H^1(\Omega_X^2)$ への作用は
$H^1(\Omega_X^2)\to H^1(\OO_X)=0$ であり、$\gamma$ の単射性から結論する。
この局所計算にPoisson条件は不要である。
\end{proof}

\subsection{残余圏を固定した形式ファイバー}

\begin{theorem}[固定した残余圏の形式ファイバー]\label{thm:fibre}
上で構成した $\Phi$ の、未変形の残余圏の上の古典的形式ファイバーは非標準的に
\[
 \Spf\CC[[a_{ij}]]/(q_{ijkl})
\]
と同型である。特に次元は7である。
\end{theorem}

\begin{proof}
環準同型を $s_i\mapsto f_i(t,a)$ と書く。
$\det(\partial f_i/\partial t_j)(0)\ne0$ なので、
$t'_i=f_i(t,a),\ a'_{ij}=a_{ij}$ は形式的な座標変換を与える。
$a$ を固定するためPlückerイデアルも保存する。
新しい座標で $\Phi$ は $t'$ への射影となり、結論を得る。
\end{proof}

$\Phi$ の接写像が双ベクトル方向を殺すことは、
任意の純Poisson族で残余圏が全次数で一定であることを意味しない。
固定残余圏の条件は一般に幾何学的パラメータ $t$ の高次補正を必要とし得る。
幾何学的sliceの座標を用いると、この補正は $t(a)=O(a^2)$ から始まる。

\subsection{明示的モデルの残余圏}

$\lambda^6\ne-27$ とする。\eqref{eq:family} のAS Gorenstein代数に
Orlovの定理\cite[Theorem 2.5]{Orlov}を適用すると半直交分解
\[
 D^b(\qgr A_\lambda)
 =\langle\mathcal K_\lambda,\OO,\OO(1)\rangle,\qquad
 \mathcal K_\lambda\simeq\Dsg^{\mathrm{gr}}(A_\lambda)
\]
を得る。埋め込みの選択によりshift・mutationの規約は変わる。
第6節の孤立特異点性と $\omega_A=A(-2)$ に、
Ueyamaのstable Cohen--Macaulay圏のSerre公式\cite{Ueyama}を適用すると、
残余圏のSerre関手は
\[
 S_{\mathcal K_\lambda}=(-2)[3]
\]
となる。ここで括弧内は特異圏側のgrade shiftを表す。
中心的な三次斉次元のmatrix factorizationの周期性 $(3)\simeq[2]$ から
\[
 S_{\mathcal K_\lambda}=(1)[1],\qquad
 \boxed{S_{\mathcal K_\lambda}^{\,3}=[5]}
\]
を得る。従って通常のcubic threefoldと同じ分数Calabi--Yau次元 $5/3$ を持つ。
grade shiftは $S_{\mathcal K_\lambda}[-1]$ で既に決まるので、
これだけを追加しても、全体圏への埋め込みやgluingの情報を補えない。
また、この関手等式だけから $\mathcal K_\lambda$ がある可換cubicの
Kuznetsov成分と同値であることは従わない。次節では別の具体的な比較により同定する。

\section{Fermat族の残余圏とHesse族の比較}

\subsection{主張と、二次類の規約}

\[
 Y_s=V(x^3+y^3+z^3+u^3+v^3+sxyz)\subset\PP^4
\]
とおく。偏微分の計算から、$Y_s$ は $s^3\ne-27$ のとき滑らかである。
以下のdg圏は常にpretriangulated hullとidempotent completionを取った
Morita圏として扱う。

\begin{theorem}[Hesse比較]\label{thm:hesse}
$\lambda^6\ne-27$ なら、dg Morita同値
\[
 \mathcal K_\lambda\simeq\Ku(Y_{\lambda^2})
\]
が存在する。また $R=\CC[[h]]$ 上では、この同値を
$h=0$ の通常のOrlov同一視に還元する同値として構成できる。
従って、$X_0=Y_0$ のJacobian環による同一視
\[
 \HH^2(\Ku(X_0))\simeq
 \left(\frac{\CC[x,y,z,u,v]}{(x^2,y^2,z^2,u^2,v^2)}\right)_3
\]
において、最初の非零の変形係数は
\[
 \boxed{\kappa_2=[xyz]\ne0.}
\]
ここで $\kappa_2$ は $h^2$ の係数を意味する。
二階微分そのものを使う規約では $2[xyz]$ となる。
\end{theorem}

証明は楕円曲線の導来圏だけでなく、grade shiftと $(3)\simeq[2]$ を比較する。
点曲線の式だけからこの命題を結論することはできない。

\subsection{楕円曲線のblockと、次数シフト}

まず体上で $C=S_\lambda/(G_\lambda)$ とする。
前述のtwisted homogeneous coordinate ringの同一視と、
Etingof--Ginzburg\cite[Theorem 3.6.2]{EG}、Orlovの定理から
\[
 \MF^{\mathbb Z}(S_\lambda,G_\lambda)
 \simeq\Dsg^{\mathrm{gr}}(C)\simeq D^b(E_\lambda)
\]
を得る。$\lambda=0$ では通常の平面三次曲線のOrlov同値である。
この同値は以下のwindow構成によりdg関手で実現する。

一般に $C=B(E,L,\sigma)$、$\deg L=3$、$\sigma=t_p$ を楕円曲線の平行移動とする。
右加群の規約
\[
 \Gamma_*(F)_n=H^0(E,F\otimes L_n),\qquad
 L_n=L\otimes\sigma^*L\otimes\cdots\otimes(\sigma^{n-1})^*L
\]
を採用すると、tailsの次数シフトは
\[
 H(F)=\sigma_*(F\otimes L)
\]
である。特異圏側の次数シフトを $\rho$ と書く。
適切なOrlov windowでは
\begin{equation}\label{eq:rotation}
 \rho=T_{\OO_E}\circ H.
\end{equation}

この公式は非可換の場合にもwindowから直接に従う。
$q:D^b(\operatorname{gr}C)\to\Dsg^{\mathrm{gr}}(C)$ とし、
$\Psi_i=qR\Gamma_{\ge i}$ とおく。Gorenstein parameterは零なので各 $\Psi_i$ は同値である。
内部次数の切断から $(1)\Psi_1=\Psi_0H$、
$q(R)=\Psi_0(\OO_E)[1]$、および三角
\[
 \Psi_1F\longrightarrow\Psi_0F\longrightarrow
 \RHom(\OO_E,F)\otimes q(R)
\]
を得る。楕円曲線のSerre双対を使うと
$\Psi_1=\Psi_0T_{\OO_E}^{-1}$ となり、\eqref{eq:rotation}を得る。
体の場合は $R$ をその体と読む。

\subsection{平行移動の消去を明示する}

$\rho_0=T_{\OO_E}\circ(-\otimes L)$ を可換な平面三次曲線のgrade shiftとする。
平行移動は $\OO_E$ を保存するので、\eqref{eq:rotation}から
\[
 \rho_{\NC}=t_{-p}^*\rho_0.
\]
原点を $o\in E$ とし $P_q=\OO_E(q-o)$、$A_q=(-\otimes P_q)$ とおく。
楕円曲線の対角はCartier因子なので、spherical twistのFourier--Mukai核は
\[
 K_{T_{\OO_E}}=I_\Delta[1],\qquad
 K_{T_{P_q}}=(P_q^{-1}\boxtimes P_q)\otimes I_\Delta[1].
\]
Theorem of the square、または両射影の各ファイバーへの制限から
\[
 (P_q^{-1}\boxtimes P_q)\otimes I_\Delta
 \simeq I_{\{y=x-q\}}
\]
となる。従って $T_{P_q}=t_q^*T_{\OO_E}$ であり、
\begin{equation}\label{eq:conjugacy}
 A_q\rho_0A_q^{-1}=t_q^*\rho_0,
 \qquad \rho_{\NC}=A_{-p}\rho_0A_{-p}^{-1}.
\end{equation}
従って平行移動のパラメータは、次数零の直線束による共役で取り除ける。

この比較では三次周期性 $\rho^3\simeq[2]$ も保存する。
二つの周期性同型の差は $\HH^0(D^b(E))=R$ の単元である。
体 $\CC$ 上ではその三乗根、形式基底上では特殊ファイバーで1となる三乗根を取り、
共役の同型を正規化する。高次のcoherenceの差も、各小拡大の核 $I$ に対する
$H^{>0}(\mathbb Z/3,I)=0$ により帰納的に消去できる。

\subsection{正則DVR上の比較とdg基底変換}

形式的な二次類を読むには、各複素ファイバーでの同値だけでは不足する。
そこで $R=\CC[[h]]$ 上で構成する。ここでは
\[
 E_R=V(x^3+y^3+z^3+h^2xyz)\subset\PP^2_R
\]
は滑らかな楕円曲線である。$S_R$ の関係行列は $E_R$ 上で階数2となる。
その核は $\sigma_R:E_R\to E_R$ を定め、転置行列から逆写像を得る。
$\sigma_0=\id$ であり、$\sigma_R$ は平行移動である。
例えばSklyaninの標準的規約では $p_h=[1:-1:-h]$、原点は $[1:-1:0]$ である。
符号の逆の規約を用いても \eqref{eq:conjugacy} の結論は同じである。

関係行列から $S_R\to B(E_R,L,\sigma_R)$ を得る。
$G_h$ は一般ファイバーで零に写るので、各次数の $R$-torsion freenessから全体で零に写る。
特殊ファイバーは通常の三次曲線の斉次座標環である。
次数 $n>0$ の両辺は階数 $3n$ の有限自由 $R$-加群なので、Nakayamaにより
\begin{equation}\label{eq:relative-thcr}
 C_R=S_R/(G_h)\simeq B(E_R,L,\sigma_R).
\end{equation}

相対Koszul複体とその双対は、各内部次数で有限自由 $R$-複体である。
$h=0$ で多項式環のKoszul複体になるので完全であり、
\[
 \RHom_{S_R}(R,S_R)=R(3)[-3],\qquad
 \RHom_{C_R}(R,C_R)=R[-2]
\]
を得る。$\operatorname{gldim}R=1$ と相対bimodule分解から
$S_R$ のgraded global dimensionは4以下、$C_R$ の両側graded injective dimensionは有限となる。

Noetherian性も相対的に確認できる。$\operatorname{gr}_hS_R=\CC[x,y,z][h]$ の
斉次イデアルの有限な生成元を持ち上げる。固定した内部次数での係数の修正は
有限自由 $R$-加群内で収束するので、任意の斉次イデアルは $S_R$ 自身で有限生成である。
非斉次イデアルには最高内部次数のイデアルを使い、次数を下げる有限回の帰納法を適用する。
これで通常の両側Noetherian性を得る。環全体が $h$-進完備であることは仮定しない。

従ってBurke--Stevenson\cite[Definition 5.1, Lemma 6.2, Theorem 6.4(3)]{BS}を
正則DVR $R$ 上で適用でき、$a=0$ のOrlov windowを得る。
非可換MFとの比較は次のように確認する。
$C_R$ のGorenstein-projective加群 $M$ についてchange of ringsから
$\Ext^{>1}_{S_R}(M,S_R)=0$ となる。
graded localな $S_R$ の有限最小分解を双対化すると
$\operatorname{pd}_{S_R}M\le1$ なので、
\[
 0\to P_1\xrightarrow\varphi P_0\to M\to0
\]
がある。$G_hM=0$ より $G_h:P_0\to P_0(3)$ を $\psi:P_0\to P_1(3)$ に持ち上げれば
$\varphi\psi=G_h$、$\psi\varphi=G_h$ となる。
分解の比較は安定射とMFのホモトピー射を同一視する。
これは\cite{CCKM}の議論をgraded localなDVR上で行ったものである。

dg構造と基底変換についても構成を固定する。
複体のdg圏、perfect複体によるdg quotient、complete resolutionを用いれば、
上のcokernel・window関手はdg準関手となる。
windowの逆関手 $qR\Gamma_{\ge i}$ は、内部次数の切断とderived global sectionからなる。
\eqref{eq:relative-thcr}の下で各成分は
$R\Gamma(E_R,F\otimes L_n)$ であり、$E_R\to\operatorname{Spec}R$ がsmooth properで
$F$ がperfectなので、任意のderived base changeと可換である。
また $C_R$ は $R$ 上flatで $\operatorname{gldim}R=1$ なので、
$-\otimes_R^{\mathbf L}\mathbb C$ は $D^b(\operatorname{gr}C_R)$ を
$D^b(\operatorname{gr}C_0)$ に送り、perfect複体をperfect複体に送る。
従って特異圏の間の特殊化関手が定まる。
内部次数の切断は成分ごとの操作であり、この関手によってwindowを比較できる。
両windowが逆であることは三角と随伴射で実現され、その射も同時に基底変換する。
従って
\[
 \MF^{\mathbb Z}(S_R,G_h)\simeq\Perf(E_R)
\]
は $h=0$ のOrlov同値と両立するdg Morita同値である。
Artin商へはこのdg圏をderived base changeする。
特異なArtin基底の\emph{絶対}特異圏を直接使うことはしない。
より具体的には、$\Perf(E_R)$ の有限個のample twistから生成子を選び、
そのwindow像の自己準同型dg代数を取ればよい。
その特殊化は $\Perf(E_0)$ の生成子の自己準同型代数となるので、
基底変換後のMFのMorita圏に生成子が欠落することもない。

相対版の \eqref{eq:conjugacy} では $p_0=o$ なので $A_{-p_0}=\id$。
従って標識を保存して、$\MF(S_R,G_h)$ を
$\MF(R[x,y,z],x^3+y^3+z^3+h^2xyz)$ と、次数シフト込みで比較できる。

\subsection{中心変数を加えるThom--Sebastiani型補題}

次に、必要な非可換テンソル積の主張を証明する。
可換の場合の対応する定理はBallard--Favero--Katzarkov\cite[Theorem 5.15, Corollary 5.18]{BFK}
だが、その定理の仮定をそのまま非可換代数へ拡張して引用することはしない。

$S$ を上の代数、$g$ を中心三次元とし、$u,v$ を中心変数とする。
\[
 \Gamma=\mathbb Z^3/\langle3e_1-3e_2,3e_1-3e_3\rangle,
 \qquad w=3e_1=3e_2=3e_3
\]
で次数を付け、$x,y,z$ の次数を $e_1$、$u$ を $e_2$、$v$ を $e_3$ とする。

\begin{lemma}\label{lem:ts}
外部テンソル積はdg Morita同値
\[
 \MF^{\mathbb Z}(S,g)\otimes_R^{\mathbf L}
 \MF^{\mathbb Z}(R[u],u^3)\otimes_R^{\mathbf L}
 \MF^{\mathbb Z}(R[v],v^3)
 \simeq\MF^\Gamma(S[u,v],g+u^3+v^3)
\]
を与える。$R=\CC$ または上の正則DVRとする。
\end{lemma}

\begin{proof}
有限自由MFの奇微分の外部テンソル和は、Koszul符号を付けると
$(d_1+d_2+d_3)^2=g+u^3+v^3$ を満たす。
完全忠実性はHom複体の直接計算である。
次数 $n$ のHomは
\[
 \bigoplus_{i+j+k=n}
 \operatorname{Hom}^i(P,P')\otimes_R
 \operatorname{Hom}^j(Q,Q')\otimes_R
 \operatorname{Hom}^k(T,T')
\]
に分解し、微分もテンソル積微分になる。
三つの $[2]$ と三次内部シフトを同一視するのが $\Gamma$ の関係式である。
有限自由MFと非負次数の代数では、固定した $n$ に寄与する分解は有限なので
無限積・完備化の問題は生じない。
$u^3,v^3$ の因子は $\Perf(RA_2)$ という有限自由なtilting模型を持つので、
相対的なテンソル積はこの模型でderived tensorとして計算できる。

生成性を示す。$A=S[u,v]/(g+u^3+v^3)$、$C=S/(g)=A/(u,v)$ とする。
MFの微分を $u,v$ で微分すると $u^2,v^2$ の作用は特異圏で零となる。
従って任意の対象 $N$ は $K(u^2,v^2;N)$ のshift・twistの直和因子である。
このKoszul複体のコホモロジーは $(u^2,v^2)$ に殺され、
$(u,v)^3\subset(u^2,v^2)$ より有限filtrationの商は $C$-加群となる。

$u,v$ は $A$ の中心正則列なので、$C$-free加群の制限は $A$ 上perfectである。
従って制限は $\Dsg(C)\to\Dsg(A)$ に降りる。
$C$-加群 $M$ の安定化を $P$ とすると、その制限の安定化は
\begin{equation}\label{eq:stabilization}
 P\boxtimes(u,u^2)\boxtimes(v,v^2)
\end{equation}
となる。これは次の分解から直接に分かる。
$M$ の長さ1の $S$-分解を $S[u,v]$ に延長し、$u,v$ のKoszul分解とテンソルする。
その上で $g$ の零ホモトピーと、$u^3=u\cdot u^2$、$v^3=v\cdot v^2$ の
Koszul零ホモトピーを足して周期化すると、\eqref{eq:stabilization}の微分そのものになる。
異なる因子の奇作用は反可換なので追加の高次ホモトピーは不要である。
これは制限加群の自由分解の周期的末尾を与える。

一変数について行列で確認しておく。
$P=(P_1\xrightarrow\varphi P_0\xrightarrow\psi P_1(3))$、$M=\operatorname{coker}\varphi$
とすると、$uM=0$ とした $M$ の $S[u]$-分解は
\[
 P_1(-1)\xrightarrow{\binom{-u}{\varphi}}
 P_1\oplus P_0(-1)\xrightarrow{(\varphi\ \ u)}P_0.
\]
$g+u^3$ の零ホモトピーは $s_0=(\psi,u^2)^t$、$s_1=(-u^2,\psi)$ で、
$ds+sd=g+u^3$、$s^2=0$ を満たす。周期化した末尾の行列は
\[
 \alpha=\begin{pmatrix}\varphi&u\\-u^2&\psi\end{pmatrix},
 \qquad
 \beta=\begin{pmatrix}\psi&-u\\u^2&\varphi\end{pmatrix},
 \qquad \alpha\beta=\beta\alpha=(g+u^3)I.
\]
この周期化が $M$ を分解することも確認できる。
元の分解を $A=S[u]/(g+u^3)$ にテンソルした複体のホモロジーは
次数0の $M$ と次数1の $M(-3)$ のみである。
周期化を周期の個数で濾過すると、$s$ が隣接する二つの $M$ を同型で結ぶ。
$ds=(g+u^3)$ によるconnecting mapでこの同型が分かる。
従ってそれらは対ごとに消え、次数0の $M$ だけが残る。
各全次数で濾過は有限なので、このspectral sequenceは収束する。
これで\eqref{eq:stabilization}は単なる因子化の恒等式でなく、制限の安定化を表す。

以上から任意の対象が外部積の像で厚く生成される。
$\Gamma$-次数の $C$-加群は $\Gamma/\mathbb Ze_1$ の9個の剰余類に分かれるが、
第二・第三因子の次数シフトがこれらをすべて与える。
直和因子を加えれば主張を得る。
\end{proof}

\subsection{通常の整数次数への降下}

$\Gamma\to\mathbb Z$、$e_i\mapsto1$ の核は
$T\simeq(\mathbb Z/3)^2$ である。
$\Gamma$-次数付きMFは、通常の次数付きMFに $H=\operatorname{Hom}(T,\CC^*)$ の
同変構造を付けた圏である。具体的には $H$ は $u,v$ を独立に三乗根倍する。
双対群の作用で再び同変化すると元の圏を回復する
\cite[Theorem 4.2]{Elagin}。これは有限群の平均化によるdg構成でも成立し、
$|H|$ が可逆な $R$ 上でも同じ証明でよい。

補題\ref{lem:ts}の左辺では、この双対作用はgrade shiftの差
\[
 \rho_2\rho_1^{-1},\qquad\rho_3\rho_1^{-1}
\]
で与えられる。三乗が恒等になる同型は各因子の三次周期性から得る。
\eqref{eq:conjugacy}は第一因子の $\rho_1$ と周期性を保存するので、
テンソル積の同値はこの群作用を保存し、通常の整数次数のMF圏へ降りる。
これで
\[
 \MF^{\mathbb Z}(S_R[u,v],G_h+u^3+v^3)
 \simeq
 \MF^{\mathbb Z}(R[x,y,z,u,v],x^3+y^3+z^3+h^2xyz+u^3+v^3)
\]
を得る。双方にGorenstein parameter 2のOrlov分解を適用すれば
命題\ref{thm:hesse}の圏同値を得る。形式的には全構成が $h=0$ の同一視を保つ。
従って二次係数は $[xyz]$ であり、この単項式はFermatのJacobianイデアルに属さない。

\subsection{dg smoothnessへの帰結}

\begin{theorem}\label{thm:dg-smooth}
$\lambda^6\ne-27$ なら $D^b(\qgr A_\lambda)$ の標準的dg enhancementは
smooth properである。
\end{theorem}
\begin{proof}
命題\ref{thm:hesse}から $\mathcal K_\lambda$ はsmooth properであり、
他の二成分も $\Perf(\CC)$ なのでsmooth properである。
$\qgr A_\lambda$ は $\chi$、有限コホモロジー次元と第6節の大域次元3を持つので、
coherent対象間のtotal Extは有限次元である。
従って半直交分解のgluing bimoduleはpointwise properである。

smoothな圏上のpointwise proper加群はperfectである。
実際、perfectな対角bimoduleを有限個のrepresentable bimoduleから構成し、
その構成に $M\otimes^{\mathbf L}-$ を作用させればよい。
従って各gluing bimoduleはperfectであり、Orlovのgluing定理
\cite[Proposition 3.8, Theorem 3.25]{OrlovGlue}から全体がsmoothとなる。
propernessは既に確認した。
\end{proof}

これは全体圏が $D^b(Y_{\lambda^2})$ と同値であるという主張ではない。
上で同定したのは残余成分であり、例外対とのgluingは別のデータである。

\section{二次類の解釈と、残余圏を固定する明示的補正}

\subsection{対称性による独立の検査}

$\omega^3=1$、$\omega\ne1$ とする。
この族は $u\mapsto\omega u$ と $v\mapsto\omega v$ を持ち、
半線形な対称性 $(h,x)\mapsto(\omega h,\omega x)$、
$(h,x,y)\mapsto(-h,y,x)$ も持つ。
Fermat Jacobian環の次数3は5変数のsquare-freeな三次単項式10個を基底とする。
最初の二対称性の不変部分は $\CC[xyz]$ のみである。

幾何学的形式変形基底との標識付き同一視を取ると、対称性を保つこの族の像は
$Y_{s(h)}$ と書け、半線形対称性から
\[
 s(\omega h)=\omega^2s(h),\qquad s(-h)=s(h)
\]
となる。従って $s(h)=h^2f(h^6)$ であり、許される次数は $2,8,14,\ldots$ である。
対称性だけでは $f(0)$ を決定できない。
命題\ref{thm:hesse}が、通常のmarked Hesse座標で $s(h)=h^2$ を与える。

\subsection{中心式の可換化が二次類を与えない理由}

対称化した中心式は
\[
 G_h=x^3+y^3+z^3-\frac{h^2}{3}\sum_{\tau\in S_3}\tau(xyz)
\]
なので、文字を可換として並べれば二次項は $-2h^2xyz$ になる。
しかしこの操作は周囲の代数の乗法の変形を捨ててしまう。
残余圏の変形はその乗法も含むため、正しい係数は命題\ref{thm:hesse}の
$+h^2xyz$ である。

補助的な楕円曲線の $j$-不変量も計算できる。
\[
 j(E_s)=\frac{s^3(216-s^3)^3}{(s^3+27)^3},\qquad
 j(E_{h^2})=512h^6+O(h^{12}).
\]
これは $s=0$ の楕円曲線の自己同型による粗モジュライ写像の分岐を反映する。
$j$ の最初の変化が6次であることは、marked変形の2次類が零であることを意味しない。
添付コードはHesse式をWeierstrass式へ変換してこの有理式を厳密検証する。

\subsection{固定した残余圏を持つ形式族}

$t=(1-h^3)^{1/3}\in\CC[[h]]$ を $t(0)=1$ の分枝とし、
\[
 S'_h=\frac{\CC[[h]]\langle x,y,z\rangle}
 {(xy-tyx-hz^2,\ yz-tzy-hx^2,\ zx-txz-hy^2)}
\]
とする。そのSklyaninパラメータは $(a,b,c)=(1,-t,-h)$ であり、
\[
 a^3+b^3+c^3=1-t^3-h^3=0.
\]
従って一般ファイバーの点曲線は固定したFermat曲線
$E_0:x^3+y^3+z^3=0$ であり、中心切断に対応する相対楕円曲線も
$E_0\times\operatorname{Spf}\mathbb C[[h]]$ と取れる。
ただし $h=0$ で $S'_0$ 自身の点スキームは $\PP^2$ である。

標準中心元\cite[Lemma 2.19]{WWY}を正規化すると
\[
 G'_h=(1+h^3)(x^3+z^3)+(1-2h^3)y^3-3h^2xyz
\]
を得る。具体的には同文献の中心元 $g$ に上のパラメータを代入し、
$yxz=t^{-1}(xyz-hz^3)$ を用いると $g=-hG'_h$ となる。
$S'_h$ は $h$-torsion freeなので $G'_h$ は中心元である。
その平坦性は、一般ファイバーのSklyanin Hilbert級数と特殊ファイバーの多項式環から、
有限生成 $\CC[[h]]$-加群の構造定理を各次数に適用して従う。

\[
 A'_h=S'_h[u,v]/(G'_h+u^3+v^3)
\]
とおく。monicな $u$ の式からこれも次数ごとに平坦であり、
前節と同じ次数シフト・テンソル積の比較により
\[
 \mathcal K(A'_h)\simeq\Ku(X_0)\otimes\CC[[h]]
\]
を得る。一方、$t=1-h^3/3+O(h^6)$ なので、元の族と半古典括弧は同じである。
元の中心元を $xyz$ 順序で書くと
\[
 G_h=(1+h^3)(x^3+z^3)+(1-h^3)y^3-2h^2xyz.
\]
従って両族の関係式を $h^3$ で比較して
\[
 G'_h-G_h\equiv-h^2xyz\pmod{h^3}
\]
となる。これは $\kappa_2=[xyz]$ を消す幾何学的補正を明示する。
高次には中心性を保つため乗法も補正しており、
元の代数のまま中心でない $h^2xyz$ だけを引いているのではない。

\section{6個の例外ファイバーは射影的に特異である}\label{sec:exceptional}

\begin{theorem}\label{thm:exceptional}
全ての $\lambda\in\CC$ について
\[
 \operatorname{gldim}(\qgr A_\lambda)=
 \begin{cases}
 3&\lambda^6\ne-27,\\
 \infty&\lambda^6=-27.
 \end{cases}
\]
ここで大域次元はcoherent対象間のExtによるホモロジー次元である。
\end{theorem}

\subsection{Fourier変換と中心切断の三角形}

$\omega^2+\omega+1=0$、$\lambda_0=\omega-\omega^2$ とすると $\lambda_0^2=-3$ である。
前節の対称性 $h\mapsto\omega h$、$h\mapsto-h$ は6個の例外ファイバーを同型にする。
従って $\lambda_0$ だけを考えればよい。
\[
 P=x+y+z,\quad Q=x+\omega y+\omega^2z,\quad R=x+\omega^2y+\omega z
\]
とおく。ここで $R$ はこの節の第三生成元を表す。
直接代入により
\[
 S_{\lambda_0}\simeq
 \CC\langle P,Q,R\rangle/(P^2+2RQ,\ Q^2+2PR,\ R^2+2QP).
\]
二次関係から $P^3=Q^3=R^3$。
一方、自由代数内のFourier展開から
\[
 P^3+Q^3+R^3
 =3\left(x^3+y^3+z^3+\sum_{\tau\in S_3}\tau(xyz)\right)
 =3G_{\lambda_0}
\]
なので、厳密に $G_{\lambda_0}=P^3$ である。

量子多項式環
\[
 T=\CC\langle P,Q,R\rangle/(PR+2RP,\ QP+2PQ,\ RQ+2QR)
\]
と位数3の自己同型 $\tau(P)=R,\tau(Q)=P,\tau(R)=Q$ をとる。
Zhang積 $a*b=a\tau^{\deg a}(b)$ で計算すると
$S_{\lambda_0}=T^\tau$、$P*P*P=-2PQR$ であり、
\[
 C:=S_{\lambda_0}/(G_{\lambda_0})\simeq(T/(PQR))^\tau.
\]
Zhang同値\cite[Theorem 1.1]{Zhang}によりgraded加群圏とtailsは保存される。

$T/(PQR)$ は三角形 $E=V(PQR)\subset\PP^2$ のtwisted homogeneous coordinate ringである。
$q=-2$ とすると、その自己同型は各成分で
\[
 \begin{array}{c|c}
 P=0 &[Q:R]\mapsto[qQ:R]\\
 Q=0 &[P:R]\mapsto[P:qR]\\
 R=0 &[P:Q]\mapsto[qP:Q]
 \end{array}
\]
となり、節点では一致する。そのgraph方程式が $T$ の多重線形関係である。
各次数の切断は、三つの頂点を検出する純冪と各辺の混合単項式で生成され、
次数 $n\ge1$ の次元は両辺とも $3n$ となる。これでTHCRとの同型を直接確認できる。
ample sequenceのmultidegreeは $(n,n,n)$ であり、
特異な射影スキームにも適用できるTHCRの定理\cite{AV,SVdB}から
\[
 \qgr C\simeq\operatorname{coh}E
\]
を得る。

$C/PC$ は辺 $L=V(P)$ の構造層に対応する。節点での完備局所模型は
\[
 B=\CC[[a,b]]/(ab),\qquad \OO_L=B/(a).
\]
この加群の自由分解は $a,b$ の掛け算が交互に現れる無限周期分解であり、
剰余体とテンソルすると全ての正次数のTorが非零となる。
従って $\OO_L$ はperfectでない。
$\pi_C C(n)$ は直線束なので、$\pi_C(C/PC)$ は
\[
 D^b(\qgr C)/\operatorname{thick}\{\pi_C C(n)\}
\]
で非零である。この段階で滑らかな楕円曲線に対するOrlov同値は用いていない。

\subsection{特異性を検出する階数4の行列因子化}

$S=S_{\lambda_0}$ 上の階数1の因子化 $(P,P^2)$ を使い、
\[
 Q=(P,P^2)\boxtimes(u,u^2)\boxtimes(v,v^2)
\]
を作る。これは $S[u,v]$ 上の $W=P^3+u^3+v^3$ のgraded MFで、各parityの階数は4である。
外積作用 $\varepsilon_i$ と縮約 $\iota_i$ により、全微分は
\[
 d=P\varepsilon_1+P^2\iota_1+
 u\varepsilon_2+u^2\iota_2+v\varepsilon_3+v^2\iota_3
\]
と書ける。中心性から異なる因子は反可換し、$d^2=W I$ となる。
内部次数シフトは通常の三次MFのテンソル積に従う。
この等式は添付コードで $8\times8$ 行列の全64成分について検証する。

$N=\operatorname{coker}Q$ とする。$(u,v)$ は $A=A_{\lambda_0}$ の中心正則列なので、
導来制限
\[
 -\otimes_A^{\mathbf L}C:D^b(\qgr A)\to D^b(\qgr C)
\]
はTor振幅2を持ち、$\pi_AA(n)$ を $\pi_CC(n)$ に送る。
graded torsion加群のTorもtorsionなので、tailsへの降下も定義される。
従って両側のline objectが生成する厚い部分圏で割った商への関手を得る。

$Q$ を $u=v=0$ に特殊化すると、$(P,P^2)$ のshift・twistを変えた4個の直和となる。
これは導来制限の像の周期的自由分解の末尾を表す。
実際、$N$ の自由分解を $C$ とテンソルすれば、Tor振幅2により十分高次の末尾は
この特殊化と一致し、有限の初期部分は上の商圏ではperfectな差にしかならない。
前小節の非perfect性からその像は非零であり、従って
\[
 \pi_A N\notin\operatorname{thick}\{\pi_AA(n)\}.
\]

もし $\qgr A$ の大域次元が有限の $d$ なら、任意のcoherent対象 $M$ の
line objectの有限和による表示を $d+1$ 段まで取り、最後の核を $K$ とする。
対応する有限複体の連結類は $\Ext^{d+1}(M,K)=0$ に入る。
従って $M$ はその有限複体の直和因子であり、上の厚い部分圏に入る。
これは $N$ と矛盾する。以上で命題\ref{thm:exceptional}を得る。

Zhang twistは中心切断 $C$ の解析にだけ用いた。
中心変数 $u,v$ の付加がそのtwistと可換するとは仮定していない。

\section{二段階の貼り合わせと単一対象の変形}\label{sec:gluing-objects}

\subsection{二つのgluing object}

$V=H^0(X,\OO_X(1))$、$\dim V=5$ とし、評価列
\[
 0\longrightarrow Q\longrightarrow V\otimes\OO_X
 \longrightarrow\OO_X(1)\longrightarrow0
\]
で $Q=\Omega^1_{\PP^4}(1)|_X$ を定める。
$Q=L_{\OO_X}\OO_X(1)[-1]$ は例外対象であり、mutationにより
\[
 \mathcal D=\Perf(X)=\langle\mathcal K,Q,\OO_X\rangle,\qquad
 \mathcal B=\OO_X^\perp=\langle\mathcal K,Q\rangle
\]
を得る。$\mathcal K=\Ku(X)$、包含を $i:\mathcal K\to\mathcal D$、
$j:\mathcal B\to\mathcal D$ と書き、
\[
 E=i^!Q\in\mathcal K,\qquad G=j^!\OO_X\in\mathcal B
\]
とおく。$E$ は $\mathcal K$ に $Q$ を加えるgluing object、
$G$ は $\mathcal B$ に最後の $\OO_X$ を加えるgluing objectである。
$E$ の構成と次の三角はFeyzbakhsh--Liu--Zhang
\cite[Section 3.2, equations (9)--(12)]{FLZ}にも記載される。
\[
 \OO_X(-1)[1]\longrightarrow E\longrightarrow Q
 \longrightarrow\OO_X(-1)[2]. \tag{E}\label{eq:E-triangle}
\]
以下、随伴の同一視および必要な1次元のtrace線の自明化を固定する。
例えば上の三角は、選択なしには
$\OO_X(-1)[1]\otimes\Ext^2(Q,\OO_X(-1))^*$ を左端に持つ。
従って以下の裸の $\Lambda^2V$ との同一視を、選択なしの
$GL(V)$-同変な同一視と解釈する必要はない。

\begin{theorem}[二つの対象の自己Ext]\label{thm:gluing-ext}
上の規約の下で
\[
 \Ext^r_{\mathcal K}(E,E)=
 \begin{cases}\CC&r=0,\\ \Lambda^2V&r=1,\\0&\text{その他},\end{cases}
 \qquad
 \Ext^r_{\mathcal B}(G,G)=
 \begin{cases}\CC&r=0,\\ \Sym^2V&r=2,\\0&\text{その他}.\end{cases}
\]
従って $E$ の古典的な標識付き対象変形の同値類は滑らかな10次元の形式空間をなし、
$G$ の対象変形は剛的で、障害の入り得る空間は15次元である。
\end{theorem}

\begin{proof}
随伴から $\RHom_{\mathcal K}(E,E)=\RHom_{\mathcal D}(E,Q)$ である。
評価列を $\OO_X(1)$ でtwistすると
\[
 \RHom(\OO_X(-1),Q)=R\Gamma(Q(1))
 =\ker(V\otimes V\to\Sym^2V)=\Lambda^2V
\]
を次数零に得る。$Q$ の例外性と\eqref{eq:E-triangle}に
$\RHom(-,Q)$ を適用すれば三角
\[
 \CC\longrightarrow\RHom(E,E)\longrightarrow
 \Lambda^2V[-1]\longrightarrow\CC[1]
\]
を得るので、最初の計算が従う。

一方、Serre関手 $S_{\mathcal D}=(-2)[3]$ により
\[
 \mathcal D=\langle\OO_X(-2)[3],\mathcal B\rangle
\]
と回転できる。
$\RHom(\OO_X,\OO_X(-2)[3])=\CC$ なので、右mutationは
\[
 \OO_X(-2)[2]\longrightarrow G\longrightarrow\OO_X
 \longrightarrow\OO_X(-2)[3] \tag{G}\label{eq:G-triangle}
\]
を与える。随伴と $\RHom(-,\OO_X)$ により
\[
 \CC\longrightarrow\RHom(G,G)\longrightarrow
 H^0(\OO_X(2))[-2]\longrightarrow\CC[1]
\]
を得て、二番目の計算も従う。

対象の変形は自己準同型dg代数のMaurer--Cartan理論で制御される
\cite[Proposition 6.1]{ELOI}\cite[Theorems 7.2, 11.8]{ELOII}。
$E$ のstrictly unitalな最小 $A_\infty$ 模型は次数0の単位と次数1しか持たない。
正次数入力だけを持つ $m_n$ の出力は次数2なので全て零であり、
次数零のscalar gaugeも可換Artin基底上では自明に作用する。
従って
\[
 \Def_E\simeq\Spf\CC[[z_1,\ldots,z_{10}]].
\]
これは対象の同値類の形式基底であり、scalar自己同型を含むstack全体ではない。
\end{proof}

\subsection{非可換接方向がこの対象を動かすことの証明}

半直交分解 $\mathcal T=\langle\mathcal A,F\rangle$ に対し、
$F$ が例外対象で $E_F=i^!F$ なら、Hochschild gluing triangleは
\begin{equation}\label{eq:gluing-triangle}
 CC^\bullet(\mathcal T)\longrightarrow
 CC^\bullet(\mathcal A)\oplus CC^\bullet(\CC)
 \longrightarrow\RHom(E_F,E_F)
 \longrightarrow CC^\bullet(\mathcal T)[1].
\end{equation}
これは\cite[Theorem 7.7, Remark 7.8]{Kuz}の完全列のdg模型である。
次数零での第二の射は、二つのscalarの差の作用である。

$\mathcal B=\langle\mathcal K,Q\rangle$ に適用すると、命題
\ref{thm:gluing-ext}と第7節の計算から
\[
 \HH^1(\mathcal B)=0,\qquad
 0\longrightarrow\Ext^1(E,E)\longrightarrow\HH^2(\mathcal B)
 \longrightarrow\HH^2(\mathcal K)\longrightarrow0,\qquad
 \HH^3(\mathcal B)=0.
\]
特に $\dim\HH^2(\mathcal B)=20$ である。
次に $\mathcal D=\langle\mathcal B,\OO_X\rangle$ に適用すると、
$\Ext^1(G,G)=0$ から
$\HH^2(\mathcal D)\to\HH^2(\mathcal B)$ は単射であり、
両辺が20次元なので同型となる。また
\[
 \Ext^2(G,G)\xrightarrow{\sim}\HH^3(X).
\]
二つの制限射の合成 $\HH^2(\mathcal D)\to\HH^2(\mathcal K)$ の核は、
補題\ref{lem:bivector-zero}から $H^0(\Lambda^2T_X)$ である。従って
\begin{equation}\label{eq:gluing-tangent}
 H^0(X,\Lambda^2T_X)\xrightarrow{\sim}
 \Ext^1_{\mathcal K}(E,E).
\end{equation}
これは単なる次元比較ではなく、制限射の核とgluing完全列による同型である。
非可換層変形から $E$ の変形への接写像は、上三角dg模型における
混合Hochschild cochainの制限により、この同型と一致する。

\begin{theorem}[滑らかな対象変形空間内の錐]\label{thm:object-cone}
命題\ref{thm:fibre}の固定残余圏の形式ファイバー $\mathfrak F_X$ から、
gluing object $E$ の変形空間への自然な写像は閉埋め込みである。
形式座標を選べば
\[
 \mathfrak F_X=
 \Spf\CC[[a_{ij}]]/(q_{ijkl})
 \lhook\joinrel\longrightarrow
 \Spf\CC[[a_{ij}]]=\Def_E.
\]
この7次元の部分空間は、原点では10次元全体を接空間に持つ。
\end{theorem}

\begin{proof}
構造層の変形で持ち上げた $\OO_X,\OO_X(1)$ から、相対mutationで $Q_A$ を作り、
右随伴で $E_A=i_A^!Q_A$ を得る。perfectな射影bimoduleを使えば、
これらはArtin基底変換と可換である。
固定残余圏の上で標識を保つ同値 $\mathcal K_A\simeq\mathcal K\otimes A$ を選び、
$E_A$ を移す。異なる選択の差は恒等に特殊化する自己同値である。
$\HH^1(\mathcal K)=0$ と第7節の形式比較により、その同型類の形式群は自明なので、
対象の同値類はこの選択に依存しない。

従って $\mathfrak F_X\to\Def_E$ は定義され、接写像は
\eqref{eq:gluing-tangent}の同型である。
対応する完備局所環の射
$\CC[[z_1,\ldots,z_{10}]]\to\CC[[a_{ij}]]/(q)$
は余接空間上同型なので全射である。
$z_i$ の像を $\CC[[a_{ij}]]$ に持ち上げると、その一次項は可逆であり、
形式座標変換を定める。この座標で核はちょうどPlückerイデアルである。
\end{proof}

\subsection{元の例外対との対応}

元の分解で使う二つの表現対象も書ける。
$\pi=i^*:\mathcal D\to\mathcal K$ とすれば、Serre双対性から
\[
 P_r:=i^!\OO_X(r)=S_{\mathcal K}\pi(\OO_X(r+2))[-3],
 \qquad r=0,1.
\]
従って $F\in\mathcal K$ に対し
$\RHom_{\mathcal D}(F,\OO_X(r))=\RHom_{\mathcal K}(F,P_r)$ である。
評価三角に $i^!$ を適用すると
\[
 E\longrightarrow V\otimes P_0\longrightarrow P_1\longrightarrow E[1]
\]
を得る。さらにrotation $\rho$ に対して
$\rho(P_0)\simeq E$ である\cite[Remark 7.2]{FLZ}。
可換の場合には二回のEuler mutationから
\[
 \pi(\OO_X(2))=\Omega^2_{\PP^4}(2)|_X[2],\qquad
 P_0=S_{\mathcal K}\bigl(\Omega^2_{\PP^4}(2)|_X\bigr)[-1]
\]
とも書ける。単一対象 $E$ の自己Extと、二対象間の射を含む全体の貼り合わせは、
変形問題として区別する必要がある。

特に第9節の固定残余圏族 $A'_h$ の一次Poisson類は非零なので、
\eqref{eq:gluing-tangent}から、固定した $\Ku(X_0)$ 内のgluing objectは
既に $h$ の一次から非自明に動く。
この方向は第4節の線形束 $\langle u,v\rangle$ に対応する分解可能双ベクトルの直線である。
残余圏が一定であることは、この追加対象まで一定であることを意味しない。

\section{gluing bimoduleの形式変形基底}\label{sec:gluing-moduli}

\subsection{対象と自己Ext}

例外ブロックを
\[
 \mathcal P=\langle\OO_X,\OO_X(1)\rangle
 \simeq\Perf(\CC Q_5)
\]
と書く。$Q_5$ は二頂点間に5本の平行な矢を持つKronecker quiverである。
$\mathcal D=\langle\mathcal K,\mathcal P\rangle$ のgluing bimoduleを
\[
 M(F,P)=\RHom_{\mathcal D}(F,P),\qquad
 F\in\mathcal K,\ P\in\mathcal P
\]
とする。これは $\mathcal K^{\mathrm{op}}\otimes\mathcal P$ 上のperfectな
双関手として扱う。以下 $\Def_M$ は、二つのdg圏を固定して、このperfect
bimoduleを可換局所Artin環上で標識付きに変形する群oidである。

$\mathcal P$ の相対Hochschild複体は、頂点の $\CC^2$ と
矢空間の $\operatorname{End}(V)$ からなり、微分は
$(a,b)\mapsto(b-a)\id_V$ である。長さ2の矢道は存在しないので
\[
 \HH^0(\mathcal P)=\CC,\qquad
 \HH^1(\mathcal P)=\mathfrak{pgl}(V),\qquad
 \HH^{\ge2}(\mathcal P)=0.
\]
gluing triangle
\begin{equation}\label{eq:full-gluing-triangle}
 CC^\bullet(\mathcal D)\longrightarrow
 CC^\bullet(\mathcal K)\oplus CC^\bullet(\mathcal P)
 \longrightarrow\RHom(M,M)\longrightarrow CC^\bullet(\mathcal D)[1]
\end{equation}
と、第7節およびHKRの計算から次を得る。

\begin{theorem}[全体の貼り合わせの接空間と障害]\label{thm:bimodule-ext}
\[
 \dim\Ext^r(M,M)=
 \begin{cases}1&r=0,\\34&r=1,\\15&r=2,\\0&\text{その他}.\end{cases}
\]
さらに完全列と同型
\begin{gather}
 0\longrightarrow\mathfrak{pgl}(V)\longrightarrow\Ext^1(M,M)
 \xrightarrow{p}H^0(\Lambda^2T_X)\longrightarrow0,
 \label{eq:bimodule-tangent}\\
 \partial:\Ext^2(M,M)\xrightarrow{\sim}\HH^3(X)
 \simeq H^0(\OO_X(2))
 \label{eq:bimodule-obstruction}
\end{gather}
を持つ。24次元の部分空間は、Kroneckerブロックの矢の座標変更の微分である。
\end{theorem}

\subsection{二次障害は五つのPlücker式である}

方向付きdg模型では、二つのブロックをまたぐHochschild cochainは、
入力にちょうど一つの $M$、出力に $M$ を持つ。
両端に $\mathcal K,\mathcal P$ のbar入力を持つこの複体は、
シフト後にはbimoduleのbar自己準同型複体となる。
混合入力へのbrace挿入はbimodule premorphismの合成なので、
\[
 \RHom_{\mathrm{bar}}(M,M)
 \longrightarrow CC^\bullet(\mathcal D)[1]
\]
はdg Lie射である。そのコホモロジー射は
\eqref{eq:full-gluing-triangle}の連結射である。

$\xi\in\Ext^1(M,M)$、$\pi=p(\xi)$ とすると、二次延長の障害は
Yoneda平方 $\xi^2$ であり、上のdg Lie射によって
\begin{equation}\label{eq:yoneda-schouten}
 \partial(\xi^2)=\frac12[\pi,\pi]\in H^0(\Lambda^3T_X)
\end{equation}
となる。ここでは大域双ベクトルの反対称化した双微分cochainとHKRを用いれば、
Hochschildコホモロジー上でGerstenhaber括弧がSchouten括弧に対応することを確認できる。
cochainそのものの厳密な等式を用いているわけではない。

$\pi=\pi(a)$ を第2節の座標で書くと、\eqref{eq:bracket}から
障害写像の15成分が張る二次式の空間は
\[
 \langle q_{0123},q_{0124},q_{0134},q_{0234},q_{1234}\rangle
\]
に一致する。$F_0,\ldots,F_4$ の線形独立性がここでも必要である。
従って独立な二次式は5本であり、その零点の余次元は3である。
また\eqref{eq:bimodule-obstruction}が同型なので、
24個の矢座標方向を含む交差項も含め、二次障害は $\pi$ だけで決まる。

\subsection{全次数の古典的形式基底}

\begin{theorem}[固定した両ブロック間の全形式gluing変形]\label{thm:full-gluing-base}
古典的な同値類の関手 $\pi_0\Def_M$ は、非標準的に
\[
 \Spf
 \frac{\CC[[b_1,\ldots,b_{24},a_{ij}\ (0\le i<j\le4)]]}
 {(q_{ijkl})}
 \simeq
 \widehat{\operatorname{PGL}(V)}_{\id}
 \times\widehat{C(\Gr(2,5))}_0
\]
で表される。次元は31、接空間次元は34である。
恒等に特殊化する矢の座標変更で割ると、
命題\ref{thm:fibre}の7次元の固定残余圏ファイバーを得る。
\end{theorem}

\begin{proof}
\textbf{Maurer--Cartanによる局所環。}
$\Ext^{-1}(M,M)=0$ で自己Extは上に有界なので、
\cite[Theorem 11.8, Proposition 11.10]{ELOII}により導来変形を
h-projective模型の変形に移せる。
\cite[Proposition 6.1]{ELOI}と\cite[Theorem 7.2]{ELOII}により、
strictly unitalな最小 $A_\infty$ 模型 $H=\Ext^*(M,M)$ を使える。
$H^0=\CC\cdot1,\ H^1=\CC^{34},\ H^2=\CC^{15}$ であり、
可換Artin基底の極大イデアル $\mathfrak m$ 上の方程式は
\[
 \kappa(z)=\sum_{n\ge2}m_n(z,\ldots,z)=0,\qquad
 z\in H^1\otimes\mathfrak m.
\]
次数零で恒等に特殊化する可逆射は $(1+c)1$ に限られる。
strict unitにより二つのMC元 $z,z'$ の間でこの射が閉じる条件は
$(1+c)(z-z')=0$ なので、$z=z'$ である。
従って同値類の関手の環は
\[
 R_M=P/I,\qquad P=\CC[[z_1,\ldots,z_{34}]],
\]
ただし $I$ は $\kappa$ の15成分が生成するイデアルである。
その二次項は前小節のYoneda平方である。

この比較で対象のperfect性も保たれる。
固定した対象対でbimoduleを評価するか、compact generatorのMorita模型を選んで考える。
Artin環 $A$ は有限長の剰余体filtrationを持つため、特殊ファイバーが
有界有限次元の複体なら、その評価の導来変形も基底上で有界coherentとなる。
有限Tor振幅の判定から $A$-perfectである。
さらに $\mathcal K^{\mathrm{op}}\otimes\mathcal P$ のperfectな対角分解を
テンソルすれば、圏上でもperfectとなる。
従ってここでのMC関手は、最初に定義したperfect bimoduleの関手を計算している。

\textbf{既存の非可換族から34方向の族を作る。}
$R_{\mathfrak F}=\CC[[a_{ij}]]/(q)$ 上の固定残余圏族を取る。
各 $R_{\mathfrak F}/\mathfrak m^n$ で残余圏の標識付き自明化を選べる。
隣接する自明化の差は恒等に特殊化する自己同値であり、
$\HH^1(\mathcal K)=0$ と第7節の比較から、その同型類は一意である。
従って各隣接段階で自然同型を選び、逐次的に整合する自明化を取れる。
$\HH^0$ は高次のscalarの選択を残すが、この逐次的な選択の存在を妨げない。
また矢空間は階数5の自由加群なので、その基底も帰納的に整合して選べる。
これにより、固定した $\mathcal K,\mathcal P$ 上の形式bimodule $M_a$ を得る。

その $\mathcal P$-作用を矢の自己同型で変える。
左作用の規約は $(g\cdot M)(F,P)=M(F,g^{-1}P)$ とする。
$GL(V)$ のscalar部分はKronecker代数の内部自己同型なので、
同値類上の有効な作用は $\operatorname{PGL}(V)$ である。こうして
\begin{equation}\label{eq:framed-family}
 \widehat{\operatorname{PGL}(V)}_{\id}\times\Spf R_{\mathfrak F}
 \longrightarrow\pi_0\Def_M
\end{equation}
を構成する。
\eqref{eq:bimodule-tangent}によって接写像は同型である。

\textbf{高次の追加方程式がないこと。}
形式座標を直して $P=\CC[[a_{ij},b_1,\ldots,b_{24}]]$、
$J=(q_{ijkl})$ とすれば、\eqref{eq:framed-family}の環射は
$P/I\to P/J$ となり、$I\subset J$ である。
前小節の計算から $I$ の二次項は5本のPlücker式を張る。
これらは $J/\mathfrak m_PJ$ の基底であり、
\[
 J\cap\mathfrak m_P^3=\mathfrak m_PJ
\]
である。実際、$f=\sum c_iq_i$ の二次項が零なら、
$q_i$ の線形独立性から全ての $c_i(0)$ が零になる。
従って $I$ 内に $f_i=\sum_j A_{ij}q_j$ で $A(0)$ 可逆となる元を選べる。
$A$ は $P$ 上可逆なので、全ての $q_j$ は $I$ に入る。よって $I=J$。

この結果、\eqref{eq:framed-family}は同型である。
構成から形式 $\operatorname{PGL}(V)$ の左移動に同変なので、
この作用で割った同値類は $\Spf R_{\mathfrak F}$ となる。
\end{proof}

\begin{remark}[古典的基底と高次構造]
命題\ref{thm:full-gluing-base}は可換局所Artin基底上の標識付き同値類についての主張である。
$M$ のscalar自己同型や導来モジュライの構造は、この通常の完備局所環には含まれない。
古典的イデアルが二次式だけで生成されることから、
自己準同型 $A_\infty$ 代数全体のformalityまでは結論していない。
\end{remark}

\subsection{最後の例外対象を付加できる正確な条件}

単一対象 $E$ の変形 $\widetilde E$ から、
一つの例外対象 $Q$ を加えた圏
$\mathcal B_{\widetilde E}=\langle\mathcal K,Q\rangle_{\widetilde E}$
を作れる。ここで $G\in\mathcal B$ が
$\mathcal B_{\widetilde E}$ のperfect対象に持ち上がる条件を考える。

$\Ext^1(G,G)=0$ なので、持ち上げが存在すれば同型類は一意である。
さらに特殊ファイバーで
$\RHom(Q,G)=\RHom(Q,\OO_X)=V^*$ は次数零だけなので、
導来Nakayamaから、任意のperfectな持ち上げでもこのHomは
次数零の階数5自由加群である。
$G$ をgluing objectとして $\OO_X$ を付加すると、
二つの例外対象のブロックは再びKronecker型となり、
mutation $\OO_X(1)=R_{\OO_X}(Q)[1]$ により元の
$\langle\OO_X,\OO_X(1)\rangle$ を回復する。
逆に、$\Def_M$ の任意の変形から、この二段階の分解と $G$ の持ち上げが得られる。
矢の基底変更は評価射のmutationで定まる $Q$ の同型類を保存するので、
$\Def_M\to\Def_E$ は形式 $\operatorname{PGL}(V)$ の商を経由する。
矢の基底を忘れれば、以上の二つの構成は同型類上で互いに逆である。

従って命題\ref{thm:full-gluing-base}と命題\ref{thm:object-cone}により、
\[
 \left\{\widetilde E\in\Def_E:
 G\text{ が }\mathcal B_{\widetilde E}\text{ のperfect対象に持ち上がる}\right\}
 =
 \widehat{C(\Gr(2,5))}_0\ \subset\ \widehat{\Aff}^{10}.
\]
すなわち、$E$ 単独の変形は障害を持たないが、
最後の例外対象を付加する条件が5本のPlücker式を課す。
これが固定残余圏の7次元特異空間のgluingによる解釈である。

\section{点加群スキームと長さ8の非被約点}\label{sec:points}

\subsection{点加群と中心変数による分離}

体 $k$ 上の次数1生成代数 $A$ に対し、可換 $k$-代数 $T$ 上のpoint moduleとは、
cyclicな次数付き右 $A\otimes_kT$-加群 $N=\bigoplus_{n\ge0}N_n$ で、
$N_0=T$、全ての $N_n$ が階数1の局所自由 $T$-加群となるものをいう。
その同型類の関手を表すスキームを $\mathscr P(A)$ と書く。
定義とtruncationによる構成については\cite[Section 4]{RZ}を参照。
以下では実際に有限段階で安定することを示すので、
strongly noetherian性からの一般的な表現可能性定理には依存しない。

$S=k\langle x,y,z\rangle/(\text{quadratic relations})$ の
二次関係の多重線形化が、$E\subset\PP^2$ 上の自己同型 $\sigma$ のgraphであるとする。
また $\sigma$ は幾何学的固定点を持たず、三次の中心元 $g$ は
このgraphに対応する普遍point moduleを零化すると仮定する。
$u,v$ を中心変数とし、
\[
 B=S[u,v],\qquad A=B/(g+u^3+v^3)
\]
とおく。$S_{\mathrm{ab}},A_{\mathrm{ab}}$ は可換化を表す。

\begin{lemma}[中心変数を付加したpoint scheme]\label{lem:point-splitting}
上の仮定の下で、スキームとして
\[
 \mathscr P(A)=E\ \sqcup\ \operatorname{Proj}A_{\mathrm{ab}}
\]
である。point moduleのshiftは、第一成分で $\sigma$、第二成分で恒等となる。
\end{lemma}

\begin{proof}
連続する二つの座標点を
$p=(x,y,z,u,v)$、$p'=(x',y',z',u',v')$ と書く。
中心性から $x_iu'-ux_i'=0$、$x_iv'-vx_i'=0$ などを得る。
$u\ne0$ のチャートでは $u'$ も局所的に可逆である。
実際、ある幾何学的点で $u'=0$ なら中心性から $p'$ の全座標が零となる。
従って両側を規格化すると $p=p'$ であり、残る二次式はちょうど
$S_{\mathrm{ab}}$ の関係式である。$v\ne0$ でも同じである。

最初の幾何学的点で $u=v=0$ なら、少なくとも一つの $x_i$ が非零なので、
中心性から $u'=v'=0$ も従う。
このような幾何学的点では、
$(x':y':z')=\sigma(x:y:z)$ である。
固定点がないため、ある $2\times2$ 小行列式
$\Delta=x_ix_j'-x_jx_i'$ は非零であり、局所近傍で可逆となる。
二本の中心性の式から $u=u'=0$ がスキームとして従い、
同様に $v=v'=0$ である。従ってこの成分の法方向に冪零構造も残らない。
これらの近傍は全ての幾何学的点を覆うので、二次関係スキームは
\[
 \operatorname{Graph}(\sigma)\ \sqcup\
 \Delta_{\operatorname{Proj}(S_{\mathrm{ab}}[u,v])}
\]
である。両者の台は交わらない。

$X_n(B)\subset(\PP^4)^n$ を $N_0,\ldots,N_n$ のtruncationスキームとする。
上の二次関係がgraphなので、全ての $n\ge2$ で
$X_n(B)$ は最初の座標だけから決まる。
三次式は $X_3(B)$ 上で、$E$ 成分では仮定から零、diagonal成分では
可換化した $g+u^3+v^3$ となる。その零点はshiftで保たれる。
従って $X_n(A)\simeq X_3(A)$ が全ての $n\ge3$ で成立し、
このスキームが無限point moduleの関手を表す。
議論は任意の冪零を含む試験環上で成り立つ。
\end{proof}

\subsection{元のFermat族の全ての非零ファイバー}

\begin{theorem}[明示的族のpoint scheme]\label{thm:point-scheme}
$\lambda\in\CC^\times$ とし、
\[
 J_0=\CC[x,y,z]/(x^2,y^2,z^2)
\]
とおく。このとき
\[
 \boxed{
 \mathscr P(A_\lambda)
 =E_{\lambda^2}\ \sqcup\
 \coprod_{\zeta^3=-1}\operatorname{Spec}J_0
 },\qquad
 E_s=V(x^3+y^3+z^3+sxyz).
\]
各非被約点の長さは8、Hilbert関数は $(1,3,3,1)$ である。
この式は $\lambda^6=-27$ の三角形のファイバーにも成立する。
$\lambda=0$ では $\mathscr P(A_0)=X_0$ である。
\end{theorem}

\begin{proof}
第5節の多重線形化行列 $M_\lambda$ の階数は、
$E_{\lambda^2}$ 上で常に2である。
実際、階数が1以下なら、二つの小行列式
$xy-\lambda z^2$ と $xy+\lambda z^2$ がともに零となり $z=0$。
巡回した式から $x=y=0$ となるので矛盾する。
逆向きの関係行列も同じ性質を持つため、関係スキームは自己同型のgraphである。
diagonalに代入した二次式は
$\lambda x^2=\lambda y^2=\lambda z^2=0$ なので、固定点はない。

$G_\lambda$ が普遍point moduleを零化することは、滑らかな一般ファイバーで
$S_\lambda/(G_\lambda)=B(E_{\lambda^2},L,\sigma)$ から従う。
非零パラメータ上の相対Hesse曲線は、一般ファイバーが幾何学的整で
基底上flatな整スキームである。
$G$ の多重線形化はその上の可逆層の切断なので、一般ファイバーでの零化から
全体で零となる。これにより三角形の6ファイバーも含めて
補題\ref{lem:point-splitting}を適用できる。

可換化は
\[
 (A_\lambda)_{\mathrm{ab}}=
 \frac{\CC[x,y,z,u,v]}
 {(x^2,y^2,z^2,\ u^3+v^3-2\lambda^2xyz)}
\]
である。そのProjの台は $[0:0:0:1:\zeta]$、$\zeta^3=-1$ の3点である。
$u=1$ とし、$T=xyz\in J_0$ とおくと $T^2=0$ である。
$\zeta$ の枝で
\[
 v=\zeta+\frac{2\lambda^2}{3\zeta^2}T
\]
と置けば、三次式は厳密に消える。
従って各局所成分は $\operatorname{Spec}J_0$ である。
$J_0$ の基底は
\[
 1,x,y,z,xy,xz,yz,xyz
\]
であり、長さとHilbert関数も従う。
\end{proof}

第一座標の $\PP^4$ への埋め込みで、その斉次イデアルは
\[
 (u,v,D+\lambda^2xyz)\ \cap\
 (x^2,y^2,z^2,u^3+v^3-2\lambda^2xyz)
\]
である。Hilbert多項式は $3n+24$ となる。
従って、楕円曲線と三つの被約点だけを記述すると、
零次元部分の長さ24の情報を失う。

\subsection{残余圏を固定する族でも非被約構造は動く}

第9節の族 $A'_h$ を $K=\CC((h))$ 上で考え、
\[
 r=(1-h^3)^{1/3},\qquad
 a=\frac{1-r}{h}=\frac{h^2}{3}+\frac{h^5}{9}+O(h^8)
\]
とおく。この $r$ は第9節の $t$ と同じ元である。
可換化した二次関係は
\[
 x^2=a yz,\qquad y^2=a zx,\qquad z^2=a xy
\]
である。また
\[
 a^3=\frac{(1-r)^2}{1+r+r^2},\qquad
 a^3-1=-\frac{3r}{1+r+r^2}\ne0.
\]
これらの二次式は $\PP^2_K$ に共通零点を持たないので、
対応するSklyaninの自己同型は幾何学的固定点を持たない。

\begin{theorem}[固定残余圏族の非被約点]\label{thm:fixed-points}
\[
 J_a=
 \frac{K[x,y,z]}
 {(x^2-a yz,\ y^2-a zx,\ z^2-a xy)}
\]
とおくと、
\[
 \mathscr P(A'_h\otimes K)
 =(E_0)_K\ \sqcup\
 \coprod_{\zeta^3=-1}\operatorname{Spec}J_a.
\]
$J_a$ は長さ8でHilbert関数 $(1,3,3,1)$ を持つ。
その局所代数に内在する三次形式は
\[
 \Phi_a(X,Y,Z)=a(X^3+Y^3+Z^3)+6XYZ.
\]
これは滑らかな平面三次曲線を定め、その $j$-不変量は
\begin{equation}\label{eq:fat-j}
 j_{\mathrm{fat}}(a)
 =\frac{110592(a^3-1)^3}{a^3(a^3+8)^3}
 =-5832h^{-6}\bigl(1+O(h^3)\bigr)
\end{equation}
である。従って残余圏と被約な楕円曲線成分が一定でも、
非被約点の局所代数は一定ではない。
\end{theorem}

\begin{proof}
二次関係から $(1-a^3)x^2y=0$ とその置換を得る。
従って指数型 $(2,1,0)$ の三次単項式は零であり、
\[
 x^3=y^3=z^3=a\,xyz,\qquad
 x\,xyz=y\,xyz=z\,xyz=0.
\]
次数4以上は零となり、$1,x,y,z,xy,xz,yz,xyz$ が張る。
三つの二次式の射影共通零点がないので、多項式環のCohen--Macaulay性から
それらは正則列であり、Hilbert級数は $(1+t)^3$ である。
従ってこの8元は独立でもあり、基底となる。
添付コードはこの積表の512個の基底三つ組について結合律を厳密に確認する。

$T=xyz$ とおくと、中心元の可換化は
\[
 (G'_h)_{\mathrm{ab}}
 =\bigl(2(1+h^3)a+(1-2h^3)a-3h^2\bigr)T
 =3(a-h^2)T.
\]
$u=1$、$\zeta^3=-1$ の枝で
\[
 v=\zeta-\frac{a-h^2}{\zeta^2}T
\]
と置けば三次式が消える。補題\ref{lem:point-splitting}から最初の分解を得る。

極大イデアルを $\mathfrak n$ と書くと、
$\mathfrak n/\mathfrak n^2$ は $x,y,z$ の3次元空間、
$\mathfrak n^3=K\cdot T$ は1次元で、これらと三重積は局所代数に内在する。
直接展開により
\[
 (Xx+Yy+Zz)^3
 =\bigl(a(X^3+Y^3+Z^3)+6XYZ\bigr)T.
\]
従って局所代数の同型は、$\Phi_a$ を線形座標変換と非零scalar倍の範囲で保存する。
$a\ne0$、$a^3\ne-8$ なのでその零点は滑らかなHesse曲線である。
第9節で検証した $j(E_s)$ に $s=6/a$ を代入すると
\eqref{eq:fat-j}を得る。
これは局所代数の同型不変量であり、$h$ に関して非定数である。
\end{proof}

この三次形式はMacaulay dual cubicとも解釈できるが、
ここでは極大イデアルの三重積から直接定義した。
$a=0$ では $\Phi_0=6XYZ$ は三角形であり、一般の $a$ における滑らかな三次曲線と異なる。
固定残余圏族の第一座標のpoint schemeの斉次イデアルは
\[
 (u,v,D)\ \cap\
 \bigl(x^2-a yz,y^2-a zx,z^2-a xy,\,
 u^3+v^3+3(a-h^2)xyz\bigr)
\]
である。

\subsection{平坦な代数族と点加群の特殊化}

$a\in R=\CC[[h]]$ として同じ $J_a$ を定義すると、
$1-a^3$ は単元なので、これは $R$ 上階数8の自由代数である。
上の $v$ の式と固定楕円曲線のgraph族は、相対truncationスキーム内で
一般ファイバーの3個の非被約点と曲線のflatな水平閉包を与える。
その特殊ファイバーは
\[
 E_0\ \sqcup\ 3\operatorname{Spec}J_0
\]
であり、Fermat三次元多様体 $X_0$ の真の閉部分スキームである。

一方、$N_0,\ldots,N_3$ を扱う相対length-4 truncationスキームは
$(\PP^4_R)^3$ の閉部分スキームとして射影的であり、
一般ファイバーは上の1次元スキーム、特殊ファイバーは対角上の $X_0$ である。
Hilbert多項式の次数が異なるため、この相対truncationスキームはflatでない。
これは $A'_h$ 自体の次数ごとのflat性と両立する。
基底全体で無限point module関手の一様な有限表現可能性を主張する必要はない。

point schemeは次数付き代数、従って偏極を持つtailsの三つ組
\[
 (\qgr A,\pi A,(1))
\]
の不変量である。実際、ここではAS Gorensteinのdepthから
$H^0_{\mathfrak m}(A)=H^1_{\mathfrak m}(A)=0$ なので、
局所コホモロジー完全列により
$A\simeq\bigoplus_{n\ge0}\operatorname{Hom}_{\qgr A}(\pi A,\pi A(n))$
と再構成できる。任意の偏極を忘れたdg Morita同値に対する不変量としては用いていない。

\section{貼り合わせの導来変形と形式弧上のtorsion}\label{sec:derived-gluing}

本節では命題\ref{thm:full-gluing-base}の古典的基底に残る導来構造を決定する。
固定する圏は $\mathcal K=\Ku(X)$ と
$\mathcal P=\langle\OO_X,\OO_X(1)\rangle$ であり、変形対象はその間の
perfect gluing bimodule $M$ である。
導来基底には非正次数の可換dg Artin $\CC$-代数を用いる。
この問題は、自己準同型の結合的な $A_\infty$ 構造全体を決定する問題とは区別する。

\subsection{scalarを除いたChevalley--Eilenberg模型}

\begin{theorem}[導来形式基底の非標準的な正規形]\label{thm:derived-gluing-normal}
$M$ のscalar自己同型を除いた可換導来形式変形は、次の完備可換dg代数で表される。
\begin{gather}
 \mathscr R_M\simeq
 \CC[[b_1,\ldots,b_{24},a_{ij}]]
 \otimes\Lambda(\theta_0,\ldots,\theta_4,\rho_1,\ldots,\rho_{10}),
 \label{eq:derived-gluing-normal}\\
 |a_{ij}|=|b_i|=0,\qquad |\theta_i|=|\rho_j|=-1,\qquad
 d\theta_i=q_i(a),\quad d\rho_j=0,\quad da_{ij}=db_i=0.
 \nonumber
\end{gather}
ここで $q_i$ は5本のPlücker式であり、符号と順序は任意に固定する。
scalarを戻した導来形式変形は、この形式空間と
$B\widehat{\mathbb G}_{m,1}$ との積である。

従って $\RHom(M,M)$ のcommutator dg Lie代数は
$L_\infty$ 同値の意味でformalである。
この結論は結合的な $A_\infty$ 代数のformalityを主張しない。
\end{theorem}

\begin{proof}
\textbf{次数とstrict unit。}
命題\ref{thm:bimodule-ext}と
\cite[Theorem 11.8, Proposition 11.10, Theorem 7.2]{ELOII}により、
strictly unitalな最小 $A_\infty$ 模型
\[
 H^0=\CC\cdot1,\qquad H^1=\CC^{34},\qquad H^2=\CC^{15}
\]
を取る。その反対称化を $L$ とし、括弧 $\ell_n$ の次数を $2-n$ とする。
strict unitはcommutatorについて中心的であり、
$n\ne2$ の括弧にunitを入力すると零になる。
残りの入力が全て次数1なら出力は次数2であり、
一つでも次数2の入力を持てば出力は次数3以上となって零になる。
従ってscalar以外の全ての演算は
\[
 \ell_n:\Sym^nH^1\longrightarrow H^2\qquad(n\ge2)
\]
に限られる。次数1の入力についてgraded反対称性は通常の対称性になる。
規約を揃えてMaurer--Cartan方程式を
\[
 \kappa(z)=\sum_{n\ge2}\frac1{n!}\ell_n(z,\ldots,z)=0
\]
と書く。この $\kappa$ は前節の結合的なMC方程式と同じである。

よって $L_+=H^1\oplus H^2$ の完備Chevalley--Eilenberg代数は
\begin{equation}\label{eq:derived-gluing-ce}
 \CC[[z_1,\ldots,z_{34}]]\otimes\Lambda(\eta_1,\ldots,\eta_{15}),
 \qquad |\eta_j|=-1,\quad d\eta_j=\kappa_j(z),\quad dz_i=0.
\end{equation}
実際、CE代数の生成元は $(L_+[1])^\vee$ であり、
次数1の自己Extは次数零の座標、次数2の自己Extは次数 $-1$ の座標を与える。
この代数から可換dg Artin環への射はMC解そのものであり、
simplex上の微分形式を用いれば同じ比較が写像空間とhomotopyにも成り立つ
\cite[Proposition 1.1, Section 4]{Getzler}。
上の次数の議論により、$\eta$ を含む未記載の微分項は存在しない。
中心的な $H^0=\CC\cdot1$ は $L_\infty$ 直和因子であり、
その形式moduliが $B\widehat{\mathbb G}_{m,1}$ である。

\textbf{15個の方程式の行変換。}
命題\ref{thm:full-gluing-base}の証明で得た形式座標を用いて
\[
 P=\CC[[a_{ij},b_1,\ldots,b_{24}]],\qquad
 (\kappa_1,\ldots,\kappa_{15})=(q_0,\ldots,q_4)\subset P
\]
とする。二次項の階数は5なので、$\eta$ の定数係数での基底変換により、
最初の5成分の二次項が独立であるとしてよい。
すると
\[
 \begin{pmatrix}\kappa_1\\ \vdots\\ \kappa_5\end{pmatrix}=Aq,
 \qquad
 \begin{pmatrix}\kappa_6\\ \vdots\\ \kappa_{15}\end{pmatrix}=Bq,
 \qquad A\in\operatorname{GL}_5(P),\quad B\in\operatorname{Mat}_{10\times5}(P).
\]
$A(0)$ が可逆なのは二次項の独立性による。
そこで15個の負次数生成元を行列
\[
 \begin{pmatrix}A^{-1}&0\\-BA^{-1}&I_{10}\end{pmatrix}
\]
で変換すると、その微分は $(q_0,\ldots,q_4,0,\ldots,0)^t$ となる。
行列の係数は $P$ に属し $dP=0$ なので、係数を微分する追加項はない。
これが\eqref{eq:derived-gluing-normal}である。

新しいCE微分は純粋に二次式である。
双対の完備coalgebraに戻せば、これは高次括弧 $\ell_{n\ge3}$ が零である
最小 $L_\infty$ 模型を与える。
残る二項括弧は、基底変更を除いて元のcohomology bracketであるから、
commutator dg Lie代数のformalityを得る。
\end{proof}

\begin{remark}[この正規形が保存する情報]
次数1の元 $x,y$ に対してcommutatorが記録するのは
$m_2(x,y)+m_2(y,x)$ であり、結合積の反対称部分は記録しない。
高次演算でも同様に、対称化されない $A_\infty$ データは残り得る。
従って上の証明から結合的formalityは従わない。
また形式座標と負次数生成元の変換は非標準的であり、
導来正規形の $\operatorname{PGL}(V)$-同変性も主張しない。
命題\ref{thm:full-gluing-base}で証明した古典的な同変性は引き続き用いられる。
\end{remark}

\subsection{形式弧を実際のperfect族として扱う}

$R=\CC[[h]]$、$F_\eta=\CC((h))$ とし、命題\ref{thm:full-gluing-base}の基底に
\begin{equation}\label{eq:gluing-formal-arc}
 b(h)\in hR^{24},\qquad a(h)\in hR^{10},\qquad q(a(h))=0
\end{equation}
という形式弧を取る。
$\Spf$ の原点完備化には、この弧の「非零の閉点」が存在するわけではない。
以下のgeneric fibreは、実際の $R$-対象を作ってから $h$ を可逆化したものである。

\begin{lemma}[形式弧のeffectivity]\label{lem:gluing-arc-effective}
\eqref{eq:gluing-formal-arc}が定める整合的なArtin変形は、
固定した双圏 $(\mathcal K^{\mathrm{op}}\otimes\mathcal P)\otimes R$ 上の
perfect bimodule $M_R$ として実現できる。
その自己準同型複体は $R$-perfectであり、$R/(h^n)$ と $F_\eta$ への
導来基底変換と可換である。
\end{lemma}

\begin{proof}
周囲の圏 $\Perf(\mathcal K^{\mathrm{op}}\otimes\mathcal P)$ はsmooth properである。
Toën--Vaquié\cite[Theorem 3.6, Corollaries 3.15, 3.21--3.22]{TV}により、
負の自己Extが零である $M$ の近傍の対象モジュライはArtin 1-stackである。
弧から得られる $M_n$ はこの同じ開部分に属する。
affineなsmooth atlas $U$ と $M$ 上の点 $u_0$ を選ぶ。
smooth性により、$M_{n+1}$ の上で $u_n$ を $u_{n+1}$ へ逐次持ち上げられる。
剰余体は代数閉体 $\CC$ なので、各持ち上げは局所Artin環自身の上で選べる。
$R=\varprojlim R/(h^n)$ と $U$ のaffine性により、
整合する $u_n$ は $\operatorname{Spec}R\to U$ を定める。
普遍対象の引戻しが求める $M_R$ である。

このmoduliでperfect性は定義に含まれる。
Morita模型で基底上のperfect性を先に確認しても、smoothな対角を用いれば
圏上のperfect性を得る\cite[Lemma 2.8]{TV}。
proper性とperfect性から自己準同型複体の有限性と基底変換が従う。
この議論は完備環 $R$ 上の実現であり、弧が有限型の複素代数曲線に由来するとは主張しない。
\end{proof}

\subsection{Jacobian複体とSmith標準形}

\begin{theorem}[形式弧上の自己Extとtorsion]\label{thm:gluing-arc-ext}
\eqref{eq:gluing-formal-arc}で $a(h)\ne0$ とし、
\[
 \nu=\min_{i<j}\operatorname{ord}_h a_{ij}(h)\ge1,\qquad
 T_\nu=(R/(h^{\nu}))^{\oplus3}
\]
とおく。固定した双圏上の自己Extは
\begin{equation}\label{eq:gluing-arc-ext}
 \Ext^0(M_R,M_R)=R,\qquad
 \Ext^1(M_R,M_R)=R^{31},\qquad
 \Ext^2(M_R,M_R)=R^{12}\oplus T_\nu
\end{equation}
であり、その他の次数では零である。
従って $M_\eta=M_R\otimes_R F_\eta$ について
\[
 \bigl(\dim_{F_\eta}\Ext^0(M_\eta,M_\eta),
       \dim_{F_\eta}\Ext^1(M_\eta,M_\eta),
       \dim_{F_\eta}\Ext^2(M_\eta,M_\eta)\bigr)=(1,31,12)
\]
となる。この計算は全ての非零の分解可能形式弧に成り立つ。
\end{theorem}

\begin{proof}
\textbf{MC方程式の微分。}
最小模型を弧のMC元 $z(h)$ でtwistする。
自己準同型複体の次数1から次数2への微分は
\[
 x\longmapsto\sum_{n\ge1}\frac1{n!}\ell_{n+1}(z,\ldots,z,x)
       =D\kappa_z(x)
\]
である\cite[Proposition 4.4]{Getzler}。
unitの微分は零であり、次数2から出る微分も零である。
従って $\RHom(M_R,M_R)$ は、次数 $0,1,2$ に置かれた
\begin{equation}\label{eq:gluing-jacobian-complex}
 [\,R\xrightarrow{0}R^{34}\xrightarrow{D\kappa_z}R^{15}\,]
\end{equation}
と同値である。
この比較は各 $R/(h^n)$ 上で最小模型の準同型をtwistして作れ、
その式は $h$-進に整合する。
左辺のperfect性と右辺の有限自由性により、逆極限でも同じ比較が成り立つ。

形式座標と方程式の行変換は、この複体上でも可逆変換になる。
実際 $\kappa(z(y))=C(y)q(y)$ なら、MC軌跡 $q=0$ の上で
\[
 D\kappa_z\,Dz_y=C(y)\,Dq_y
\]
であり、$DC$ に由来する項は $q=0$ により消える。
可逆な定数項を持つ全ての変換行列は $R$ 上でも可逆である。
よって\eqref{eq:gluing-jacobian-complex}の非零となり得る微分は、
適当な基底で
\begin{equation}\label{eq:gluing-wedge-differential}
 R^{24}\oplus\Lambda^2R^5
 \longrightarrow\Lambda^4R^5\oplus R^{10},\qquad
 (b,v)\longmapsto(a(h)\wedge v,0)
\end{equation}
となる。ここでは $q(a)=\tfrac12a\wedge a$ という規格化を用いた。

\textbf{DVR上の階数2交代行列。}
$a=h^{\nu} a'$ と書くと、$a'$ はprimitiveであり、$R$ は整域なので $q(a')=0$ である。
基底を置換して $a'_{01}$ がunitとしてよい。
添字 $0,1$ を含む三つのPfaffian式により、残る
$a'_{23},a'_{24},a'_{34}$ を $a'_{01},a'_{0i},a'_{1i}$ で表せる。
従って $a'$ は $R$ 上の独立なprimitiveな二ベクトルのwedgeであり、
$\operatorname{GL}_5(R)$ による変換で $e_0\wedge e_1$ になる。

$R^5=U\oplus W$、$U=\langle e_0,e_1\rangle$、$\rk W=3$ とする。
$(e_0\wedge e_1)\wedge-$ は
$\Lambda^2U\oplus(U\otimes W)$ 上で零であり、
$\Lambda^2W$ を $\Lambda^2U\otimes\Lambda^2W$ に同型に写す。
後者は階数3の直和因子なので、
\eqref{eq:gluing-wedge-differential}の非零のSmith因子は
\[
                         h^{\nu},\quad h^{\nu},\quad h^{\nu}
\]
に尽きる。kernelの階数は $24+7=31$、cokernelは $R^{12}\oplus T_\nu$ であり、
主張を得る。
\end{proof}

\begin{remark}[特殊ファイバーで戻る三つの接方向]
中心の $(1,34,15)$ と弧上の計算は、導来基底変換によって整合する。
実際、完全列
\[
 0\longrightarrow\Ext^1(M_R,M_R)\otimes_R\CC
 \longrightarrow\Ext^1(M,M)
 \longrightarrow\operatorname{Tor}_1^R(\Ext^2(M_R,M_R),\CC)
 \longrightarrow0
\]
の次元は $31,34,3$ である。
また $\Ext^2(M_R,M_R)\otimes_R\CC$ の次元は $12+3=15$ である。
従って障害加群のtorsionが、原点で増える三つの接方向を正確に記録している。
$a(h)$ が恒等的に零なら微分も零であり、自己Extの階数は $(1,34,15)$ のままである。
\end{remark}

\subsection{全体圏のHochschild cohomology}

\begin{theorem}[形式弧上のHochschild加群]\label{thm:gluing-arc-hh}
命題\ref{thm:gluing-arc-ext}の仮定のもとで
$\mathcal D_R=\langle\mathcal K_R,\mathcal P_R\rangle_{M_R}$ とする。
以下のHochschild cohomologyは $R$ 線形のものである。
\begin{equation}\label{eq:gluing-arc-hh}
 \HH^i(\mathcal D_R)=
 \begin{cases}
 R&i=0,\\
 R^{17}&i=2,\\
 R^{12}\oplus(R/(h^{\nu}))^{\oplus3}&i=3,\\
 0&\text{その他}.
 \end{cases}
\end{equation}
従ってgeneric fibreのHochschild次元は、次数 $0,2,3$ でそれぞれ $1,17,12$ である。
\end{theorem}

\begin{proof}
perfect bimoduleによるtriangular gluingはsmooth proper性を保つ。
\eqref{eq:full-gluing-triangle}は相対bar模型でも成り立ち
\cite[Theorem 7.7, Remark 7.8]{Kuz}、全ての複体は $R$-perfectである。
固定したブロックについて
\[
 \HH^*(\mathcal K_R)=R\text{ in degree }0\;\oplus\;R^{10}\text{ in degree }2,
 \qquad
 \HH^*(\mathcal P_R)=R\text{ in degree }0\;\oplus\;R^{24}\text{ in degree }1.
\]
次数零の射 $R^2\to\operatorname{End}(M_R)=R$ はscalarの差を取る全射である。
また
\[
 \HH^1(\mathcal P_R)=R^{24}\longrightarrow\Ext^1(M_R,M_R)=R^{31}
\]
は直和因子への単射であり、そのcokernelは $R^7$ である。
これは命題\ref{thm:full-gluing-base}の古典的な同変性から分かる。
左移動の微分は $b$ 因子上の可逆行列であり、$a$ 成分を持たない。
ここで導来正規形の同変性を仮定する必要はない。

残る問題はcharacteristic map
\begin{equation}\label{eq:gluing-characteristic-arc}
 \chi_R:\HH^2(\mathcal K_R)=R^{10}\longrightarrow\Ext^2(M_R,M_R)
\end{equation}
である。この射は次元の比較だけでは決まらない。
命題\ref{thm:fibre}の証明で、全NC族の形式基底を
\[
 \CC[[s_1,\ldots,s_{10},a_{ij}]]/(q)
\]
とし、残余圏への写像を $s$ への射影に直した。
$a=a(h)$ と矢のframingを固定し、
\[
 s_j=\epsilon c_j(h),\qquad\epsilon^2=0,\qquad c_j(h)\in R
\]
と引き戻す。この族は $\HH^2(\mathcal K_R)$ の任意の基底方向を
全体圏の一次変形へ持ち上げる。従って、その障害である $\chi_R$ は零である。

入力のNC族を形式族のまま扱う場合にも、この議論は厳密に行える。
各 $R/(h^n)[\epsilon]$ 上で整合するKodaira--Spencer類を取り、逆極限に移ればよい。
ここでは $\mathcal D_R\otimes_RR/(h^n)$ を、
$(s=0,a(h),b(h))$ に沿う既存のframed NC族と整合的に同一視し、
その同一視によって類を運ぶ。
$CC^\bullet(\mathcal D_R)$ は完備環上のperfect複体であり、
有限段階のcohomologyは有限長なのでMittag--Leffler条件を満たす。
Milnor完全列の $\varprojlim{}^1$ は零になり、
\[
 \HH^i(\mathcal D_R)\simeq
 \varprojlim_n\HH^i(\mathcal D_R\otimes_RR/(h^n))
\]
を得る。従って変動する圏について別の代数化定理を仮定せずに、
必要な $R$ 上の類と $\chi_R=0$ を得られる。

gluingの長完全列は
\[
 0\longrightarrow R^7\longrightarrow\HH^2(\mathcal D_R)
 \longrightarrow R^{10}\longrightarrow0,
 \qquad
 \HH^3(\mathcal D_R)\simeq\Ext^2(M_R,M_R)
\]
を与える。最初の列は $R$-加群として分裂する。
次数零の差写像と矢座標写像の単射性から $\HH^0=R,\HH^1=0$ も従い、
全ての主張を得る。
\end{proof}

特殊ファイバーに戻るときには
\[
 0\longrightarrow\HH^2(\mathcal D_R)\otimes_R\CC
 \longrightarrow\HH^2(X)
 \longrightarrow\operatorname{Tor}_1^R(\HH^3(\mathcal D_R),\CC)
 \longrightarrow0
\]
があり、$17+3=20$ を再現する。同様に次数3は $12+3=15$ となる。
このように自己ExtとHochschild cohomologyの両方で、
同じ三つの $h^{\nu}$-torsionが原点での跳びを記録する。

\subsection{滑らかな古典的moduliに残る12個の導来方向}

$F_\eta$ 上では $a(h)$ は非零の分解可能双ベクトルである。
$a_{01}\ne0$ のchartでは、三つのPfaffian式が正則列となり、
残る二つは関係式 $Aq=0$ によってその線形結合で表される。
三つの方程式を局所座標とし、対応するKoszul生成元との組を除くと、
$M_\eta$ におけるscalarを除いた導来局所模型は
\begin{equation}\label{eq:generic-gluing-derived}
 F_\eta[[u_1,\ldots,u_{31}]]\otimes\Lambda(v_1,\ldots,v_{12}),
 \qquad |v_i|=-1,\quad d=0
\end{equation}
となる。従って古典的moduliは31次元で滑らかであり、
commutator $L_\infty$ 代数はこの点でhomotopy abelianである。
12個の方向のうち10個は\eqref{eq:derived-gluing-normal}の $\rho$、
残り2個は5本のPfaffian式が余次元3を定めることによる。

これは、非零の $\Ext^2$ が必ず古典的な障害方程式を生むわけではないことを示す。
scalarを除いたvirtual dimensionは $34-15=19$ のままであり、
generic fibreでは古典的次元31との差が12個の導来方向として残る。
本節では全体圏のHochschild dg Lie代数のformalityや、
導来moduliの $\operatorname{PGL}(V)$ 商の積表示までは結論していない。


\section{安定対象モジュライ上の適合条件とBrill--Noether接錐}\label{sec:bn-algebraic}

\subsection{既知の10次元モジュライの中のgluing object}

第\ref{sec:gluing-objects}節の記号
$\mathcal K=\Ku(X)$、$Q=\Omega^1_{\PP^4}(1)|_X$、
$E=i^!Q$、$G=j^!\OO_X$ を引き続き用いる。
$\ell\subset X$ を直線とし、$v=[I_\ell]$ とおくと
\[
 \operatorname{ch}(v)=\left(1,0,-\frac{H^2}{3},0\right),\qquad
 \operatorname{ch}(E)=3\operatorname{ch}(v)=(3,0,-H^2,0).
\]
Serre-invariantな安定性条件 $\sigma$ を選び、対象の位相を固定する。
Feyzbakhsh--Liu--Zhang\cite[Section 3.2, Proposition 5.5]{FLZ}によれば、
$E$ は $\sigma$-stableであり、class $3v$ の他のstable対象は、
shiftを除いてGieseker-stableな層である。

\begin{theorem}[安定対象モジュライでの位置づけ]\label{thm:stable-ambient}
半安定対象の粗モジュライ空間
$M^{\mathrm{ss}}_\sigma(\mathcal K,3v)$ は既約かつ射影的であり、
そのstable部分
\[
 M^{\mathrm s}=M^{\mathrm s}_\sigma(\mathcal K,3v)
\]
は滑らかな10次元空間である。
$e=[E]$ はこのstable部分の点で、その完備局所環は
$\CC[[z_1,\ldots,z_{10}]]$ である。
階数3のnormalized Ulrich束は $M^{\mathrm{ss}}_\sigma(\mathcal K,3v)$ の
非空な稠密開部分をなし、$M^{\mathrm s}$ との交わりも非空かつ稠密である。
$E$ 自身はこのUlrich束の部分には属さない。
\end{theorem}

\begin{proof}
既約性、射影性とUlrich束の開埋め込みは
\cite[Section 6, Corollary 6.5]{FP}による。
stable部分の滑らかさは\cite[Theorem 1.2]{PY}であり、
その次元は
\[
 1-\chi(3v,3v)=1-9\chi(v,v)=10
\]
である。$e$ での完備局所環は命題\ref{thm:gluing-ext}からも直接得られる。
非空なUlrichの開部分と非空なstable部分は、既約な
$M^{\mathrm{ss}}_\sigma(\mathcal K,3v)$ 内で交わる。
一方、\eqref{eq:E-triangle}から
$\mathcal H^{-1}(E)=\OO_X(-1)$、$\mathcal H^0(E)=Q$ なので、
$E$ は層ではなく二項複体である。
\end{proof}

ここでnormalized Ulrich束とは、次数付き切断加群の生成元を次数1に置く
\cite{FP}の規約であり、Chern characterは $(3,0,-H^2,0)$ である。
生成元を次数0に置く規約とは $\OO_X(1)$ によるtwistだけ異なる。
また、射影的と述べたのは半安定対象の粗モジュライであり、
stable部分そのもののpropernessを主張してはいない。

\subsection{相対対象モジュライによる適合条件の代数化}

第\ref{sec:gluing-moduli}節では、$E$ の変形 $\widetilde E$ に対し、
$G$ が $\mathcal B_{\widetilde E}$ のperfect対象へ持ち上がる条件を
閉形式部分空間 $\mathfrak F_X\subset\Def_E$ として決定した。
この条件を実際の代数的な局所部分空間として実現する。

\begin{theorem}[代数的な適合部分空間]\label{thm:algebraic-compatibility}
$e\in M^{\mathrm s}$ の適当な\'etale近傍 $(U,u)$ と、
有限型の閉部分スキーム $(C,c)\subset(U,u)$ を選べて、次を満たす。
$U$ は滑らかな10次元スキームで、$\mathcal K$ 内のuniversal perfect対象
$\mathcal E_U$ を持ち、$\mathcal E_u\simeq E$ である。
$C$ は、相対圏 $\mathcal B_{\mathcal E_U}$ のsimple perfect対象の
モジュライのうち、$g=[G]$ を通る一つの\'etale局所的な枝である。
completed closed immersionは
\[
 \widehat C_c\longrightarrow\widehat U_u
 \ \simeq\
 \mathfrak F_X\longrightarrow\Def_E
\]
と一致する。従って適当な形式座標の下で
\[
 \widehat{\OO}_{C,c}
  \simeq\CC[[a_{ij}]]/(q_{ijkl}),\qquad
 \widehat{\OO}_{U,u}\simeq\CC[[a_{ij}]].
\]
特に $C$ は $c$ で7次元、余次元3であり、$T_cC=T_uU$ は10次元である。
\end{theorem}

\begin{proof}
\textbf{相対圏の構成。}
To\"en--Vaqui\'e\cite[Corollary 3.22]{TV}によれば、
negative自己Extを持たないsimple対象のstackは、
局所有限表示なalgebraic space上の、\'etale局所自明な
$\mathbb G_m$-gerbeである。
stable部分の\'etaleスキームchartを選び、このgerbeを自明化すれば、
universal対象 $\mathcal E_U$ を持つ $U$ が得られる。
安定性はここでは開条件であり、\cite[Proposition 6.3]{FP}の
Clifford加群のGiesekerモジュライによる記述を使っても確認できる。

$U$ をaffineに縮小する。
$\mathcal K$ のcompact generatorの自己準同型dg代数を $A$ とすれば、
$A_U=A\otimes_\CC\OO(U)$ は $\OO(U)$ 上smooth properである。
$\mathcal E_U$ に対応するperfectな $A_U$-加群を使い、
$\mathcal K_U$ と一つの例外対象 $Q$ を上三角に貼り合わせる。具体的には
\[
 \RHom_{\mathcal B_{\mathcal E_U}}(F,Q)
   =\RHom_{\mathcal K_U}(F,\mathcal E_U),\qquad
 \RHom_{\mathcal B_{\mathcal E_U}}(Q,F)=0,\qquad
 \RHom(Q,Q)=\OO_U
\]
と定め、pretriangulatedかつidempotent completeな包を取る。
この相対圏はsmooth properである。
propernessは各Hom複体の $\OO(U)$-perfect性による。
smoothnessは、上三角dg代数の対角bimoduleを、
二つの対角成分のperfectな対角bimoduleとperfectな混合加群から構成できる
ことによる。この構成は体に限らず可換基底環上で成立する。

\textbf{相対simple対象のモジュライ。}
\cite[Corollaries 2.13, 3.15, 3.17, 3.22]{TV}を基底環 $\OO(U)$ 上で適用する。
$\mathcal B_{\mathcal E_U}$ のsimple perfect対象のstackからscalar自己同型を
除いたalgebraic spaceを $\mathcal H$ と書くと、
$f:\mathcal H\to U$ は局所有限表示である。
$g=[G]$ における\emph{相対}接空間は
\[
 T_g(\mathcal H/U)=\Ext^1_{\mathcal B}(G,G)=0.
\]
残差体は $\CC$ なので、局所有限表示とこの消滅から、
$f$ は $g$ でunramifiedである。

$\mathcal H$ の $g$ を通る\'etaleスキームchartを取り、
unramifiedな射の\'etale局所構造
\cite[Lemma 41.17.1, Tag 04HH]{StacksUnramified}を適用する。
$U$ をさらに\'etale局所化すると、指定点を通る開いた枝 $C$ が
$U$ へ閉埋め込みになる。必要に応じて有限型の近傍へ縮小する。

\textbf{完備化の同定。}
この\'etale局所化は、指定した $\CC$-点での形式変形関手を変えない。
$g$ での $\mathcal H$ の完備化は、
$\mathcal E_A$ に沿って $G$ を
$\mathcal B_{\mathcal E_A}$ のsimple perfect対象へ持ち上げる変形関手である。
special fiberでnegative自己Extがなく $\operatorname{End}(G)=\CC$ なので、
局所Artin基底上のperfectな持ち上げはこのsimple locusに入る。
第\ref{sec:gluing-moduli}節によれば、その持ち上げは存在すれば同型類が一意で、
存在条件はちょうど $\mathfrak F_X$ である。
従って完成された閉埋め込みは
$\mathfrak F_X\subset\Def_E$ と一致する。
\end{proof}

$C$ 上のuniversalな $G$ 自身が必要なら、
simple対象の $\mathbb G_m$-gerbeをさらに\'etale局所自明化すればよい。
この命題は指定対象 $G$ を通る局所的な枝を構成している。
相対モジュライの任意の有限型部分の像はconstructibleだが、
その全体のpropernessや大域的な像の閉性はここでは結論しない。
また、一般の形式イデアルの代数化を仮定したのではなく、
実際の相対対象モジュライを先に構成した点がこの証明の要点である。

\subsection{Brill--Noether incidenceの局所構造}

$p\in X$ に対し $P_p=i^*\OO_p\in\mathcal K$ とおく。
相対射影を使えば、これらは $X$ を基底とするperfectな族になる。
$F\in\mathcal K$ に対し、随伴とSerre双対性から
\[
 \Ext^3_{\mathcal K}(P_p,F)=\Ext^3_X(\OO_p,F)
 \simeq\operatorname{Hom}_X(F,\OO_p\otimes\omega_X)^*.
\]
従って
\[
 b_p(F):=\dim_\CC\Ext^3_{\mathcal K}(P_p,F)
       =\dim_\CC\operatorname{Hom}_X(F,\OO_p).
\]
$E$ の場合は、任意の $p$ について
\[
 \dim\Ext^2_{\mathcal K}(P_p,E)=1,\qquad b_p(E)=4.
\]
一方、$F$ が階数3のtorsion-free層なら
\[
 b_p(F)\ge4
 \quad\Longleftrightarrow\quad
 F\text{ は }p\text{ で局所自由でない}.
 \tag{BN}\label{eq:bn-nonfree}
\]
実際、$b_p(F)$ は局所加群 $F_p$ の最小生成元数である。
3個の生成元による $\OO_{X,p}^{\,3}\twoheadrightarrow F_p$ の核は、
torsion-free加群の階数零の部分加群なので零となる。

\begin{theorem}[滑らかなincidenceとその法束]\label{thm:bn-incidence}
命題\ref{thm:algebraic-compatibility}の $U$ を縮小し、
$X\times U$ 上のcohomology-jump部分スキーム
\[
 \mathcal I
   =\{(p,t)\in X\times U:b_p(\mathcal E_t)\ge4\}
\]
をdeterminantalなスキーム構造で定める。
すると $\mathcal I$ は $X\times\{u\}$ の近傍で滑らかな9次元スキームであり、
$\mathcal I_u=X$ はスキームとして成立する。
さらに、固定したtrace線の自明化の下で
\[
 \mathcal N_{X/\mathcal I}
 \simeq
 \ker\bigl(\Ext^1(E,E)\otimes\OO_X\longrightarrow Q(1)\bigr)
 \simeq\Lambda^2Q.
\]
\end{theorem}

\begin{proof}
$E$ の $p$ における導来ファイバーは、次数 $-1$ に1次元、次数0に4次元を持つ。
perfect複体の局所表示から可逆な部分を消去すると、
$(p,u)$ の近傍でuniversal対象は
\[
 [\,\OO\xrightarrow{\,d\,}\OO^{\,4}\,]
 \qquad\text{（次数 }-1,0\text{）}
\]
と表示できる。この表示で $\mathcal I$ は $d$ の4成分の零点である。
special fiberではこれらの成分が全て零なので、$\mathcal I_u=X$ となる。

基底方向の微分はcohomologyのKodaira--Spencer写像
\[
 \Ext^1(E,E)\longrightarrow
 \operatorname{Hom}(\mathcal H^{-1}(E),\mathcal H^0(E))
 =H^0(Q(1))
\]
を $p$ で評価したものである。
この写像は命題\ref{thm:gluing-ext}で計算した同型と一致する。
実際、\eqref{eq:E-triangle}の射を
$\iota:\OO_X(-1)[1]\to E$、$\beta:E\to Q$ と書くと、
その写像は
\[
 \alpha\longmapsto
 \beta[1]\circ\alpha\circ\iota:
 \OO_X(-1)[1]\longrightarrow Q[1].
 \tag{KS}\label{eq:bn-edge-map}
\]
随伴 $\RHom(E,E)\simeq\RHom(E,Q)$ は $\beta$ との合成で与えられ、
三角へ $\RHom(-,Q)$ を適用する次の写像は $\iota$ との合成なので、
これは既に用いた自己Extの同型そのものである。

評価列の外積から
\[
 0\longrightarrow\Lambda^2Q
 \longrightarrow\Lambda^2V\otimes\OO_X
 \longrightarrow Q(1)\longrightarrow0
\]
を得る。従って $H^0(Q(1))=\Lambda^2V$ の評価は各点で全射である。
$d$ の4成分の微分が全射となるので、
$\mathcal I$ は13次元の $X\times U$ 内で滑らかな余次元4の部分スキームである。
法束の計算もこの接空間列と外積の評価列から従う。
$X$ のpropernessにより、必要な縮小を $U$ 上で一斉に行える。
\end{proof}

\subsection{Yoneda積による一次条件とFermatでの接錐}

$\alpha\in\Ext^1(E,E)$ に対して、Yoneda積
\[
 m_{\alpha,p}:\Ext^2_{\mathcal K}(P_p,E)
       \longrightarrow\Ext^3_{\mathcal K}(P_p,E),
 \qquad b\longmapsto\alpha[2]\circ b
\]
を考え、その零点を
$Z_\alpha=\{p\in X:m_{\alpha,p}=0\}$ とする。
これは $\mathcal K$ 内の対象と射だけで書かれた一次Brill--Noether条件である。
局所双対性と導来ファイバーの関手性により、
$m_{\alpha,p}$ は\eqref{eq:bn-edge-map}の $p$ における評価に一致する。

\begin{lemma}[一次条件の階数による記述]\label{lem:bn-rank}
\eqref{eq:bn-edge-map}によって $\alpha$ を
$a\in H^0(Q(1))=\Lambda^2V$ と同一視し、対応する切断を $s_a$ と書く。
$X\subset\PP(V^*)$ の規約の下で
\[
 Z_\alpha=Z(s_a)
   =X\cap\PP\bigl(\ker(a:V^*\to V)\bigr).
\]
$a\ne0$ の階数が2ならこの零点は平面三次曲線であり、
階数が4なら高々1点である。従って集合として
\[
 \{a\in\Lambda^2V:\dim Z(s_a)\ge1\}
   =C(\Gr(2,V)).
 \tag{*}\label{eq:first-bn-cone}
\]
ここでは $a=0$ も右辺に含める。
\end{lemma}

\begin{proof}
$p\in\PP(V^*)$ を評価functionalと見れば、$s_a(p)$ は contraction
$\iota_p(a)$ である。5次元の交代形式の非零の階数は2または4である。
階数2ならkernelの射影化は平面であり、階数4なら1点である。
滑らかなcubic threefoldは平面を含まないので、前者の交叉は次元1の平面三次曲線となる。
例えば平面の理想を $(l_1,l_2)$ として $F=l_1q_1+l_2q_2$ と書けたなら、
平面上の $q_1,q_2$ の共通零点で全ての偏微分が零になり、滑らかさに反する。
最後に階数2以下の交代形式は分解可能な外積であり、
その錐はPl\"ucker埋め込みのGrassmannian錐である。
\end{proof}

一般の $X$ に対して、上の $\Lambda^2V$ との同一視と、
\eqref{eq:gluing-tangent}を通じてPoisson空間から得る同一視の比較には、
次元の一致以上の議論が必要である。
ここではまずFermatの場合の群作用による証明を与える。
全ての滑らかな $X$ に対する比較は、第\ref{sec:general-tangent-comparison}節で証明する。

\begin{theorem}[FermatでのBrill--Noether接錐]\label{thm:fermat-bn-cone}
$X_0=V(x^3+y^3+z^3+u^3+v^3)$ とする。
命題\ref{thm:algebraic-compatibility}の適合部分空間 $C$ の接錐は、
\eqref{eq:bn-edge-map}の座標でPl\"ucker錐である。
その台は圏論的に
\[
 \left|\operatorname{TC}_cC\right|
 =\{\alpha\in\Ext^1(E,E):\dim Z_\alpha\ge1\}
\]
と特徴付けられる。これは10次元の接空間全体ではなく、7次元の接錐についての主張である。
\end{theorem}

\begin{proof}
被覆群
$\widetilde\Gamma=\mu_3^5\rtimes S_5$ の線形作用で計算し、
最後に $\Gamma=(\mu_3^5/\mu_{3,\mathrm{diag}})\rtimes S_5$ へ降ろす。
$V$ の座標を $x_0,\ldots,x_4$ とし、
$g^*x_i=\zeta_i x_i$ とする。Fermat三次式はこの作用で不変なので、
同変なadjunctionは
\[
 \omega_{X_0}\simeq\OO_{X_0}(-2)\otimes\det V
\]
を与える。よってPoisson側の自然な表現は
\[
 H^0(\Lambda^2T_{X_0})
 \simeq H^0(\Omega^1_{X_0}(2))\otimes(\det V)^*
 \simeq\Lambda^2V\otimes(\det V)^*.
\]
最後の同型は、conormal列
$0\to\OO_{X_0}(-1)\to Q(1)\to\Omega^1_{X_0}(2)\to0$
と $\OO_{X_0}(-1)$ のコホモロジーの消滅から従う。

対象側ではtrace線
\[
 L=\Ext^2(Q,\OO_{X_0}(-1))
   \simeq H^3(\OO_{X_0}(-2))
   \simeq(\det V)^*
\]
を保持する。\eqref{eq:E-triangle}の左端は正確には
$\OO_{X_0}(-1)[1]\otimes L^*$ なので、
\[
 \Ext^1(E,E)
 \simeq H^0(Q(1))\otimes L
 \simeq\Lambda^2V\otimes(\det V)^*.
\]
従って二つの側のadjunctionおよびtraceのcharacterは一致する。
各 $(i,j)$ のweight lineのcharacterは
\[
 \chi_{ij}(\zeta)
  =\frac{\zeta_i\zeta_j}{\zeta_0\zeta_1\zeta_2\zeta_3\zeta_4}.
\]
対角 $\mu_3$ 上では $\zeta^{\,2-5}=1$ なので、この表現は $\Gamma$ に降りる。
10個のweightは相異なり、$S_5$ はその直線を推移的に置換する。
従ってこの群作用と可換する自己準同型はscalarだけである。

\eqref{eq:gluing-tangent}は自然な同型なので、
上の二つの同一視の下では非零scalar倍となる。
命題\ref{thm:object-cone}および命題\ref{thm:algebraic-compatibility}の
接錐が、この座標でもPl\"ucker錐であることが従う。
その台の記述は補題\ref{lem:bn-rank}である。
\end{proof}

右辺のHilbert多項式の跳躍条件から別途定めたスキーム構造まで
Pl\"uckerイデアルと一致するとは主張していない。
また、$C$ 自体の全ての点を\eqref{eq:bn-nonfree}の条件で特徴付けたわけではない。
一般の $X$ についてHochschild写像と\eqref{eq:bn-edge-map}を比較する証明は、
命題\ref{thm:general-tangent-comparison}で与える。

\subsection{明示的な固定残余圏族から得る層の変形}

\begin{theorem}[非定数な階数3の層の族と特殊化の制限]\label{thm:gluing-sheaf-arc}
第9節の固定残余圏族 $A'_h$ は、$R=\CC[[h]]$ 上で、
固定した $X_0$ の $\mathcal K$ 内のperfect対象の非定数な族
$\mathcal E_R$ を与える。special fiberは $E$ であり、
その幾何学的generic fiberはGieseker-stableなtorsion-free層で、
\[
 \rk=3,\qquad c_1=0,\qquad c_2=H^2,\qquad c_3=0
\]
を持つ。
generic fiberの非局所自由な点の集合の水平閉包を
$S\subset X_0\times\operatorname{Spec}R$ とすれば、集合として
\[
 S_0\subset
 X_0\cap\{u=v=0\}
 =V(x^3+y^3+z^3)\subset\PP^2.
\]
\end{theorem}

\begin{proof}
既存の形式族は
$\Def_E\simeq\widehat U_u$ への形式弧を与える。
$U$ をaffineに取れば、対応する連続準同型
$\widehat{\OO}_{U,u}\to R$ を $\OO(U)$ に制限して
$\operatorname{Spec}R\to U$ を得る。
universal perfect対象を引き戻した $\mathcal E_R$ は、
全てのArtin切断で元の形式族と一致する実際の $R$ 上の族である。

$A'_h$ の一次Poisson類は非零であり、
\eqref{eq:gluing-tangent}によりこの弧の接ベクトルも非零である。
従ってgeneric pointの像は $e$ ではない。
$U$ をstable locus内に取り、universal対象の局所Tor-amplitudeを
$[-1,0]$ に縮小しておく。
さらに $U$ を縮小し、$e$ の上の点を $u$ だけにしておく。
\cite[Proposition 5.5]{FLZ}によれば、$U\setminus\{u\}$ の閉点は
Gieseker-stableな層を表す。このamplitudeと正の階数によりshiftは零である。
層であることとGieseker-stabilityの開性、および有限型 $\CC$-スキームの
Jacobson性から、この結論は $U\setminus\{u\}$ の全ての点に成り立つ。
従って非定数な弧のgeneric fiberについても同じ結論が従う。
purityと正の階数からtorsion-freeであり、Chern類はclass $3v$ から得られる。

命題\ref{thm:fermat-bn-cone}の座標比較により、この弧の接方向は、
非零scalar倍を除いて $a=u\wedge v$ である。
命題\ref{thm:bn-incidence}の局所二項表示では、微分は
$h\,s_a+O(h^2)$ と書ける。
generic fiberの非局所自由な部分はその微分の4成分で定義される。
水平閉包のイデアルはこれらの成分を $h$ で割ったものを含むので、
$h=0$ で $s_a$ は零になる。従って $S_0\subset Z(s_a)$。
最後に
\[
 \iota_p(u\wedge v)=u(p)v-v(p)u
\]
から $Z(s_a)=X_0\cap\{u=v=0\}$ を得る。
\end{proof}

この段階では非局所自由部分の特殊化の包含までを得た。
第\ref{sec:evaluation-gluing}節では、この族のgeneric fiberを評価核として記述する。
それにより、非局所自由部分が固定された平面三次曲線そのものであり、
局所自由でもUlrichでもないことを証明する。

\section{非被約点からの平行移動パラメータの復元}\label{sec:reconstruction}

第\ref{sec:points}節の局所代数 $J_a$ が記録する情報を、
Sklyaninの平行移動パラメータと比較する。
まず代数閉体 $k$、$\operatorname{char}k=0$ 上の幾何学的ファイバーを扱い、
最後に $\CC[[h]]$ 上の形式族へ戻る。
以下の一般ファイバーについての個数は幾何学的な個数であり、
$\CC((h))$ 上の非自明なformまで $j$ だけで分類するという主張ではない。
本節では $t=a^3$ と書き、第9節の三乗根は $r$ と書く。

\subsection{Jacobian環のassociated formとしての三次形式}

\[
 f_a=x^3+y^3+z^3-3axyz,\qquad
 J_a=k[x,y,z]/(x^2-ayz,y^2-azx,z^2-axy)
\]
とおく。$a^3\ne1$ のとき $f_a$ は滑らかな平面三次曲線を定め、
$J_a$ はそのJacobian環である。

\begin{theorem}[associated formと逆数関係]\label{thm:fat-associated}
$a^3\ne1$ とする。局所代数 $J_a$ のsocle三次形式
\[
 \Phi_a(X,Y,Z)=a(X^3+Y^3+Z^3)+6XYZ
\]
は、$f_a$ のMacaulay inverse system、従ってそのassociated formである。
さらに $a\ne0$、$a^3\ne-8$ のとき、通常の楕円曲線の $j$ の規約で
\begin{equation}\label{eq:associated-reciprocal-j}
 j(V(f_a))=\frac{27a^3(a^3+8)^3}{(a^3-1)^3},\qquad
 j(V(f_a))\,j_{\mathrm{fat}}(a)=1728^2.
\end{equation}
\end{theorem}

\begin{proof}
$\partial_x f_a=3(x^2-ayz)$ とその巡回置換からJacobian環の表示を得る。
第\ref{sec:points}節の積表を使い、$T=xyz$ とおくと
\[
 (Xx+Yy+Zz)^3=\Phi_a(X,Y,Z)T.
\]
また $(\partial_X^2-a\partial_Y\partial_Z)\Phi_a=0$ とその巡回置換が成り立つ。
$\Phi_a$ のapolar algebraのHilbert関数は $(1,3,3,1)$ であり、
$J_a$ も同じ長さ8を持つので、この三つの二次作用素がそのannihilatorを生成する。
これがinverse systemとしての同一視を与える。
$j$ の式はHesse三次曲線の公式と\eqref{eq:fat-j}から従う。
これはAlper--Isaev--Kruzhilin\cite[Section 3.2, (3.11), (3.18)]{AIK}の
associated-form involutionの具体例である。
同論文の規格化 $J=j/1728$ では、この関係は $J\mapsto1/J$ と書かれる。
\end{proof}

associated formはHessianとは区別される\cite[Remark 3.1]{AIK}。
実際、直接微分すると
\begin{equation}\label{eq:fat-hessian-distinction}
 \operatorname{Hess}(f_a)
 =-54a^2(x^3+y^3+z^3)+54(4-a^3)xyz
\end{equation}
となり、$\Phi_a$ とは異なる。
この逆数関係を、固定曲線 $E_0$ と $V(\Phi_a)$ の間の
3-isogenyと解釈することはできない。
後で構成する3-isogenyは $V(\Phi_a)$ 自身を始域とする。

\subsection{固定Fermat曲線の次数2の商}

Sklyaninの平行移動の原点を
\[
 E_0=V(x^3+y^3+z^3),\qquad O=[1:-1:0]
\]
とする。第9節のSklyaninパラメータは
\[
 p(h)=[1:-r:-h],\qquad r^3+h^3=1,
 \qquad a=\frac{1-r}{h}=-\frac{x+y}{z}
\]
であり、点曲線の自己同型は $p(h)$ による平行移動である
\cite[Section 2.3]{WWY}。

\begin{theorem}[markedな係数からの復元]\label{thm:translation-double-cover}
有理関数 $a=-(x+y)/z$ は次数2の射
$E_0\to\PP^1_a$ に延長し、$E_0/\{\pm1\}$ を与える。
その分岐値は $a=0$ および $a^3=4$ である。
一般の $a$ から平行移動点を復元する式は
\begin{equation}\label{eq:translation-quadratic}
 (1-a^3)h^2+3a^2h-3a=0,\qquad r=1-ah.
\end{equation}
従って $a$ は平行移動点を符号まで決定する。
\end{theorem}

\begin{proof}
$\omega$ を原始三乗根とし、$Q_\pm=[1:-\omega^{\pm1}:0]$ とおく。
直線 $x+y=0$ は $O$ における変曲接線であり、
$z=0$ の交点は $O,Q_+,Q_-$ なので
\[
 \operatorname{div}(a)=2O-Q_+-Q_-.
\]
従って $a$ は次数2である。
$[x:y:z]\mapsto[y:x:z]$ は $a$ と $O$ を固定する非自明なinvolutionなので、
楕円曲線の $[-1]$ に等しい。
その固定点は $O$ と $x=y$ 上の三点であり、後者では $a^3=4$ となる。
また $r=1-ah$ を $r^3+h^3=1$ に代入し、一般点で $h$ を消去すると
\eqref{eq:translation-quadratic}を得る。その判別式は
\[
 9a^4+12a(1-a^3)=3a(4-a^3).
\]
二根の交換は $(r,h)\mapsto(1/r,-h/r)$、すなわち上の $[-1]$ である。
\end{proof}

ここでの分岐値 $a^3=4$ と、socle三次曲線が特異になる $a^3=-8$ は異なる。
前者は平行移動点の符号を忘れることに由来し、後者は局所代数の三次形式に由来する。

\subsection{粗い局所代数からは四つのSklyanin同型類が残る}

\begin{theorem}[幾何学的一般点における有限多価の復元]\label{thm:four-sklyanin-classes}
$t=a^3$ とおくと、$E_0\to\PP^1_t$ は
$\operatorname{Aut}(E_0,O)\simeq\mu_6$ による商である。
socle三次曲線の不変量は
\begin{equation}\label{eq:fat-j-degree-four}
 \mathcal J(t)=\frac{110592(t-1)^3}{t(t+8)^3}
\end{equation}
という次数4の写像であり、合成 $E_0\to\PP^1_j$ の次数は24である。

同じ埋め込まれた曲線 $E_0$ に付随するSklyanin代数を $S(p)$ と書く。
$p,q\in E_0\setminus E_0[3]$ に対して
\begin{equation}\label{eq:fixed-sklyanin-isom}
 S(p)\simeq_{\mathrm{gr}}S(q)
 \quad\Longleftrightarrow\quad
 q\in\operatorname{Aut}(E_0,O)\cdot p.
\end{equation}
従って、一般の非被約点の局所代数の同型類からは、
平行移動点として24個、Sklyanin代数の次数付き同型類として4個の候補が残る。
\end{theorem}

\begin{proof}
$[x:y:z]\mapsto[y:x:z]$ と $[x:y:z]\mapsto[x:y:\omega z]$ は
$\operatorname{Aut}(E_0,O)$ を生成し、$a$ に対してそれぞれ
$a\mapsto a$、$a\mapsto\omega^{-1}a$ と作用する。
$t=a^3$ は不変で次数6なので、この群の商座標である。
\eqref{eq:fat-j-degree-four}の分子と分母は互いに素であり、その最大次数は4である。

幾何代数の同型判定\cite[Theorem 2.8]{IMRelations}によれば、
次数付き同型は点曲線の平行移動自己同型の射影的共役に対応する。
$\OO_{E_0}(1)=\OO_{E_0}(3O)$ なので
\[
 \operatorname{Aut}(\PP^2,E_0)
 =E_0[3]\rtimes\operatorname{Aut}(E_0,O)
\]
である\cite[Theorem 4.12]{IMRelations}。
平行移動を $T_p$ と書くと、$T_s\alpha$ による共役は
\[
 (T_s\alpha)T_p(T_s\alpha)^{-1}=T_{\alpha(p)}
\]
だから\eqref{eq:fixed-sklyanin-isom}を得る。
これは\cite[Theorem 4.16 and Lemma 4.17]{IMRelations}の $i=0$ の場合でもある。

最後に、$\Phi_a$ が滑らかな一般点では、$J_a$ の局所代数としての同型類と
$\Phi_a$ の射影同値類は同値である。
一方向は極大イデアルの次数付き環と三重積を取り、逆方向はinverse systemを使う。
滑らかな平面三次曲線の射影同値類は、代数閉体上で $j$ により決まる。
実際、同じ $j$ を持つ楕円曲線間の同型を平行移動で調整すれば、
二つの次数3の線形束を一致させられる。
これは $[3]:E\to E$ の全射性による。
従って一般の局所代数の同型類は $\mathcal J$ の一般ファイバーを与え、
そこには四つの $t$、それぞれに六つの平行移動点がある。
\end{proof}

候補は具体的に
\[
 j\,t(t+8)^3-110592(t-1)^3=0
\]
を解いて求められる。さらに一つの $t$ を選ぶと
\[
 (t-1)r^2+(t+2)r+(t-1)=0,\qquad h^3=1-r^3
\]
から六つの点を復元できる。
最初の二次式の根は互いに逆数で、判別式は $3t(4-t)$ である。
ここまでの分類は三生成元のSklyanin代数についてのものである。
第\ref{sec:five-generator}節では、これを五生成元の三次元代数 $A'_h$ の
次数付き同型分類へ拡張する。

\subsection{socle三次曲線の3-isogenyと \texorpdfstring{$X_0(3)$}{X0(3)}}

$a\ne0$ として $b=-2/a$ と置けば、$V(\Phi_a)$ は
\[
 C_b:\ x^3+y^3+z^3-3bxyz=0
\]
と書ける。以下 $b^3\ne1$ とする。

\begin{theorem}[明示的な巡回3-isogeny]\label{thm:fat-three-isogeny}
$C_b$ の自己同型
\[
 \delta[x:y:z]=[x:\omega y:\omega^2z]
\]
は非零の3-torsion点による平行移動であり、その商は
\begin{equation}\label{eq:fat-isogeny-quotient}
 [x:y:z]\longmapsto[X:Y:Z:T]=[x^3:y^3:z^3:xyz],\qquad
 X+Y+Z=3bT,\quad XYZ=T^3
\end{equation}
で与えられる。商曲線の $j$ は
\begin{equation}\label{eq:fat-isogeny-j}
 j_3(t)=-\frac{110592(t+9)^3}{t^3(t+8)}.
\end{equation}
幾何学的一般点で $t$ は、楕円曲線と選ばれた位数3の巡回部分群の組の
粗モジュライ $X_0(3)$ の座標となる。
$s=-27t/(t+8)$ と書けば、二つの $j$-写像は
\begin{equation}\label{eq:x03-j-pair}
 j_{\mathrm{fat}}=\frac{(s+3)^3(s+27)}s,\qquad
 j_3=\frac{(s+27)(s+243)^3}{s^3}.
\end{equation}
双対isogenyへ移るFricke involutionは
\begin{equation}\label{eq:fat-fricke}
 s\longmapsto\frac{729}s,
 \qquad\text{すなわち}\qquad t\longmapsto-t-8
\end{equation}
である。
\end{theorem}

\begin{proof}
$\delta$ の射影空間内の固定点は三つの座標頂点だけで、どれも $C_b$ 上にはない。
楕円曲線上の有限位数の自己同型が平行移動でなければ、
$1-\alpha$ の全射性から固定点を持つ。従って $\delta$ は位数3の平行移動である。
\eqref{eq:fat-isogeny-quotient}の一般ファイバーは $\delta$ の三点軌道だから、
これはその巡回部分群による商である。
$Z$ を消去すると $XY(3bT-X-Y)=T^3$ を得る。
$x_W=-T$、$y_W=X$、$z_W=Y$ とすればWeierstrass形は
\[
 y_W^2z_W+3b x_Wy_Wz_W+y_Wz_W^2=x_W^3
\]
となり、その不変量は
\[
 c_4=9b(9b^3-8),\qquad\Delta=27(b^3-1).
\]
従って $j=c_4^3/\Delta$ に $b^3=-8/t$ を代入して
\eqref{eq:fat-isogeny-j}を得る。
商曲線の原点を始域の原点の像に選べば、この次数3の射はisogenyである。
Weierstrass表示で選んだ変曲点を原点に変更しても $j$ は変わらない。

$a\mapsto\omega a$ に対して $b\mapsto\omega^2b$ となり、
$[x:y:z]\mapsto[\omega x:y:z]$ が曲線と巡回対称性を同一視する。
従ってこの組の粗モジュライは $t=a^3$ に降下する。
一般の楕円曲線の $E[3]\simeq(\mathbb Z/3)^2$ には位数3の部分群が四つあり、
一般の自己同型 $\pm1$ はそれぞれを保つ。
従って $X_0(3)\to\PP^1_j$ の次数は4である。
$t$ の写像も次数4なので、$\PP^1_t$ からこの粗モジュライへの写像は次数1となる。

\eqref{eq:x03-j-pair}は直接代入で従う。
また $s\mapsto729/s$、すなわち $t\mapsto-t-8$ は二つの $j$ を交換する。
幾何学的一般点では異なる位数3の部分群から得られる商の $j$ は異なるので、
これが双対isogenyに対応する。
この最後の事実は、一般の楕円曲線の自己準同型環が $\mathbb Z$ であることからも分かる。
二つの異なる商が同型なら、一方のisogenyと他方の双対を合成して得る次数9の
自己準同型は $\pm[3]$ であり、双対を取ることで二つの核が一致してしまう。
\end{proof}

Hesse pencilの共通する3-torsionとそのsymplectic構造については
\cite[Theorem 1.1 and Sections 2--4]{ATT}を参照されたい。
ここで $X_0(3)$ が記録する巡回部分群は、socle三次曲線上に選んだ対称性である。
固定Sklyanin曲線 $E_0$ は、平行移動パラメータ $p$ を動かすための別の曲線である。

\subsection{形式近傍と、一般のSklyaninパラメータ}

$O$ の形式近傍では
\[
 t=a(h)^3=\frac{h^6}{27}+\frac{h^9}{27}+O(h^{12}),\qquad
 \frac1{\mathcal J(t)}=-\frac{t}{216}+O(t^2).
\]
後者は $t$ に関して逆形式級数を持つ。
従って $O$ に近づく枝を指定すれば、$1/j_{\mathrm{fat}}$ から $t$ は一意に復元される。
大域的な四つの $t$ の曖昧さはこの形式近傍では消え、形式商の座標環は
\[
 \widehat{\OO}_{E_0,O}^{\mu_6}=\CC[[t]]
\]
となる。整域上の非零の形式弧を復元するときに残る曖昧さは $\mu_6$ の作用である。
ただし粗い商のArtin値点は、群作用による軌道の関手とは異なる。
例えば $\CC[\epsilon]/(\epsilon^2)$ 上の $h=0$ と $h=\epsilon$ は
同じ $t=0$ を与えるが、同じ $\mu_6$ 軌道には属さない。

一般のSklyaninパラメータ $[A:B:C]$ でも、$C=-1$ と規格化して
\[
 d=A+B,\qquad w=AB,\qquad
 \tau=\frac{A^3+B^3+C^3}{ABC}
\]
とおけば、$A^3+B^3=d^3-3wd$ より
\begin{equation}\label{eq:general-sklyanin-recovery}
 w=\frac{1-d^3}{\tau-3d},\qquad
 A,B:\ U^2-dU+\frac{1-d^3}{\tau-3d}=0
\end{equation}
を得る。点曲線は $x^3+y^3+z^3-\tau xyz=0$、
対角上の二次関係は $J_d$ のものである。
従ってmarkedな $(\tau,d)$ は、一般点で $A\leftrightarrow B$ までパラメータを復元する。
この二次式の判別式は
\[
 \frac{d^3+\tau d^2-4}{\tau-3d}
\]
である。元のFermat族では $d=0$、$\tau=-\lambda^2$、
固定残余圏族では $\tau=0$、$d=a(h)$ となる。
両族で異なる不変量が変化することを、この座標でまとめて記述できる。

\subsection{三つの検出次数：貼り合わせ、局所環、粗モジュライ}

第9節の固定残余圏族を $R=\CC[[h]]$ 上で考える。
第13節の水平な非被約成分の局所環は、$a=a(h)$ とした $J_a$ である。
この成分のflat性と、相対point-module truncation全体の非flat性は区別する。

\begin{theorem}[固定残余圏族で最初に見える次数]\label{thm:orders-126}
この族に沿って、次の三つの次数が得られる。
\begin{enumerate}
 \item gluing object $E_h$ の変形の一次類は非零である。
 \item $J_{a(h)}$ のmarkedな可換局所代数としての変形は
 $\bmod\,h^2$ では定数であるが、$\bmod\,h^3$ では定数でない。
 \item socle三次形式の粗モジュライの局所座標 $\tau=1/j_{\mathrm{fat}}$ は
 \[
   \tau=-\frac{h^6}{5832}+O(h^9)
 \]
 であり、最初の非零項は6次である。
\end{enumerate}
\end{theorem}

\begin{proof}
第一の主張は第11節の接写像の同型と、非零の一次Poisson類による。
第二を確認する。まず $a=\epsilon$、$\epsilon^2=0$ とした
$J_\epsilon$ のKodaira--Spencer類は、complete intersection
$J_0=\CC[x,y,z]/(x^2,y^2,z^2)$ の余接複体によって
\[
 T^1_{J_0}=\operatorname{coker}
 \bigl[J_0^3\xrightarrow{\operatorname{diag}(2x,2y,2z)}J_0^3\bigr]
\]
における $(-yz,-zx,-xy)$ で表される。
$yz\notin xJ_0$ なのでこれは非零である。

$a(h)=h^2/3+O(h^5)$ であるから、$\bmod\,h^2$ では定数である。
さらに、$\bmod\,h^3$ で定数族へのmarkedな同型があったと仮定する。
定数族の生成元 $x$ の像は
$x+h f+h^2g$、$f,g\in J_0$ と書ける。
その平方の $h$ の係数は $2xf$ なので、$f\in\operatorname{ann}_{J_0}(x)=xJ_0$。
従って $f^2=0$ である。$h^2$ の係数は
\[
 \frac13yz+2xg
\]
となり、$yz\notin xJ_0$ に反する。よって定数化できない。
第三は\eqref{eq:fat-j}の逆数と $a(h)$ の展開から従う。
\end{proof}

従って、同じ非可換族に対しても、貼り合わせ対象、水平な非被約成分の局所環、
そのsocle三次形式の粗モジュライは、それぞれ1次・2次・6次から変化を検出する。
これは普遍的な写像 $\Def_E\to\operatorname{Def}_{J_0}$ を構成したという意味ではない。
ここで比較したのは、次数付き代数模型も指定された一本の形式族である。

\section{一般のcubic threefoldにおける二つの接写像の比較}
\label{sec:general-tangent-comparison}

第\ref{sec:bn-algebraic}節でFermatの場合に証明した接錐の同定を、
全ての滑らかな三次三次元多様体へ拡張する。
ここで比較するのは、Hochschild gluingから得る
\[
 \kappa_F:H^0(X,\Lambda^2T_X)\xrightarrow{\sim}\Ext^1(E,E)
\]
と、cohomologyのKodaira--Spencer写像
\[
 e_F:\Ext^1(E,E)\xrightarrow{\sim}H^0(Q(1))\otimes L_F,
 \qquad L_F=\Ext^2(Q,\OO_X(-1))
\]
である。後者は\eqref{eq:bn-edge-map}のtrace線を保持した形である。
conormal列とadjunctionが与える同型を
\[
 \lambda_F:H^0(Q(1))\otimes L_F
       \xrightarrow{\sim}H^0(X,\Lambda^2T_X)
\]
と書く。$L_F\simeq H^3(X,\OO_X(-2))$ の規約を一貫して用いる。

\begin{theorem}[一般の $X$ における接座標の比較]
\label{thm:general-tangent-comparison}
HKR、連結射、およびadjunctionの符号・正規化を固定すると、
$F$ に依存しない定数 $c\in\CC^\times$ が存在して
\[
               \lambda_F\circ e_F\circ\kappa_F
                    =c\,\id_{H^0(\Lambda^2T_X)}
\]
となる。従って任意の滑らかなcubic threefoldに対して、
命題\ref{thm:algebraic-compatibility}の適合部分空間 $C$ の接錐は、
edge写像 $e_F$ の座標でもPl\"ucker錐である。
特に、その台は
\[
 |\operatorname{TC}_cC|
 =\{\alpha\in\Ext^1(E,E):\dim Z_\alpha\ge1\}
\]
である。
\end{theorem}

証明では、滑らかな三次式の開集合だけで同変性を使うことはしない。
その開集合上には判別式を分母に持つ正則関数が存在するためである。
まず、特異なものも含む全ての非零三次式の上に比較写像を延長する。

\subsection{全非零三次式上の相対gluing}

$S=\Sym^3V\setminus\{0\}$ とし、
$p:\mathcal X\to S$ を $\PP(V^*)\times S$ 内の普遍三次超曲面とする。
この射はflat、proper、Gorensteinである。
以下のHochschild複体と自己Extは全て基底 $S$ に相対的なものとする。
以下の $\Perf(\mathcal X/S)$ は、$\mathcal X$ 上でperfectな複体の圏を
$S$-線形dg圏として見るという意味であり、単に $S$-perfectな複体全体を意味しない。

\begin{lemma}[特異三次式を含む普遍的な持ち上げ]
\label{lem:universal-tangent-lift}
$S$ 上で相対圏
\[
 \mathcal D=\Perf(\mathcal X/S)
       =\langle\mathcal K,Q,\OO_{\mathcal X}\rangle,
 \qquad\mathcal B=\langle\mathcal K,Q\rangle
\]
と $E=i^!Q$、$G=j^!\OO_{\mathcal X}$ を構成できる。
$L=R^3p_*\OO_{\mathcal X}(-2)$ とおくと、相対自己Ext層は
\[
 \mathcal E xt^{0,1}(E,E)=\OO_S,\ \Lambda^2V\otimes L,
 \qquad
 \mathcal E xt^{0,2}(G,G)=\OO_S,\ \Sym^2V\otimes L
\]
であり、記載した以外の次数では零である。
gluingの連結射
\[
 \delta:\Lambda^2V\otimes L
      \longrightarrow\mathcal H^2CC^\bullet(\mathcal B/S)
\]
は、一意に
\[
 \widetilde\delta:\Lambda^2V\otimes L
      \longrightarrow\mathcal H^2CC^\bullet(\mathcal D/S)
\]
へ持ち上がる。
\end{lemma}

\begin{proof}
$\OO,\OO(1),Q$ の例外性は、射影空間のEuler列と
三次超曲面の制限列から従い、超曲面の滑らかさを必要としない。
また $\mathcal X/S$ はGorensteinなので、perfect対象に対する相対Serre双対性がある。
従って第\ref{sec:gluing-objects}節と同じmutationが定義され、
trace線を保持すると
\[
 \OO(-1)[1]\otimes L^*\longrightarrow E\longrightarrow Q,
 \qquad
 \OO(-2)[2]\otimes L^*\longrightarrow G\longrightarrow\OO
\]
という三角を得る。Euler列と低次数の制限列の計算は相対基底上でも成立し、
これらの三角から主張した自己Ext層を得る。
それらは局所自由であり、任意の基底変換と可換である。

Hochschild gluing triangleは、上三角dg圏のbar複体において
混合cochainを分けることで得られる。
この構成は任意の可換基底環上で成立し、smoothnessを仮定しない
\cite[Section 4.5, (4.5.1)--(4.5.2)]{KellerHoch}。
$\mathcal D=\langle\mathcal B,\OO\rangle$ に適用し、
$\mathcal E xt^1(G,G)=0$ を用いると、完全列
\[
 0\longrightarrow\mathcal H^2CC^\bullet(\mathcal D/S)
 \longrightarrow\mathcal H^2CC^\bullet(\mathcal B/S)
 \longrightarrow\Sym^2V\otimes L
\]
を得る。ここで $\Perf(S)$ の相対Hochschild複体は $\OO_S$ のみである。

$\delta$ と最後の射との合成は、局所自由層の間の射
\[
 \chi:\Lambda^2V\otimes L\longrightarrow\Sym^2V\otimes L
\]
である。滑らかな三次式の開集合上では、
第\ref{sec:gluing-objects}節の二段階のgluing計算により
$\HH^2(X)\to\HH^2(\mathcal B)$ は同型であるから、$\chi$ は零である。
$S$ は整かつ被約であり、この開集合は稠密なので、$\chi=0$ は $S$ 全体で成立する。
完全列から $\widetilde\delta$ の存在と一意性を得る。

相対複体の扱いを明記しておく。
$CC^\bullet(\mathcal D/S)$ は対角の相対自己Extと比較できる
\cite[Section 8.1, Theorem 8.9]{ToenMorita}。
普遍族はflatかつseparatedであり、対角の構造層はpseudo-coherentである。
従って各次数の相対Hochschild cohomologyはcoherentであり、flatな局所化と可換である。
$\mathcal B,\mathcal K$ についてもgluing triangleを使って局所化できる。
上の完全列は各affine基底上で構成して層化したものである。
特異ファイバーでHochschild cohomologyの任意の基底変換が同型になるとは仮定しない。
必要なのは自然な基底変換射と、上で計算した二つの自己Extの基底変換だけである。
\end{proof}

\subsection{反対称化とambient双ベクトル}

特異な可換スキームに対しても、Hochschildの二次類には
一次交換子を取る自然な射
\[
 \mathcal H^2CC^\bullet(\mathcal D/S)
 \longrightarrow
 p_*\mathcal H om_{\OO_{\mathcal X}}
   (\Lambda^2\Omega^1_{\mathcal X/S},\OO_{\mathcal X})
\]
がある。局所Hochschild 2-cocycle $\mu$ に対して
$\mu(f,g)-\mu(g,f)$ を取る射である。
正確には、対角のExtに対する局所大域スペクトル系列のedge射で
局所の次数2のExt層の大域切断へ送り、その上で反対称化する。
結合性の一次式からこれはbiderivationになり、
1-cochainのcoboundaryは可換代数上で対称なので、この射はcohomology上で定義される。
重なりでのcocycleの差もcoboundaryなので、局所のbiderivationは貼り合わさる。

この相対記述は通常のbar模型で確認できる。
$S$ を $F(v)\ne0$ の開集合で覆い、固定した線形座標変換で $v=e_0$ とする。
この開集合上の $\mathcal X$ は $x_j\ne0$、$1\le j\le4$ のaffine開集合で覆われる。
各座標環とその交叉は、$x_0$ についてmonicな三次式の商なので基底環上freeである。
従ってrelative bar/\v Cech模型で上の反対称化を行える。
ここでは特異スキームについてのHKR同型を用いていない。

$\Omega^1_{\PP^4/S}|_{\mathcal X}\twoheadrightarrow\Omega^1_{\mathcal X/S}$ から、
biderivationをambient双ベクトルへ単射で送れる。
さらに
\[
 \Lambda^2T_{\PP^4}\simeq\Omega^2_{\PP^4}(5)\otimes(\det V)^*,
 \qquad H^*(\PP^4,\Omega^2_{\PP^4}(2))=0
\]
なので、制限写像は
\[
 W\otimes\OO_S\xrightarrow{\sim}
 p_*(\Lambda^2T_{\PP^4}|_{\mathcal X}),
 \quad W=H^0(\PP^4,\Omega^2_{\PP^4}(5))\otimes(\det V)^*
\]
という同型である。消滅はEuler列の外積の大域切断が
$0\to\Lambda^2V\to V\otimes V\to\Sym^2V\to0$ となることから直接確認できる。

補題\ref{lem:universal-tangent-lift}の持ち上げと合成して、
全ての非零三次式上で正則な、$GL(V)$-同変射
\[
 \Psi:\Lambda^2V\otimes L\longrightarrow W\otimes\OO_S
\]
を得る。滑らかなファイバーで、ambientへの単射を $j_F$ と書けば、
符号とHKRの一定の正規化因子を除いて
\[
                         \Psi_F=j_F\circ\kappa_F^{-1}\circ e_F^{-1}
\]
である。すなわち、これはgluingの逆写像のHKR双ベクトルをambient空間へ送る写像である。

\subsection{三次式への一次性と表現の重複度1}

\begin{proof}[命題\ref{thm:general-tangent-comparison}の証明]
普遍超曲面の完全列
\[
 0\longrightarrow\OO_{\PP^4}(-5)
 \xrightarrow{\ F\ }\OO_{\PP^4}(-2)
 \longrightarrow\OO_{X_F}(-2)\longrightarrow0
\]
の境界により、
\[
 \partial_F:L_F\xrightarrow{\sim}
 H^4(\PP^4,\OO(-5))\simeq(\det V)^*
\]
と自明化する。この境界は
\[
                         \partial_{\lambda F}=\lambda^{-1}\partial_F
\]
を満たす。実際、上の完全列で中央と右端を恒等に保つ比較射は、
左端に $\lambda^{-1}$ を掛ける射である。

$F$ と $\lambda F$ は同じ超曲面と同じ内在的gluingを与えるので、
この自明化で $\Psi_F$ を
$\Lambda^2V\otimes(\det V)^*\to W$ と書けば
\[
                            \Psi_{\lambda F}=\lambda\Psi_F
\]
となる。
$\Psi$ は $\Sym^3V\setminus\{0\}$ 全体で正則であり、
この35次元affine空間から除いた点の余次元は2以上である。
従って各行列成分は原点へ正則に延長する。
上のhomogeneityから各成分は $F$ の係数について一次式である。
よって $\Psi$ は $GL(V)$-同変な線形射
\[
 \Sym^3V\otimes\Lambda^2V\otimes(\det V)^*
                   \longrightarrow W
\]
に対応する。

Euler--Koszul列とPieri則によって
\[
 H^0(\PP^4,\Omega^2(5))\simeq\mathbb S_{(3,1,1)}V,
 \qquad
 \Sym^3V\otimes\Lambda^2V
  \simeq\mathbb S_{(4,1)}V\oplus\mathbb S_{(3,1,1)}V.
\]
特に目標の既約表現の重複度は1である。
ここで最初の表現は、完全列
\[
 0\longrightarrow H^0(\Omega^2(5))
 \longrightarrow\Lambda^2V\otimes\Sym^3V
 \longrightarrow V\otimes\Sym^4V
 \longrightarrow\Sym^5V\longrightarrow0
\]
の核としても得られる。
従って、この形の非零の同変射はscalar倍を除いて一意である。

標準的な同変射を具体的に書く。
$a=v\wedge w$ に対して $\eta_a=v\,dw-w\,dv$ とし、線形に延長する。
$\eta_a$ は $\Omega^1_{\PP^4}(2)$ の切断である。
affine錐のEulerベクトル場を $\mathsf E$ とすると
\[
 \iota_{\mathsf E}\eta_a=0,\quad
 \iota_{\mathsf E}d\eta_a=2\eta_a,\quad
 \iota_{\mathsf E}dF=3F.
\]
従って
\[
                 \Omega(F,a)=3F\,d\eta_a-2\,dF\wedge\eta_a
\]
は水平でweight 5の二形式であり、$H^0(\Omega^2_{\PP^4}(5))$ に属する。
これは $F,a$ に双線形かつ $GL(V)$-同変である。
$X_F$ 上では
\[
                         \Omega(F,a)|_{X_F}=-2\,dF\wedge\eta_a
\]
となる。conormal列の外積とadjunctionにより、これは
$\lambda_F(a)$ のambient双ベクトルへの像の一定非零scalar倍である。
滑らかな $F$ では $\lambda_F$ は同型なので、この同変射は非零である。

重複度1により $\Psi$ はこの射の一定scalar倍となる。
また滑らかなファイバーで $\kappa_F^{-1}$ は同型であり、
HKR反対称化とambientへの単射も双ベクトル上で単射なので、$\Psi$ は非零である。
従って二つの写像を結ぶscalarは非零であり、
$\lambda_F e_F\kappa_F=c\,\id$ を得る。

命題\ref{thm:object-cone}の形式座標の一次部分は $\kappa_F$ である。
従ってedge座標への変更もPl\"ucker錐をscalar倍で保ち、
接錐はスキームとして同じPl\"ucker錐になる。
その台のBrill--Noetherによる記述は補題\ref{lem:bn-rank}から従う。
\end{proof}

この証明は、一般の $X$ について二つの接座標の比較を与える。
一方、適合部分空間 $C$ の高次の方程式や大域的な閉包を、
Brill--Noether条件だけで特徴付けたわけではない。
また、零点集合の次元跳躍から別途作ったスキーム構造とPl\"uckerイデアルの
一致までは主張していない。
第\ref{sec:picard-compatibility}節では別の全次数の持ち上げを用い、
適合部分空間そのものを相対Picardの形式像として記述する。

\subsection{任意の適合形式弧への特殊化の帰結}

\begin{theorem}[一般の $X$ と任意の非定数な適合形式弧]
\label{thm:general-compatible-arc-bn}
$R=\CC[[h]]$ とし、$C$ を通る非定数な形式弧を取る。
Pl\"ucker形式座標でその最初の非零項を
\[
 a(h)=h^\nu a_\nu+O(h^{\nu+1}),\qquad a_\nu\ne0
\]
と書く。すると $a_\nu$ は分解可能である。
対応する実際の $R$-perfect対象 $\mathcal E_R$ の幾何学的一般ファイバーは、
class $3[I_\ell]$ のGieseker-stableなtorsion-free層である。
その非局所自由な部分をFittingイデアルで定め、
そのスキーム論的水平閉包を $S_R\subset X\times\operatorname{Spec}R$ とすると、
\[
 (S_R)_0\ \subset\
 C_0:=X\cap\PP\bigl(\ker(a_\nu:V^*\to V)\bigr)
\]
は閉部分スキームの包含である。$C_0$ は平面三次曲線である。

さらに $C_0$ が滑らかで、一般ファイバーの非局所自由な部分が1次元なら、
$(S_R)_0=C_0$ であり、$S_R/R$ は平面三次曲線の滑らかな族になる。
\end{theorem}

\begin{proof}
$q(a(h))=0$ の最初の非零となり得る係数から $q(a_\nu)=0$ を得る。
従って $a_\nu$ は分解可能である。
命題\ref{thm:algebraic-compatibility}のuniversal対象を弧へ引き戻せば、
実際の $R$-perfect対象が得られる。
非定数性と\cite[Proposition 5.5]{FLZ}から、
命題\ref{thm:gluing-sheaf-arc}と同じ議論で一般ファイバーは主張した層になる。

局所二項表示 $[\OO\xrightarrow d\OO^4]$ において、
命題\ref{thm:general-tangent-comparison}から、
非零scalarと局所自明化を除けば
\[
                         d=h^\nu s_{a_\nu}+O(h^{\nu+1})
\]
である。一般ファイバーの非局所自由部分のFittingイデアルは
$d$ の四つの成分で生成される。
水平閉包はこのイデアルの $h$-saturationで定まるので、
$d$ の各成分を $h^\nu$ で割ったものを含む。
特殊化により $S_0\subset Z(s_{a_\nu})=C_0$ をスキームとして得る。

$h$-saturationにより $\OO_{S_R}$ は $R$-torsion-free、従って $R$-flatである。
$S_R$ はproperで、1次元の一般ファイバーを持つなら特殊ファイバーも1次元である。
滑らかな平面三次曲線 $C_0$ は整であり、その1次元の閉部分スキームは
$C_0$ 全体に等しい。従って $S_0=C_0$。
flatnessと特殊ファイバーの滑らかさから $S_R\to\operatorname{Spec}R$ は
$S_0$ の近傍で滑らかであり、propernessから全体で滑らかになる。

最後に $H^1(C_0,\OO_{C_0}(1))=0$ と $h^0(\OO_{C_0}(1))=3$ より、
$H^0(S_R,\OO_{S_R}(1))$ は階数3の自由 $R$-加群であり、
$V\otimes R$ からの制限は全射である。
その核は階数2の直和因子で、$S_R$ を含む相対平面を定める。
$X$ とこの平面との交叉は $R$-flatな三次曲線族で、特殊ファイバーは $C_0$。
その構造層から $\OO_{S_R}$ への全射の核は、flatnessとNakayamaにより
特殊ファイバーの近傍で零であり、propernessにより全体で零である。
よって $S_R$ はこの平面切断の族に一致する。
\end{proof}

最後の結論では、一般ファイバーの非局所自由な部分が1次元であることを仮定した。
その部分が空または0次元である可能性を、接錐の計算だけで除いてはいない。
命題\ref{thm:all-compatible-arcs}では相対Picardによる記述を使い、
全ての非定数な適合形式弧で、その水平閉包が平面三次曲線の族に等しいことを示す。

\subsection{Brill--Noether境界の六次接錐}\label{sec:sextic-boundary}

本節は命題\ref{thm:bn-incidence}のincidenceから得られる、
局所自由でない対象の境界についての計算である。
適合部分空間 $C$ との包含や同一視は仮定しない。
$\mathcal I\subset X\times U$ を同命題のincidenceとし、
$D\subset U$ を射 $\mathcal I\to U$ のscheme-theoretic imageとする。

$W=\Lambda^2V$ とし、$a\in W$ の5個の符号付きPfaffianを
$q(a)=(q_0(a),\ldots,q_4(a))$ と書く。
符号は、階数4の場合に $[q(a)]\in\PP(V^*)$ が $a$ のkernel点になるように取る。
$X=V(F)$ の方程式を用いて
\[
 H_X=V\bigl(F(q_0(a),\ldots,q_4(a))\bigr)\subset\PP(W)
\]
と定める。

\begin{theorem}[境界の重複度と接錐]\label{thm:sextic-boundary}
$U$ を $e=[E]$ の周りで縮小すると、$D$ は既約かつ被約な因子である。
cohomologyの接座標 $T_eU=\Lambda^2V$ に関して
\[
 \operatorname{TC}_eD
 =\operatorname{Spec}
 \frac{\CC[a_{ij}]}{\bigl(F(q_0(a),\ldots,q_4(a))\bigr)}
\]
がスキームとして成立する。特に $\operatorname{mult}_eD=6$ である。
$H_X$ は正規な既約六次超曲面で、特異部分の台は $\Gr(2,V)$ である。
また
\[
 \PP_X(\Lambda^2Q)\longrightarrow H_X
\]
は解消であり、階数2の $[a]$ 上のファイバーは
$X\cap\PP(\ker a)$ という平面三次曲線である。
\end{theorem}

\begin{proof}
\textbf{像とblowup。}
$X$ は射影的なので $\mathcal I\to U$ はproperである。
$\mathcal I$ はcentral fiber $X$ の近傍で滑らかである。
非滑らかな点の閉集合のproperな像をまず除き、
次に $X$ に会わない成分のproperな像を除けば、滑らかな既約9次元スキームとしてよい。
そのscheme-theoretic imageはintegralである。
実際、affine開集合 $\operatorname{Spec}A\subset U$ 上で、そのイデアルは
\[
 A\longrightarrow\Gamma(\mathcal I\times_U\operatorname{Spec}A,\OO)
\]
の核であり、右辺は空でなければ整域である。
proper性によりこの核はcoherentなイデアルを与える。

$\mathcal I_e=X$ がスキームとして成立するので
$\mathfrak m_e\OO_{\mathcal I}=\mathcal I_{X/\mathcal I}$ である。
従ってblowupの普遍性からproperな射
\[
 f:B=\operatorname{Bl}_X\mathcal I\longrightarrow\operatorname{Bl}_eU
\]
を得る。sourceの例外因子を $J=\PP_X(\Lambda^2Q)$、targetの例外因子を
$E_U=\PP(T_eU)$ とすると、Cartier因子として
\begin{equation}\label{eq:sextic-exceptional-order}
                         f^*E_U=J
\end{equation}
であり、係数は1である。
$B\setminus J$ は $\mathcal I\setminus X$ と同型なので、
$f(B)$ は $D$ のstrict transform $\widetilde D$ である。

\textbf{例外因子上の射。}
normal bundleの包含から、$f|_J$ は
\[
 J=\{(p,[a])\in X\times\PP(W):\iota_p(a)=0\}
   \longrightarrow\PP(W)
\]
という射である。
階数4の部分では $[a]$ のkernelは一点であり、その点は $[q(a)]$ である。
kernel点への写像
\[
 \PP(W)_{\rk4}\longrightarrow\PP(V^*)
\]
は、ファイバーが $\PP(\Lambda^2\ker p)$ 内の開集合となる滑らかな射である。
従って $X$ の逆像は滑らかである。
階数2の場合のファイバーは平面三次曲線であり、必ず非空である。
ここで滑らかな三次三次元多様体は平面を含まない。
実際、平面のイデアルを $(u,v)$ として $F=uQ_1+vQ_2$ と書ければ、
平面上の二つの二次式 $Q_1,Q_2$ の共通零点で $X$ が特異になってしまう。
よって $J$ の像は $H_X$ であり、階数4の部分では $J\to H_X$ は同型である。

$J$ は $X$ 上の $\PP^5$-束なので既約である。
Pfaffianの共通零点 $\Gr(2,V)$ は余次元3であり、そこに別の超曲面成分は存在し得ない。
階数4の部分での滑らかさと合わせて、$F(q)$ は既約かつ被約な六次式である。
また $\widetilde D\cap E_U$ の台は8次元の $H_X$ なので、$D$ は9次元の因子となる。

\textbf{接錐の重複度。}
$H_X$ の階数4の部分の閉点 $z$ を取る。
$f^{-1}(z)$ は $J$ 上の一点 $y$ だけであり、
$J$ の接方向の写像は8次元の同型である。
\eqref{eq:sextic-exceptional-order}により、$J$ の法方向も
$E_U$ の法方向へ非零に写る。従って $df_y$ の階数は9であり、
$f$ は $y$ でunramifiedである。
proper性とファイバーの一意性から、$z$ の近傍で $f$ をfiniteにできる。
そのファイバーは被約な一点 $\operatorname{Spec}\CC$ なので、
finite加群に対するNakayamaにより局所的な構造環の射は全射である。
従って $f$ はそこで閉埋め込みとなり、$\widetilde D$ は滑らかで、
$E_U$ と横断的に交わる。

$U$ は正則なので $D$ の局所方程式を $d$ と書ける。
その最低次斉次項の台は既約な $H_X$ だから
$\operatorname{in}(d)=c\,F(q)^m$、$m\ge1$ である。
上の横断性は例外因子上での重複度が1であることを意味し、$m=1$ を得る。
これがスキームとしての接錐と重複度6を証明する。

\textbf{特異部分と解消。}
階数4の部分は既に滑らかである。
階数2では $q=0$ であり、$F$ は三次式なので $F(q)$ の全ての一次偏微分が零になる。
従って特異部分の台は $\Gr(2,V)$ である。
これは $H_X$ 内で余次元2なので、超曲面の $S_2$ 性と合わせて正規性を得る。
$J$ は滑らかであり、$J\to H_X$ はproperかつ双有理なので解消である。
\end{proof}

\begin{remark}[横断的な三次錐]
階数2の点の近傍で $a_{01}\ne0$ とし、交代行列をblock表示する。
可逆な $2\times2$ blockに対するSchur complementを取ると、
残る $3\times3$ 交代行列の三成分がGrassmannianに横断する座標になる。
この座標でPfaffianのkernel vectorは三成分に関して正確に線形である。
従って $H_X$ の横断sliceは、対応する平面切断
$X\cap\PP(\ker a)$ を底とする三次錐になる。
平面切断が滑らかな場合には、横断sliceの解消の例外曲線はその楕円曲線である。
本節は $D$ の接錐を決定しており、$D$ 自身をその錐へ形式同型にする主張ではない。
\end{remark}

命題\ref{thm:general-tangent-comparison}と合わせると、同じ接空間 $T_eU$ の中で
\[
 \bigl(\operatorname{Sing}(\operatorname{TC}_eD)\bigr)_{\mathrm{red}}
                  =\operatorname{TC}_eC
\]
が成立する。すなわち、7次元の適合条件の接錐は、
非局所自由な境界の六次接錐の特異部分を被約化したものに一致する。
この接錐の等式だけでは、$C\subset D$ や
$C=\operatorname{Sing}(D)_{\mathrm{red}}$ は導かれない。
命題\ref{thm:compatibility-in-bn-boundary}では相対Picardの全次数の記述から
形式包含 $\widehat C_c\subset\widehat D_e$ を別に証明する。
後者の特異部分全体との等式は、第7版の命題\ref{thm:global-bn-singular}で証明する。


\section{固定残余圏族の貼り合わせ対象：楕円曲線上の評価核}\label{sec:evaluation-gluing}

第\ref{sec:bn-algebraic}節では、明示的な固定残余圏族のgluing objectが、
generic fiberで階数3の安定torsion-free層になることを示した。
本節ではFermatの対称性と対象の変形を用い、この層を具体的に記述する。

\subsection{平面三次曲線上の次数3の線形束}

$X\subset\PP^4$ は滑らかなcubic threefold、
$C=X\cap\PP^2$ は滑らかな平面三次曲線とし、$j:C\hookrightarrow X$ を包含とする。
$L\in\Pic^3(C)$ に対して
\begin{equation}\label{eq:evaluation-kernel}
 0\longrightarrow F_L\longrightarrow H^0(C,L)\otimes\OO_X
 \xrightarrow{\mathrm{ev}}j_*L\longrightarrow0
\end{equation}
で $F_L$ を定める。次数3なので $h^0(C,L)=3$、$h^1(C,L)=0$ であり、
$L$ は大域切断で生成される。

\begin{theorem}[評価核の安定性と特異部分]\label{thm:evaluation-sheaf}
$F_L$ はGieseker-stableな階数3のtorsion-free層であり、
\[
 \operatorname{ch}(F_L)=(3,0,-H^2,0),\qquad
 F_L^{**}=H^0(C,L)\otimes\OO_X,\qquad F_L^{**}/F_L=j_*L.
\]
非局所自由な部分は、Fittingイデアルによるスキーム構造を含めて $C$ である。
さらに $L\ne\OO_C(1)$ なら $F_L\in\Ku(X)$ であり、
\[
 h^2(X,F_L(-2))=3.
\]
従ってこの層はUlrichではない。
\end{theorem}

\begin{proof}
$C$ は余次元2の局所完全交叉である。
$C$ 上で $L$ を局所自明化し、評価射の一成分を単元に直すと、
\[
 F_L\simeq\mathcal I_C\oplus\OO_X^{\oplus2}
\]
という局所表示を得る。
従ってtorsion-free性、二重双対とその商、
$\operatorname{Fitt}_3(F_L)=\mathcal I_C$ が従う。
$\chi(C,L)=3$ と $H\cdot C=3$ から
$\operatorname{ch}(j_*L)=(0,0,H^2,0)$ であり、Chern characterも得られる。

安定性を示す。$A\subset F_L$ を階数 $r=1,2$ の非零部分層とする。
$F_L\subset\OO_X^3$ なので $\mu_H(A)\le0$ である。
$\mu_H(A)<0$ なら既にreduced Hilbert polynomialは $p_A<p_{F_L}$ となる。
$\mu_H(A)=0$ とする。$\OO_X^3$ の中での飽和を $\overline A$ とすると、
飽和で生じるeffectiveな因子の次数は $\OO_X^3$ の半安定性により零である。
従って $c_1(\overline A)=0$。
$\det(\overline A)^{**}=\OO_X$ のPl\"ucker座標は定数なので、
そのgenericな $r$ 次元部分空間は固定される。
飽和性から $\overline A=W\otimes\OO_X$、$\dim W=r$ と書ける。

$K_W=F_L\cap(W\otimes\OO_X)$ とおく。
$0\ne W\subset H^0(C,L)$ の評価像は、整な曲線 $C$ 上でgeneric rank 1を持つので、
$\operatorname{ch}_2(K_W)=-[C]$ である。
$A\subset K_W$ の商は余次元2以上のtorsion層だから
\[
 \frac{H\cdot\operatorname{ch}_2(A)}r
 \le-\frac3r<-1
 =\frac{H\cdot\operatorname{ch}_2(F_L)}3.
\]
$c_1=0$ のもとでHilbert多項式の一次係数を比較すれば、
再び $p_A<p_{F_L}$ を得る。

\eqref{eq:evaluation-kernel}の大域切断の射は同型なので $H^*(X,F_L)=0$。
$L\ne\OO_C(1)$ なら $L(-1)$ は非自明な次数0の線形束であり、
$H^*(C,L(-1))=0$ である。
$\OO_X(-1)$ もacyclicなので $H^*(X,F_L(-1))=0$、従って $F_L\in\Ku(X)$。
最後に同じ完全列を $\OO_X(-2)$ でtwistすると
\[
 H^2(X,F_L(-2))\simeq H^1(C,L(-2))\simeq\CC^3
\]
となる。これはnormalized UlrichのACM条件\cite[Section 6]{FP}を満たさない。
\end{proof}

\subsection{Picardの原点をgluing objectへ送る射影}

$V=H^0(X,\OO_X(1))$、$U\subset V$ を $C$ の平面を切り出す2次元部分空間とし、
$\pi=i^*=L_{\OO_X}L_{\OO_X(1)}$ と書く。
$L$ の族に対して
\[
 \mathcal P_L=\pi(j_*L)[-1]
\]
を考える。$L\ne\OO_C(1)$ では命題\ref{thm:evaluation-sheaf}と
\eqref{eq:evaluation-kernel}から $\mathcal P_L\simeq F_L$ である。

trace線を
\[
 T=\Ext^2(Q,\OO_X(-1))\simeq H^3(\OO_X(-2))
\]
とする。$C$ のKoszul列をtwistした複体
\[
 [\,\OO_X(-1)\otimes\det U\longrightarrow
       \OO_X\otimes U\longrightarrow\OO_X(1)\,]
 \simeq j_*\OO_C(1)
\]
は次数 $-2,-1,0$ に置かれる。
$\pi$ は右の二項を殺し、\eqref{eq:E-triangle}から
$\pi(\OO_X(-1))[1]\simeq T\otimes E$ なので
\begin{equation}\label{eq:picard-special-object}
 \mathcal P_{\OO_C(1)}\simeq\det U\otimes T\otimes E.
\end{equation}
従って全ての平面についてPicardのこの点は、粗モジュライの同じ点 $[E]$ へ写る。
固定した $C$ では1次元因子を自明化して $\mathcal P_{\OO_C(1)}\simeq E$ とする。

\subsection{二対象の最小模型とPicard方向の接写像}

\begin{lemma}[射影の局所的な積の公式]\label{lem:mutation-product}
$A=\OO_X(-1)$、$B=Q$ とし、trace線を自明化する。
$0\ne a\in W=\Hom(A,B)=\Lambda^2V$ として
$F_a=\operatorname{Cone}(a:A\to B)$ とおく。
$F_a$ の形式変形は、$a$ の一座標を1に規格化した $f\equiv a$ と、
一つのパラメータ $t$ を用いて記述できる。
$E$ の形式変形の座標を $w\in W$ とすれば、
射影 $\pi$ が誘導する写像は適切なこの座標で
\[
                       (f,t)\longmapsto w=tf
\]
である。特に原点での接写像は
$(\dot f\bmod\CC a,\dot t)\mapsto\dot t\,a$ である。
\end{lemma}

\begin{proof}
Euler列とSerre双対性から、二対象 $A,B$ のfull dg部分圏のcohomologyは
\[
 \begin{aligned}
 \Hom^*(A,A)=\CC\id_A,\quad \Hom^*(B,B)=\CC\id_B,
 &\qquad\text{in degree }0,\\
 \Hom^*(A,B)=W\text{ in degree }0,
 &\qquad\Hom^*(B,A)=\CC g\text{ in degree }2.
 \end{aligned}
\]
である。strictly unitalな最小 $A_\infty$ 模型を取る。
非単位の合成可能な入力は $W$ と $g$ が交互に並ぶものに限られる。
$m_n$ の次数 $2-n$ から、偶数個の入力の出力は対角Homの次数2、
奇数個の入力の出力は $A\to B$ の次数1または $B\to A$ の次数3となる。
いずれも零なので、radical同士の積と全ての高次演算は零である。

$F_a$ は $A[1]\oplus B$ に微分 $a$ を入れたもので、
その自己Extは次数0で $\CC$、次数1で $(W/\CC a)\oplus\CC g$、他で零である。
$E$ の次数1の $W$ 成分は、三角の射 $\iota,\beta$ による
$\alpha\mapsto\beta[1]\alpha\iota$ そのものなので、
第\ref{sec:bn-algebraic}節のcohomologyの接座標と一致する。
従って最小模型によるArtin変形\cite{ELOI,ELOII}を
\[
 F_{f,t}:\ \delta=f+t g
\]
と書ける。一方 $E$ は微分 $g$ の対象であり、その変形を
$E_w:\delta=g+w$ と書ける。
これらの記述は可換Artin基底およびその形式逆極限上で用いる。
非triangularな行列を、任意の通常のパラメータで自動的にperfectな対象と
みなすことはしていない。

$A,B\in\OO_X^\perp$ なので、ここで $\pi$ は $L_B$ と一致する。
$\RHom(B,F_{f,t})$ は次数 $0,1$ の二項複体 $[R\xrightarrow{t}R]$ である。
left mutationの評価射のconeを消去すると $E_{tf}$ が残る。
より具体的には、$A[1]\oplus B$ 上の
\[
 \operatorname{diag}(1,t):F_{f,t}\longrightarrow E_{tf}
\]
は閉射であり、そのconeから $A$ 上の恒等写像を消去すると
$\operatorname{Cone}(t:B\to B)$ が残る。
$\pi$ はこのconeを殺し、$E_{tf}$ は $B$ と $\OO_X$ の右直交にある。
従って $\pi(F_{f,t})\simeq E_{tf}$ となる。
\end{proof}

\begin{lemma}[Picard方向はPl\"ucker直線を与える]\label{lem:picard-gluing-ks}
$\Pic^3(C)$ の $\OO_C(1)$ における形式近傍から $\Def_E$ への写像
$L\mapsto\mathcal P_L$ の接写像は非零である。
$\Ext^1(E,E)=\Lambda^2V\otimes T$ の同一視の下で、その像は
$\det U\otimes T$ というPl\"ucker直線である。
\end{lemma}

\begin{proof}
Euler列とKoszul列を比較すると
\[
 0\longrightarrow\OO_X(-1)\otimes\det U
   \xrightarrow{a}Q\longrightarrow F_{\OO_C(1)}\longrightarrow0
\]
を得る。ここで $a$ は $\Lambda^2U$ の非零元に対応する。
従って補題\ref{lem:mutation-product}を適用できる。

評価核の完全列から
\[
 \RHom(\OO_X(1),F_L)\simeq R\Gamma(C,L(-1))[-1].
\]
$L=\OO_C(1)$ では右辺は次数1と2に1次元ずつを持つ。
$0\ne\xi\in H^1(C,\OO_C)$ による線形束の一次変形は、その二つの間の微分を
\[
 H^0(C,\OO_C)\xrightarrow{\,\cup\xi\,}H^1(C,\OO_C)
\]
という非零写像に変える。
補題の模型では $f$ 方向はこの微分を変えず、$g$ 方向が非零に変える。
後者は $\delta=g$ の対象 $E$ が $\Ku(X)$ に属することからも確認できる。
従ってPicard方向の $\dot t$ は零でなく、射影後の接方向は $\dot t\,a$ である。
traceを戻すと、$H^1(\OO_C)\simeq\det U\otimes T$ から
$\Lambda^2V\otimes T$ へのtautologicalな直線の包含となる。
\end{proof}

\subsection{Fermatの固定部分空間と明示的族の同定}

\begin{theorem}[明示的族のgeneric gluing sheaf]\label{thm:explicit-gluing-sheaf}
第9節の固定残余圏族 $A'_h$ を $R=\CC[[h]]$ 上で考え、
\[
 C_0=X_0\cap\{u=v=0\}=V(x^3+y^3+z^3)\subset\PP^2
\]
とする。$L_0=\OO_{C_0}(1)$ を特殊値とする次数3の線形束の族 $L_R$ が存在し、
$\Pic^3(C_0)$ へのその形式弧の一次項は非零である。
適切な基底上の1次元因子の自明化の下で
\[
 \mathcal E_R\simeq i^*(j_*L_R)[-1].
\]
特に $L_\eta\ne\OO_{C_0}(1)$ であり、generic fiberは
\begin{equation}\label{eq:explicit-gluing-evaluation}
 0\longrightarrow\mathcal E_\eta\longrightarrow
 H^0(C_0,L_\eta)\otimes\OO_{X_0}
 \longrightarrow j_*L_\eta\longrightarrow0
\end{equation}
で表される。従って
\[
 \mathcal E_\eta^{**}\simeq\OO_{X_0}^{\oplus3},\qquad
 \mathcal E_\eta^{**}/\mathcal E_\eta\simeq j_*L_\eta,
 \qquad\operatorname{Fitt}_3(\mathcal E_\eta)=\mathcal I_{C_0}.
\]
この層はGieseker-stableであり、局所自由でなく、Ulrichでもない。
\end{theorem}

\begin{proof}
$\Gamma=\mu_3(u)\times\mu_3(v)$ とする。
$\Gamma$ は $C_0$ を点ごとに固定するので、$\Pic^3(C_0)$ と
$L\mapsto i^*(j_*L)[-1]$ の粗モジュライの像を固定する。
\eqref{eq:picard-special-object}の1次元因子のcharacterは
\[
 \det\langle u,v\rangle\otimes T:
 (\zeta_u\zeta_v)(\zeta_u\zeta_v)^{-1}=1
\]
であり、特殊値も $\Gamma$-同変に $E$ と同一視できる。

命題\ref{thm:fermat-bn-cone}のtraceを保持した表現を用いると
\[
 \Ext^1(E,E)=\Lambda^2V\otimes(\det V)^*,\qquad
 \dim\Ext^1(E,E)^\Gamma=1.
\]
固定部分は $u\wedge v$ の直線である。
有限群の作用は、characteristic zeroでReynolds作用素を使って形式的に線形化できる。
従って滑らかな $\Def_E$ の固定部分 $(\Def_E)^\Gamma$ は滑らかな形式曲線である。
補題\ref{lem:picard-gluing-ks}と形式逆関数定理から
\[
 \widehat{\Pic^3(C_0)}_{\,\OO_{C_0}(1)}
       \xrightarrow{\sim}(\Def_E)^\Gamma.
\]

$A'_h$ の関係式は $\Gamma$-作用を保ち、例外対象の持ち上げとmutationも自然である。
残余圏を固定した $\mathcal K_R$ と同一視した後、輸送された $\Gamma$ の作用と
一定の $\Gamma$ の作用との差は、特殊ファイバーで恒等となる自己同値である。
第7節の形式比較と $\HH^1(\mathcal K)=0$ により、その同型類は自明である。
従って $\mathcal E_R$ の標識付き変形の同値類は $\Gamma$ で固定される。
ここでは古典的な同値類上の固定性だけを使い、coherentな群作用の厳密化は仮定しない。

ゆえにこの弧は、上のPicard形式曲線の一意な弧に対応する。
既存の非零の一次gluing類と補題\ref{lem:picard-gluing-ks}から、
Picardへの弧も一次項が非零である。
Picard曲線のaffine chartに完備局所環の射を制限し、Poincar\'e束を引き戻せば、
実際の $R$ 上の $L_R$ を得る。
simple対象の普遍\'etale chartで比較すれば、必要な基底上の線形束の因子を除いて
$\mathcal E_R\simeq i^*(j_*L_R)[-1]$ となる。この因子は局所環 $R$ 上では自由である。

非定数性により $L_\eta\ne\OO_{C_0}(1)$ なので、generic fiberでは射影が不要となる。
\eqref{eq:explicit-gluing-evaluation}と残りの結論は
命題\ref{thm:evaluation-sheaf}から従う。
\end{proof}

従って命題\ref{thm:gluing-sheaf-arc}の特殊化の包含は、
この明示的な固定残余圏族については、固定曲線 $C_0$ そのものへ強化される。
一方、線形束 $L_R$ のPicard座標を、Sklyaninの平行移動点 $p(h)$ と直接結ぶ
明示公式はここでは得ていない。
この段階の証明は一本のFermat族を扱うものである。
命題\ref{thm:all-compatible-arcs}では、任意の滑らかな $X$ の
全ての非定数な適合形式弧へ、この評価核表示を拡張する。

\section{五生成元代数の内在的復元と次数付き分類}\label{sec:five-generator}

本節では、第\ref{sec:reconstruction}節の三生成元Sklyanin代数の復元を、
中心変数を二つ加えた五生成元代数の分類へ拡張する。
まず $k$ は標数零の代数閉体とする。
\[
 S=S(\alpha,\beta,\gamma)
 =\frac{k\langle x,y,z\rangle}
 {(\alpha yz+\beta zy+\gamma x^2,\
   \alpha zx+\beta xz+\gamma y^2,\
   \alpha xy+\beta yx+\gamma z^2)}
\]
の点曲線 $E$ が滑らかな楕円曲線であると仮定する。
具体的には $\alpha\beta\gamma\ne0$ および
\[
 (\alpha^3+\beta^3+\gamma^3)^3\ne27\alpha^3\beta^3\gamma^3
\]
を満たす通常の楕円Sklyaninの場合を扱う。
$L=\OO_E(1)$、$\sigma=T_p$ とし、
標準的な点加群の商の核を生成する次数3の中心元を $g_S$ と書く。
\cite[Definition--Lemma 2.13, Lemma 2.19]{WWY}により
\[
 S/(g_S)\simeq B(E,L,\sigma)
\]
であり、$g_S$ は非零scalar倍の選択を除いて定まる。
以後
\[
 A(S,g_S)=S[u,v]/(g_S+u^3+v^3),\qquad \deg u=\deg v=1
\]
とおく。$u,v$ は中心変数である。

\subsection{次数1の中心とquadratic cover}

\begin{theorem}[中心平面からのSklyanin代数の復元]\label{thm:intrinsic-sklyanin}
$A=A(S,g_S)$ の次数1の中心は
\[
 U_A:=Z(A)\cap A_1=ku\oplus kv
\]
である。従って $A$ から内在的に
\[
 B_A=A/(U_A)=S/(g_S)
\]
を回復できる。さらに、次数2までの積から作るquadratic cover
\begin{equation}\label{eq:intrinsic-quadratic-cover}
 \mathcal Q(B_A):=
 \frac{T_k((B_A)_1)}
 {\left(\ker\bigl[(B_A)_1\otimes(B_A)_1\longrightarrow(B_A)_2\bigr]\right)}
 \simeq S
\end{equation}
により、元の三生成元代数そのものを回復できる。
この構成は次数付き同型に関して関手的である。
\end{theorem}

\begin{proof}
$W=S_1$ とし、三つの二次関係が張る空間を $R_S\subset W\otimes W$ とする。
各関係はそれぞれ $\gamma x^2,\gamma y^2,\gamma z^2$ を含み、
$\gamma\ne0$ である。従って
\[
 R_S\cap\Lambda^2W=0.
\]
もし $w\in W$ が中心的なら、全ての $w'\in W$ について
$w\otimes w'-w'\otimes w\in R_S\cap\Lambda^2W$ となる。
$\dim W=3$ なので $w=0$ である。よって $Z(S)\cap S_1=0$。

三次関係による商は次数2までを変えないので、
\[
 A_1=W\oplus(ku\oplus kv),\qquad
 A_2=S_2\oplus(W\otimes(ku\oplus kv))\oplus\Sym^2(ku\oplus kv)
\]
である。$w+\lambda u+\mu v$ が中心的であるための条件は、
$w$ が $S$ の次数1の中心に入ることに等しい。
従って $U_A=ku\oplus kv$ を得る。
これで $B_A$ の表示も従う。

$g_S$ は次数3なので、$S$ と $B_A$ の次数1、次数2およびその間の積は同じである。
$S$ は二次関係で定義されるため、\eqref{eq:intrinsic-quadratic-cover}を得る。
中心、商、積のkernelはいずれも内在的な構成なので、関手性も従う。
\end{proof}

特に、この復元は中心全体の生成元を決定する必要がなく、
$\sigma$ が有限位数か無限位数かにも依存しない。

\subsection{次数3の中心元の線と全ての同型の形}

\begin{lemma}[次数3の中心元の線の不変性]\label{lem:canonical-central-line}
上のクラスのSklyanin代数の次数付き同型
$\phi:S\xrightarrow{\sim}S'$ は
\[
 \phi(g_S)=\lambda g_{S'}\qquad(\lambda\in k^\times)
\]
を満たす。
\end{lemma}

\begin{proof}
$S$ の次数3の元で、全てのpoint module $N$ に対して
$N_0$ を零化するものの空間を $I_3(S)$ とする。
$S\twoheadrightarrow B(E,L,\sigma)$ の核は $(g_S)$ である
\cite[Lemma 2.19]{WWY}。
商の次数3の元は $L_3=L\otimes\sigma^*L\otimes(\sigma^2)^*L$ の切断であり、
各点のpoint moduleへの作用はその評価である。
非零切断はある点で非零なので
\[
 I_3(S)=k g_S.
\]
次数付き同型はpoint moduleとその次数零の部分を保存するから、
$I_3(S)$ を $I_3(S')$ へ写す。
従って次数3の中心元の線は保存される。
\end{proof}

\begin{theorem}[五生成元代数の次数付き同型判定]\label{thm:five-generator-isom}
上のクラスの $S,S'$ と次数3の中心元の選択について、
\begin{equation}\label{eq:five-generator-isom}
 A(S,g_S)\simeq_{\mathrm{gr}}A(S',g_{S'})
 \quad\Longleftrightarrow\quad
 S\simeq_{\mathrm{gr}}S'
\end{equation}
が成り立つ。
さらに全ての次数付き同型は、次数1の分解
\[
 S_1\oplus(ku\oplus kv),\qquad S'_1\oplus(ku'\oplus kv')
\]
についてblock diagonalである。
すなわち、同型の全体は
\[
 \phi:S\xrightarrow{\sim}S',\qquad
 \tau:ku\oplus kv\xrightarrow{\sim}ku'\oplus kv'
\]
であって、ある $\lambda\in k^\times$ に対し
\begin{equation}\label{eq:five-generator-iso-pairs}
 \phi(g_S)=\lambda g_{S'},\qquad
 \Sym^3(\tau)(u^3+v^3)=\lambda(u'^3+v'^3)
\end{equation}
を満たす組と一致する。
\end{theorem}

\begin{proof}
\textbf{必要性と混合項の消滅。}
$\Phi:A(S,g_S)\to A(S',g_{S'})$ を次数付き同型とする。
命題\ref{thm:intrinsic-sklyanin}により中心平面を保存し、商とquadratic cover上で
$\phi:S\xrightarrow{\sim}S'$ を誘導する。
$X=\phi(x),Y=\phi(y),Z=\phi(z)$ とおけば
\[
 \Phi(x)=X+\ell_x,\quad \Phi(y)=Y+\ell_y,\quad
 \Phi(z)=Z+\ell_z,\qquad \ell_i\in ku'\oplus kv'
\]
と書ける。
$\gamma x^2+\alpha yz+\beta zy=0$ の像を取り、
$S'_1\otimes(ku'\oplus kv')$ 成分だけを見ると
\[
 2\gamma X\ell_x+(\alpha+\beta)(Y\ell_z+Z\ell_y)=0.
\]
$X,Y,Z$ は独立であり、$2\gamma\ne0$ なので $\ell_x=0$。
巡回した二つの関係から $\ell_y=\ell_z=0$ も従う。
よって混合項は存在しない。
なお $S\hookrightarrow A$ は、$u$ のmonicな三次関係による
$S[v]$-加群分解から単射である。

補題\ref{lem:canonical-central-line}により $\phi(g_S)=\lambda g_{S'}$ である。
$S'[u',v']$ の次数3で商の核は
$k(g_{S'}+u'^3+v'^3)$ なので、三次関係の像が零になる条件は
\eqref{eq:five-generator-iso-pairs}に一致する。

\textbf{十分性。}
$\phi:S\xrightarrow{\sim}S'$ を取り、
$\phi(g_S)=\lambda g_{S'}$ とする。
$k$ は代数閉体なので $\delta^3=\lambda$ を満たす $\delta\in k^\times$ を選べる。
$u\mapsto\delta u'$、$v\mapsto\delta v'$ とすれば、三次関係の像は
$\lambda(g_{S'}+u'^3+v'^3)$ であり、$\phi$ は求める同型へ延長する。
一般の\eqref{eq:five-generator-iso-pairs}についても同じ議論が成り立つ。
\end{proof}

同じ証明は、$u^3+v^3$ を相異なる三つの零点を持つ二変数三次式に置き換えても成り立つ。
代数閉体上では任意の相異なる三点の組は射影同値であり、
三次式のscalarも吸収できるため、この二変数三次式から追加の同型類のmodulusは生じない。

\subsection{固定Fermat曲線と四つの非可換三fold}

第\ref{sec:reconstruction}節と同様に
$E_0=V(x^3+y^3+z^3)$、$O=[1:-1:0]$ とする。
$p\in E_0\setminus E_0[3]$ に対して $S(p)=A(E_0,T_p)$ と書き、
その標準次数3の中心元の任意の非零倍を $g_p$ として
\[
 \mathcal A(p)=S(p)[u,v]/(g_p+u^3+v^3)
\]
とおく。その次数付き同型類は $g_p$ の非零scalar倍に依存しない。
第9節の $A'_h$ の幾何学的一般ファイバーは、この形に属する。

\begin{theorem}[固定楕円曲線上での完全な次数付き分類]\label{thm:four-five-generator}
$p,q\in E_0\setminus E_0[3]$ に対して
\[
 \mathcal A(p)\simeq_{\mathrm{gr}}\mathcal A(q)
 \quad\Longleftrightarrow\quad q\in\mu_6\cdot p.
\]
従って $t=a^3$ は、この構成に属する五生成元代数の幾何学的な次数付き同型類を
分離する座標である。
一般の非被約点の局所代数 $J_a$ の同型類に対して、
対応する五生成元代数の次数付き同型類はちょうど四つである。
\end{theorem}

\begin{proof}
命題\ref{thm:five-generator-isom}と
\eqref{eq:fixed-sklyanin-isom}から最初の同値が従う。
Sklyaninの同型判定そのものは、幾何対の射影共役による
\cite[Theorem 2.8, Theorem 4.16]{IMRelations}である。
$t=a^3$ は $E_0/\mu_6$ の座標であり、
\eqref{eq:fat-j-degree-four}は次数4であるから、最後の個数も従う。
三乗根の選択から新たな同型類が増えないことは、
命題\ref{thm:five-generator-isom}の十分性が保証している。
\end{proof}

この結果はpoint schemeの抽象スキームだけでは五生成元代数を分類できないことも示す。
一般の同じ $j_{\mathrm{fat}}$ に対応する四つの代数のpoint schemeは、
shiftを忘れれば全て $E_0\sqcup3\operatorname{Spec}J_a$ と同型である。
その楕円曲線成分上のpoint-module shiftは異なる平行移動を記録している。

また、第\ref{sec:points}節の再構成
\[
 A\simeq\bigoplus_{n\ge0}
 \operatorname{Hom}_{\qgr A}(\pi A,\pi A(n))
\]
により、構造対象と次数シフトを保つtailsの三つ組
$(\qgr A,\pi A,(1))$ の同値についても同じ同型判定が成り立つ。
ここで $S[u,v]$ はAS-regularであり、中心的な正則三次式による商 $A$ は
AS-Gorenstein、depth 4なので、この再構成に必要な低次数の局所コホモロジーは消滅する。
次数シフトの自然同型は積との整合性を含めて保つものとする。
これは、構造対象や偏極を指定しない圏同値の分類ではない。

\subsection{自己同型群の大きさと二重化する点}

$\operatorname{PAut}_{\mathrm{gr}}(A)
=\operatorname{Aut}_{\mathrm{gr}}(A)/k^\times$ と書く。
分母は全ての次数1生成元を同じscalar倍にする作用である。
$G_p=\{\alpha\in\operatorname{Aut}(E,O):\alpha(p)=p\}$ とおく。

\begin{theorem}[射影的な次数付き自己同型群]\label{thm:graded-automorphisms}
滑らかな平面三次曲線 $E$、$\OO_E(1)=\OO_E(3O)$ と
$p\in E\setminus E[3]$ に対して、上の代数 $A$ は完全列
\[
 1\longrightarrow\mu_3
 \longrightarrow\operatorname{PAut}_{\mathrm{gr}}(A)
 \longrightarrow(E[3]\rtimes G_p)\times S_3
 \longrightarrow1
\]
を持つ。特に
\[
 |\operatorname{PAut}_{\mathrm{gr}}(A)|=162\,|G_p|.
\]
固定Fermat曲線の場合、この位数は $t\ne4$ で162、$t=4$ で324である。
後者では $j_{\mathrm{fat}}=432$ となる。
\end{theorem}

\begin{proof}
命題\ref{thm:five-generator-isom}により全ての自己同型はblock diagonalである。
$S$ の射影的自己同型は、$(E,T_p)$ の射影的自己同型であり、
\cite[Theorem 4.12]{IMRelations}から
\[
 \operatorname{PAut}_{\mathrm{gr}}(S)=E[3]\rtimes G_p
\]
となる。二変数三次式の三つの零点の射影的自己同型群は $S_3$ である。
これらの任意の組の線形持ち上げを選ぶと、
$g$ と $u^3+v^3$ のscalar倍率は、第二の持ち上げを三乗根倍することで一致させられる。
従って右側への射は全射である。

両方の射影的作用が恒等なら、線形作用は
$S_1$ 上で $\lambda\id$、$\langle u,v\rangle$ 上で $\mu\id$ である。
三次関係を保つ条件は $\lambda^3=\mu^3$。
共通scalarを除けば $\mu/\lambda\in\mu_3$ だけが残るので、kernelを得る。

Fermatの場合は $\operatorname{Aut}(E_0,O)=\mu_6$ である。
位数3の自己同型で固定される点は3-torsionに入り、位数6の自己同型の固定点は $O$ だけである。
従って $p\notin E_0[3]$ で非自明なstabilizerを持つのは
$p\in E_0[2]\setminus\{O\}$ の場合だけで、そのstabilizerは $\{\pm1\}$ である。
命題\ref{thm:translation-double-cover}から、この三点の共通する $t$ は4である。
最後に $\mathcal J(4)=432$ を代入して得る。
\end{proof}

従って $j_{\mathrm{fat}}=432$ の四つの次数付き同型類のうち、
$t=4$ に対応する一つは射影的自己同型群が324元であり、残る三つは162元である。
実際、対応する四次式は
\[
 432\,t(t+8)^3-110592(t-1)^3
 =432(t-4)(t^3-228t^2+48t-64)
\]
と因数分解する。右側の三次式の判別式は $-2^{14}3^{11}\ne0$ であり、
$t=4$ での値は $-3456\ne0$ なので、四つの $t$ は相異なる。
同じ非被約局所環を持つ代数でも、自己同型群の大きさは異なり得る。

\subsection{DVR上の形式族への適用と基底体の制限}

$R=\CC[[h]]$ 上の次数ごとに有限自由な二つのSklyanin型族 $S_R,S'_R$ を考える。
ここでは、これらは次数1で生成され、次数2の関係のみで定義される $R$-代数とする。
その幾何学的一般ファイバーの点曲線は滑らかとし、
次数3の中心元 $G,G'$ は一般ファイバーで標準次数3の中心元の線を張るとする。
さらに $G,G'$ は次数3の自由加群のprimitiveな元であると仮定する。
この条件は第9節の $G'_h$ で成立する。

このとき、相対次数付き代数
\[
 A_R=S_R[u,v]/(G+u^3+v^3),\qquad
 A'_R=S'_R[u',v']/(G'+u'^3+v'^3)
\]
についても
\[
 A_R\simeq_{R,\mathrm{gr}}A'_R
 \quad\Longleftrightarrow\quad
 S_R\simeq_{R,\mathrm{gr}}S'_R
\]
が成り立つ。
実際、一般ファイバーの中心の計算と $R$-torsion freenessから
$Z(A_R)_1=Ru\oplus Rv$ となり、相対quadratic coverで必要性を得る。
逆に $\phi:S_R\to S'_R$ が与えられると、一般ファイバーで
$\phi(G)=\lambda G'$、$\lambda\in\CC((h))^\times$ である。
両辺のprimitive性を $S'_{R,3}$ の基底で比較すれば
$\lambda\in R^\times$ と分かる。
全ての $R$ の単元は三乗根を持つので、中心変数の共通スケールで十分性を得る。
特殊ファイバーでの同一視を保つ場合には、$\lambda(0)=1$ に対して
定数項1の三乗根を選べばよい。

このDVR上の主張を、任意の非被約Artin基底上の同型関手の同一視へ拡張してはいない。
例えば第9節の族を $R_N=R/(h^N)$、$N\ge2$ へ制限すると、
交換子が $h$ で割り切れるため
\[
 h^{N-1}x,\qquad h^{N-1}y,\qquad h^{N-1}z
\]
も次数1の中心に入る。従って中心の形成はこの非平坦な基底変換と可換ではない。
また、任意の基底体では必要な三乗根が存在するとは限らず、
\eqref{eq:five-generator-iso-pairs}のscalar条件を保持しなければならない。
一般の $\CC((h))$ 上のformを、幾何学的な $t$ または $j$ だけで分類する主張ではない。

\section{適合部分空間の相対Picard記述}\label{sec:picard-compatibility}

本節では、7次元の適合部分空間を、平面三次曲線上の線形束の変形から
高次まで特徴付ける。第\ref{sec:evaluation-gluing}節の一本のFermat族についての記述を、
任意の滑らかなcubic threefoldと任意の非定数な適合形式弧へ拡張する。

\subsection{全ての平面切断にわたる形式Picard}

$\mathcal G=\Gr(2,V)$、$\mathcal S\subset V\otimes\OO_{\mathcal G}$ をtautological束とし、
\[
 T=H^3(X,\OO_X(-2)),\qquad
 W=\Ext^1(E,E)=\Lambda^2V\otimes T,\qquad
 \mathcal D=\det\mathcal S\otimes T
\]
とおく。$U\in\mathcal G$ の定める平面三次曲線を
$C_U=X\cap\PP((V/U)^*)$ と書く。
$X$ は平面を含まないので、全ての $C_U$ は平面三次曲線である。
本節では、その被約性、既約性、滑らかさを仮定しない。

Universal曲線 $p:\mathcal C\to\mathcal G$ はflatで、
$\mathcal C\subset X\times\mathcal G$ は余次元2の相対正則埋め込みである。
Koszul列とcohomology and base changeから
\[
 p_*\OO_{\mathcal C}=\OO_{\mathcal G},\qquad
 R^1p_*\OO_{\mathcal C}=\mathcal D,\qquad
 R^{\ge2}p_*\OO_{\mathcal C}=0
\]
を得る。相対Picardの $L=\OO_{C_U}(1)$ というsectionに沿う形式完備化を
$\widehat{\mathcal P}$ とする。
このsectionでtranslateすると、$\widehat{\mathcal P}$ は相対次元1の可換形式群である。
曲線上の線形束の障害群 $H^2$ は零であり、Picardの滑らかさは
\cite[Tag 0DPJ]{StacksPicardSmooth}によっても得られる。
標数0の形式対数\cite[Appendix A.2, Proposition A.2.2]{PicardFormal}により、
接写像が恒等となる一意の形式群同型を得る：
\begin{equation}\label{eq:relative-picard-log}
 \widehat{\mathcal P}
 \simeq\widehat{\operatorname{Tot}(\mathcal D)}_{\mathcal G}
\end{equation}

線形束 $L$ の各Artin変形に対して
\[
 \mathcal P_L=i^*(j_*L)[-1]
\]
を取る。$j$ は正則埋め込みなので $j_*L$ はrelative perfectであり、
相対mutationはperfect性と基底変換を保つ。
\eqref{eq:picard-special-object}は、特異または非被約な $C_U$ に対しても
Koszul列からそのまま成立する。従ってzero section全体が $e=[E]$ に写り、
\[
 \Phi:\widehat{\mathcal P}\longrightarrow\Def_E
\]
を得る。全体上のPoincar\'e束の存在は仮定しない。
局所的な線形束の選択の差は基底からの線形束であり、粗対象モジュライへの射は降下する。

\begin{theorem}[相対Picard像とそのblowup]\label{thm:picard-formal-image}
$\Phi$ のscheme-theoretic imageは、$\Def_E$ の中の閉形式部分スキーム
$P$ である。$P$ は整な7次元空間で、抽象的にはGrassmannian錐の頂点完備化に同型である。
また
\[
 \widehat{\mathcal P}\simeq
 \widehat{\operatorname{Bl}_eP}_{\mathcal G}
\]
となり、$\Phi$ はそのblowdownと $P\hookrightarrow\Def_E$ の合成である。
cohomologyの接座標 $W=\Lambda^2V\otimes T$ では、
$\operatorname{TC}_eP$ はPl\"ucker錐である。
\end{theorem}

\begin{proof}
補題\ref{lem:picard-gluing-ks}の接写像の計算は、全ての平面三次曲線に適用できる。
実際、$h^0(\OO_{C_U})=h^1(\OO_{C_U})=1$ であり、
$0\ne\xi\in H^1(\OO_{C_U})$ とのcup積は $1\mapsto\xi$ という非零写像である。
この点を二対象模型で確認する。
Euler列とKoszul列から
$F_0=\operatorname{Cone}(\OO_X(-1)\otimes\det U\to Q)$ であり、
線束のArtin変形の評価核 $F_L$ は、このperfect対象の変形となる。
full dg部分圏 $\langle\OO_X(-1),Q\rangle$ とambient圏の自己Hom dg代数は同じなので、
一次Kodaira--Spencer類は、同じ最小模型の $(\dot f,\dot t)$ として計算できる。
ここで必要なのは $\epsilon^2=0$ 上の一次計算だけである。
$F_L\in\OO_X^\perp$ なので、Euler列と評価列から自然に
\[
 \RHom(Q,F_L)
  \simeq\RHom(\OO_X(1),F_L)[1]
  \simeq R\Gamma(C_U,L(-1))
\]
となる。左辺は二対象模型で $[R\xrightarrow{t}R]$ であり、$f$ は微分に寄与しない。
右辺の一次微分は上のcup積なので、$t$ の一次項は非零である。
$\det U$ を保持すると、この方向は $\det U\otimes T$ に値を取る。
二対象のmutationの公式から、zero sectionの法方向の写像は
\[
 \mathcal D=\det\mathcal S\otimes T
    \lhook\joinrel\longrightarrow W\otimes\OO_{\mathcal G}
\]
というtautologicalな包含になる。

\eqref{eq:relative-picard-log}により
\[
 R_P:=\Gamma(\widehat{\mathcal P},\OO)
   =\prod_{n\ge0}H^0(\mathcal G,\mathcal D^{-n})
\]
は、Pl\"ucker section ringの頂点完備化である。
GrassmannianのPl\"ucker埋め込みのprojective normalityと次数1での生成から、
$R_P\simeq\CC[[a_{ij}]]/(q_{ijkl})$ は正規な整域であり、
$\dim R_P=7$、$\mathfrak m_{R_P}/\mathfrak m_{R_P}^2=W^*$ である。

$\Def_E=\Spf A$、$A\simeq\CC[[z_1,\ldots,z_{10}]]$ と書く。
$\Phi$ が誘導する連続環準同型 $A\to R_P$ は、上の接写像から余接空間上同型になる。
完備局所環のNakayamaによりこの射は全射である。
従って $P=\Spf R_P\hookrightarrow\Def_E$ が閉形式像となる。

Pl\"ucker錐の頂点のblowupはtautological線束の全空間である。
trace線を含めれば、そのzero sectionに沿う完備化は
$\widehat{\operatorname{Tot}(\mathcal D)}_{\mathcal G}$ となる。
affine形式空間への射は大域切断の環準同型で決まるので、
元の $\Phi$ はこの標準的な縮約を通じて $P$ へ写る。
これによりblowupの主張も得られる。
\end{proof}

\subsection{Picard族に沿う最後のgluing対象の持ち上げ}

ここでは接錐の比較を用いず、Picard族そのものが適合条件を満たすことを証明する。
最後のgluing対象の障害を、特異な三次式も含む普遍族上で考えることが要点である。

\begin{theorem}[Picard族の全次数の適合性]\label{thm:picard-compatible}
$X\subset\PP(V^*)$ を滑らかなcubic threefold、
$U\subset V$ を2次元部分空間、$C=X\cap\PP((V/U)^*)$ を平面切断とする。
$j:C\hookrightarrow X$、$\pi=i^*$、
\[
 T=H^3(X,\OO_X(-2)),\qquad D_U=\det U\otimes T
\]
とおく。可換局所Artin $\CC$-代数 $R$ 上で、
$L_R$ を $\OO_C(1)$ の線形束としての変形とする。
このとき
\[
 E_{L_R}=D_U^{-1}\otimes\pi(j_*L_R)[-1]\in\mathcal K_R
\]
に沿って、$G$ は $\mathcal B_{E_{L_R}}$ のperfect対象に持ち上がる。
従ってPicard形式近傍の射は、閉形式部分空間
$\widehat C_c\subset\Def_E$ を経由する。
この主張は平面切断 $C$ が特異な場合にも成立する。
\end{theorem}

\begin{proof}
\textbf{普遍族を取る基底。}
$V=U\oplus Z$、$\dim Z=3$ と分解し、
\[
 \Sigma_U=\{F\in\Sym^3V:F|_Z\ne0\}
 \simeq\Aff^{25}\times(\Aff^{10}\setminus\{0\}),
 \qquad H_U=GL(U)\times GL(Z)
\]
とおく。$H_U$ はcharacteristic zeroでreductiveである。
除いた部分は余次元10の線形部分空間なので、局所cohomologyの計算から
\begin{equation}\label{eq:picard-universal-base-cohom}
 \Gamma(\Sigma_U,\OO)=\CC[\Sym^3V],\qquad
 H^r(\Sigma_U,\OO)=0\quad(1\le r\le8)
\end{equation}
を得る。特に係数を判別式で割ることは許していない。

$\mathcal X\to\Sigma_U$ を普遍三次超曲面、
$q:\mathcal C\to\Sigma_U$ をその $\PP(Z^*)$ による切断とする。
両者はflatかつproperであり、$\mathcal C\hookrightarrow\mathcal X$ は
余次元2の正則埋め込みである。
Euler列、制限列、Gorenstein双対性による第\ref{sec:gluing-objects}節と
補題\ref{lem:universal-tangent-lift}の計算は、
特異・非被約なファイバーについてもperfect対象に対して成立する。
trace線を保持すると、相対自己Extは
\[
 \mathcal E xt^r_{\mathcal B/\Sigma_U}(G,G)=
 \begin{cases}
  \OO_{\Sigma_U}&r=0,\\
  \Sym^2V\otimes T&r=2,\\
  0&\text{その他}
 \end{cases}
\]
であり、任意の基底変換と可換である。
普遍制限列の境界写像から、同変な同一視
\[
 T\simeq(\det V)^*\otimes\OO_{\Sigma_U},\qquad
 R^1q_*\OO_{\mathcal C}\simeq D_U
       =\det U\otimes(\det V)^*\simeq(\det Z)^*
\]
を得る。また $q_*\OO_{\mathcal C}=\OO_{\Sigma_U}$、
$R^rq_*\OO_{\mathcal C}=0$ ($r\ge2$) である。

\textbf{同変なPicardパラメータ。}
$R_n=\bigoplus_{r=0}^nD_U^{-r}$ を次数 $n$ で切った対称代数とし、
$\Sigma_n=\operatorname{Spec}_{\Sigma_U}R_n$ とする。
$R_n\twoheadrightarrow R_{n-1}$ の核は $I_n=D_U^{-n}$ で、
正次数のイデアルにより零化される。
$\OO_{\mathcal C}(1)$ の線形束としての持ち上げの障害は
\[
 H^2(\mathcal C,q^*I_n)^{H_U}=0
\]
に入る。消滅はLeray列と\eqref{eq:picard-universal-base-cohom}による。
同変な場合も、不変部分を取る関手が完全なので同じ結論を得る。
一次変形の類を $D_U\otimes D_U^*$ の恒等元に取れば、
$H_U$-同変な線形束の整合的な族 $L_n$ を逐次構成できる。
各ファイバーでは $h^1(\OO_C)=1$、$H^2(\OO_C)=0$ なので、
この族はPicard形式近傍の座標を与える。
ここでは相対Picardスキーム上の大域的なPoincar\'e束の存在を仮定していない。

\eqref{eq:picard-special-object}によって
\[
 \mathcal E_n=D_U^{-1}\otimes\pi(j_*L_n)[-1]
\]
は $E$ の同変な変形である。
$j$ が正則埋め込みなので $j_*L_n$ はperfectであり、相対mutationもperfect性を保つ。
これを使って $\mathcal B_n=\langle\mathcal K_{\Sigma_n},Q\rangle_{\mathcal E_n}$
を構成する。

\textbf{変化するdg圏での持ち上げ障害。}
square-zero拡大に沿うdg加群の持ち上げを、局所的なsemifree模型で計算する。
次数付きの持ち上げに微分と加群作用を持ち上げると、その平方の失敗は
$I_n$ を係数とする次数2の自己準同型cocycleである。
持ち上げを変えるとcoboundaryだけ変わり、その類が零なら次数1の補正で
平方を零にできる。
これは微分の持ち上げの障害計算\cite[Theorem 3.8]{LowenObstruction}、
およびMaurer--Cartan模型\cite[Proposition 7.1]{ELOI}のcurvature計算である。
ここでは既に指定された $\mathcal B_n$ 上のbar模型に適用するので、
圏が一定であることを仮定していない。

この構成の還元が導来還元であることも確認する。
affine基底上で $\mathcal B_n$ のcofibrant模型を取り、前段階の加群をsemifreeに取る。
その各評価は基底上h-projectiveかつperfectなので、有界な有限projective複体と
homotopy同値である。
微分を補正した後も次数付き評価はprojectiveであり、
crude lifting lemma\cite[Corollary 3.11]{LowenObstruction}により、
このhomotopy同値を有界な有限projective複体の持ち上げへ輸送できる。
従って新しい評価はK-flatで、通常の還元と導来還元が一致する。
圏のcofibrant模型の基底変換と導来テンソル積の結合則から、
加群としての導来還元も前段階の加群になる。
無限複体の次数付きflat性だけからこの結論を出しているわけではない。

\v Cech descentも含めると、$G_{n-1}$ の持ち上げ障害は
\[
 \mathbb H^2\bigl(\Sigma_U,
       \RHom_{\mathcal B/\Sigma_U}(G,G)\otimes I_n\bigr)^{H_U}
\]
に入り、存在すれば同型類の差は同じ複体の次数1に入る。
相対自己Extと\eqref{eq:picard-universal-base-cohom}から、次数1は零、次数2は
\begin{equation}\label{eq:picard-obstruction-equivariants}
 \left[\CC[\Sym^3V]\otimes\Sym^2V\otimes(\det V)^*
                         \otimes D_U^{-n}\right]^{H_U}
\end{equation}
となる。$H^3(\Sigma_U,\OO)=0$ により、local-to-global列の
次数2からのtransgressionもない。
各段階を同変な加群の圏で構成しているので、同型類の一意性だけから
coherentな群作用の存在を推測する必要はない。

\textbf{全ての次数で障害が消えること。}
\eqref{eq:picard-obstruction-equivariants}の元を、多項式同変写像
\[
 b:\Sym^3V\longrightarrow
       \Sym^2V\otimes(\det V)^*\otimes D_U^{-n}
\]
と見る。$a_s=(s\id_U,\id_Z)$ の作用について、
$F$ の $\Sym^3Z$、$U\otimes\Sym^2Z$、$\Sym^2U\otimes Z$、$\Sym^3U$
成分のweightは、それぞれ $0,1,2,3$ である。
一方、$D_U$ のweightは0であり、目標の
\[
 \Sym^2U\otimes(\det V)^*\otimes D_U^{-n},\quad
 U\otimes Z\otimes(\det V)^*\otimes D_U^{-n},\quad
 \Sym^2Z\otimes(\det V)^*\otimes D_U^{-n}
\]
のweightは $0,-1,-2$ である。
従って多項式写像 $b$ の負weight成分は零であり、残る成分は
$F|_Z$ のみを用いて $\Sym^2U$ の成分に値を取る。
ところが $SL(U)$ は $F|_Z$、$D_U$、$(\det V)^*$ に自明に作用し、
$\Sym^2U$ は $SL(U)$-不変ベクトルを持たない。従って $b=0$ である。
この議論は全ての $n\ge1$ に成立する。

\textbf{perfect性とArtin基底への特殊化。}
以上から $G_0=G$ を出発点として、整合的な同変dg加群 $G_n$ を得る。
特異な三次式を含む普遍基底上では、各評価が基底上perfectであるという
pseudo-perfect性だけを用いれば十分である。
評価複体の導来還元は $G$ の評価でありperfectなので、
冪零イデアルに関するperfect性の判定\cite[Lemma 15.80.4, Tag 07LU]{StacksNilPerfect}
により、この性質は各段階で保たれる。

最後に滑らかな $X_F$ に制限する。
$\mathcal K_F$ はsmooth properで、$\mathcal E_n$ はperfectなので、
$\mathcal B_n|_F$ もsmooth properである。
このgluingのsmoothnessはArtin基底上でも成立する。
実際、$A$ を $\mathcal K_F\otimes R_n$ のdg模型、$M$ を $\mathcal E_n$ の
perfect加群とし、上三角代数を $T_{\mathrm{tri}}$、その頂点冪等元を $e,f$ と書く。
対角bimoduleの三角
\[
 \begin{split}
 T_{\mathrm{tri}}e\otimes_A^{\mathbf L}M\otimes_{R_n}^{\mathbf L}fT_{\mathrm{tri}}
 &\longrightarrow
 (T_{\mathrm{tri}}e\otimes_A^{\mathbf L}eT_{\mathrm{tri}})
 \oplus(T_{\mathrm{tri}}f\otimes_{R_n}^{\mathbf L}fT_{\mathrm{tri}})\\
 &\longrightarrow\Delta_{T_{\mathrm{tri}}}
 \end{split}
\]
の左項と中項は、$A$ のsmoothnessと $M$ のperfect性によりperfectである。
従って対角もperfectであり、体上という仮定は必要ない。
perfectな対角bimoduleの有限分解とテンソルすれば、
pseudo-perfectな $G_n|_F$ はperfectな $\mathcal B_n|_F$-加群になる。
任意のArtin線形束変形は、上のPicardパラメータ族のある有限段階の引き戻しなので、
その引き戻しが所要のperfectな $G$ の持ち上げを与える。
\end{proof}

\subsection{適合条件との高次までの一致}

\begin{theorem}[適合部分空間の幾何学的記述]\label{thm:picard-compatibility-equality}
命題\ref{thm:algebraic-compatibility}の適合部分空間の完備化 $\widehat C_c$ は、
命題\ref{thm:picard-formal-image}の $P$ に、
$\Def_E$ 内の閉形式部分スキームとして一致する。
従って
\[
 \widehat{\mathcal P}
 \simeq\widehat{\operatorname{Bl}_c C}_{\mathcal G}.
\]
\end{theorem}

\begin{proof}
命題\ref{thm:picard-compatible}により、固定した各 $U\in\mathcal G$ のPicard形式曲線は
適合条件を満たす。
適合イデアルの任意の元を $\widehat{\mathcal P}$ に引き戻し、形式対数の次数ごとに展開する。
その $n$ 次係数は $H^0(\mathcal G,\mathcal D^{-n})$ の切断であり、
全ての閉点 $U$ で零になる。$\mathcal G$ は被約なので、この係数自体が零である。
従って $\Phi$ 全体が $\widehat C_c$ を経由し、$P\subset\widehat C_c$ を得る。

既存の形式変形の計算により $\widehat C_c$ は整な7次元空間であり、
命題\ref{thm:picard-formal-image}により $P$ も整な7次元空間である。
正則完備局所環 $A$ 内の対応する二つの素イデアルは、包含関係を持ち高さが等しいので一致する。
blowupの主張は命題\ref{thm:picard-formal-image}から従う。
この証明は、全次数での包含と等次元性を使っており、接錐の一致だけからは結論していない。
\end{proof}

\subsection{相対cohomologyの単純零点}

\begin{lemma}[Picardのidentity sectionの重複度]\label{lem:relative-picard-theta}
$L=\OO_{C_U}(1)$ のsectionの形式近傍で、$Rp_*L(-1)$ は局所的に
\[
 [\,\mathcal L_0\xrightarrow{\,d\,}\mathcal L_1\,]
 \qquad\text{（次数 }0,1\text{）}
\]
という二つの線束の間の複体で表せる。
zero sectionの局所形式パラメータを $t$ とすると、線束を自明化した表示で
$d=t\cdot u$、$u$ は単元である。
この主張は特異または非被約な平面三次曲線を含む。
\end{lemma}

\begin{proof}
$p$ はproper flatな曲線の族であり、$L(-1)$ は線束なので、
$Rp_*L(-1)$ はperfectかつTor amplitude $[0,1]$ を持つ。
zero sectionの各点では、導来ファイバーは $R\Gamma(\OO_{C_U})$ であり、
次数0と1にそれぞれ1次元のcohomologyを持つ。
局所的なperfect表示から可逆な微分成分を消去すると、上の二項表示になる。

zero sectionに沿ってcohomologyは $\OO_{\mathcal G}$ と $\mathcal D$ であり、
基底変換と可換なので $d$ はそこで恒等的に零になる。従って $d=tu$ と書ける。
$d$ の法方向の一次項はPicardのKodaira--Spencer類とのcup積
\[
 H^0(\OO_{C_U})\longrightarrow H^1(\OO_{C_U}),\qquad1\longmapsto\xi
\]
である。これは非零の法方向に対して同型なので、$u$ はzero section上で単元であり、
形式近傍全体でも単元となる。
\end{proof}

この補題は、被約でない曲線上の任意の非自明な次数0の線束がacyclicであると
仮定するものではない。必要な消滅は、identity sectionの近傍の相対複体から直接得られる。

\subsection{全ての非定数な適合形式弧の一般ファイバー}

\begin{theorem}[一般gluing sheafと非自由支持の決定]\label{thm:all-compatible-arcs}
$X$ を任意の滑らかなcubic threefold、$R=\CC[[h]]$、
$\gamma:\Spf R\to\widehat C_c$ を非定数の適合形式弧とする。
対応する貼り合わせ対象を $E_R$ とすると、次を満たす平面の族 $U_R$ と
線形束 $L_R$ が得られる。
\[
 C_R=X_R\cap\PP((V_R/U_R)^*),\qquad
 L_R|_{C_0}=\OO_{C_0}(1).
\]
一般ファイバーでは
\begin{equation}\label{eq:all-arcs-evaluation}
 0\longrightarrow E_\eta
 \longrightarrow H^0(C_\eta,L_\eta)\otimes\OO_{X_\eta}
 \longrightarrow j_*L_\eta\longrightarrow0
\end{equation}
となる。$E_\eta$ はGieseker-stableな階数3のtorsion-free層であり、
\[
 \operatorname{Fitt}_3(E_\eta)=\mathcal I_{C_\eta},\qquad
 E_\eta^{**}\simeq\OO_{X_\eta}^{\oplus3},\qquad
 E_\eta^{**}/E_\eta\simeq j_*L_\eta.
\]
従って非局所自由部分の水平閉包は、スキームとして $C_R$ に等しい。

さらに、cohomologyの接座標で $\gamma$ の最初の非零係数を $h^m a_m$ とすると、
$a_m$ は分解可能であり、対応する $U_0\in\Gr(2,V)$ に対して
\[
 C_0=X\cap\PP(\operatorname{ann}U_0)=Z(s_{a_m}).
\]
一般の層は局所自由でもACMでもなく、normalized Ulrichではない。
\end{theorem}

\begin{proof}
$\gamma$ が非定数なので、$\mathfrak m_c\OO_R=(h^m)$ は非零の可逆イデアルである。
blowupの普遍性により $\gamma$ は一意に $\operatorname{Bl}_cC$ へ持ち上がる。
命題\ref{thm:picard-compatibility-equality}により、この持ち上げは相対Picardの形式点となり、
$U_R$ と $L_R$ の形式族を与える。
Grassmannianの射影性と、properな $C_R$ 上のGrothendieck existenceにより、
これらは実際の $R$ 上の平面族と線束に実現する。
基底からの線束の曖昧さは局所環 $R$ 上では自明である。

Picardの形式パラメータ $t_R$ は非零である。
補題\ref{lem:relative-picard-theta}から
\[
 R\Gamma(C_R,L_R(-1))\simeq[R\xrightarrow{\,t_Ru\,}R]
\]
となり、一般ファイバーではacyclicである。
また $h^0(\OO_{C_0}(1))=3$、$h^1(\OO_{C_0}(1))=0$ なので、
$p_*L_R$ は階数3の自由加群で、$R^1p_*L_R=0$。
特殊ファイバーでの大域生成性、Nakayamaとproper性により評価射は全射となる。
その核を $F_R$ とすれば
\[
 0\longrightarrow F_R\longrightarrow p_*L_R\otimes\OO_{X_R}
 \longrightarrow j_*L_R\longrightarrow0
\]
は $R$-flatな層の完全列であり、$\pi F_R\simeq E_R$ である。
一般ファイバーで $F_\eta$ と $F_\eta(-1)$ はacyclicなので
$F_\eta\in\Ku(X_\eta)$、従って $E_\eta=F_\eta$ を得る。
安定性は、既存のstable対象モジュライのchartと
\cite[Proposition 5.5]{FLZ}から従う。

$C_R\subset X_R$ は余次元2の正則埋め込みである。
線束を局所自明化し、評価の一成分を単元に直せば
$F_R\simeq\mathcal I_{C_R}\oplus\OO_{X_R}^2$ となる。
これがFittingイデアル、二重双対とその商を与える。
$C_R$ は $R$-flatなので、$C_\eta$ のスキームとしての閉包は $C_R$ である。
blowupの例外点は $[a_m]$ であり、法方向のtautologicalな接写像から
特殊平面は $U_0$、特殊曲線は表示の $C_0$ となる。

最後に、特殊ファイバーの $L_0(-2)=\OO_{C_0}(-1)$ は大域切断を持たない。
これは非被約な平面三次曲線についても、$\PP^2$ 上の完全列
$0\to\OO(-4)\to\OO(-1)\to\OO_{C_0}(-1)\to0$ から従う。
半連続性により $H^0(C_\eta,L_\eta(-2))=0$ であり、
$\chi(C_\eta,L_\eta(-2))=-3$ だから $h^1(C_\eta,L_\eta(-2))=3$。
\eqref{eq:all-arcs-evaluation}と $H^1(\OO_X(-2))=H^2(\OO_X(-2))=0$ から
\[
 h^2(X_\eta,E_\eta(-2))=h^1(C_\eta,L_\eta(-2))=3
\]
となり、ACM消滅を満たさない。
\end{proof}

\begin{theorem}[適合部分空間は非局所自由境界に含まれる]
\label{thm:compatibility-in-bn-boundary}
第\ref{sec:sextic-boundary}節のBrill--Noether境界 $D$ に対して、
$\widehat C_c\subset\widehat D_e$ が閉形式部分スキームとして成立する。
\end{theorem}

\begin{proof}
固定した各 $U\in\mathcal G$ のPicard形式曲線 $\CC[[t]]$ を取る。
補題\ref{lem:relative-picard-theta}と評価核の表示により、
その一般対象の非局所自由部分は非空な平面三次曲線 $C_U$ である。
従って一般点はBN incidenceの像 $D$ に属する。
$D$ の局所方程式 $d$ は $\CC((t))$ で零になるので、整域 $\CC[[t]]$ でも零である。
$\Phi^*d$ の形式対数の各次数の係数は、$\mathcal G$ 上の切断であり、
全ての閉点 $U$ で零になる。従って各係数は零であり、
$d$ は $P=\widehat C_c$ 上で零となる。
この包含は全次数のPicard記述を使う帰結であり、接錐の包含だけから推論したものではない。
\end{proof}

\begin{remark}
全 $R$ 上の等式は $E_R\simeq\pi F_R$ である。
特殊ファイバーでは $F_0$ は層、$E$ は二項複体なので、
一般ファイバーでの $E_\eta=F_\eta$ と区別する。
本節は適合部分空間そのものの相対Picard記述を与えるが、
Brill--Noether境界の特異部分全体が $C$ に一致することまでは主張しない。
\end{remark}

